\documentclass[11pt]{article}
\usepackage[T1]{fontenc}
\usepackage{lmodern}
\usepackage[margin=1.02in]{geometry}
\usepackage{amsmath,amssymb,amsthm,mathtools,mathrsfs}
\usepackage{bm,enumitem,booktabs,array,tabularx}
\usepackage{microtype,aliascnt,needspace}
\usepackage[hidelinks]{hyperref}
\usepackage[nameinlink,capitalise]{cleveref}
\setlength{\emergencystretch}{2em}
\raggedbottom
\numberwithin{equation}{section}
\newtheorem{theorem}{Theorem}[section]
\newaliascnt{proposition}{theorem}
\newtheorem{proposition}[proposition]{Proposition}
\aliascntresetthe{proposition}
\newaliascnt{lemma}{theorem}
\newtheorem{lemma}[lemma]{Lemma}
\aliascntresetthe{lemma}
\newaliascnt{corollary}{theorem}
\newtheorem{corollary}[corollary]{Corollary}
\aliascntresetthe{corollary}
\theoremstyle{definition}
\newaliascnt{definition}{theorem}
\newtheorem{definition}[definition]{Definition}
\aliascntresetthe{definition}
\newaliascnt{assumption}{theorem}
\newtheorem{assumption}[assumption]{Assumption}
\aliascntresetthe{assumption}
\newaliascnt{example}{theorem}
\newtheorem{example}[example]{Example}
\aliascntresetthe{example}
\theoremstyle{remark}
\newaliascnt{remark}{theorem}
\newtheorem{remark}[remark]{Remark}
\aliascntresetthe{remark}
\crefname{theorem}{Theorem}{Theorems}
\Crefname{theorem}{Theorem}{Theorems}
\crefname{proposition}{Proposition}{Propositions}
\Crefname{proposition}{Proposition}{Propositions}
\crefname{lemma}{Lemma}{Lemmas}
\Crefname{lemma}{Lemma}{Lemmas}
\crefname{corollary}{Corollary}{Corollaries}
\Crefname{corollary}{Corollary}{Corollaries}
\crefname{definition}{Definition}{Definitions}
\Crefname{definition}{Definition}{Definitions}
\crefname{assumption}{Assumption}{Assumptions}
\Crefname{assumption}{Assumption}{Assumptions}
\crefname{example}{Example}{Examples}
\Crefname{example}{Example}{Examples}
\crefname{remark}{Remark}{Remarks}
\Crefname{remark}{Remark}{Remarks}
\newcommand{\R}{\mathbb R}
\newcommand{\C}{\mathbb C}
\newcommand{\N}{\mathbb N}
\newcommand{\Z}{\mathbb Z}
\newcommand{\dd}{\mathrm d}
\newcommand{\ii}{\mathrm i}
\newcommand{\id}{\mathrm{id}}
\newcommand{\Tr}{\operatorname{Tr}}
\newcommand{\Rea}{\operatorname{Re}}
\newcommand{\supp}{\operatorname{supp}}
\newcommand{\rank}{\operatorname{rank}}
\newcommand{\Ker}{\operatorname{Ker}}
\newcommand{\Ran}{\operatorname{Ran}}
\newcommand{\diag}{\operatorname{diag}}
\newcommand{\Sym}{\operatorname{Sym}}
\newcommand{\End}{\operatorname{End}}
\newcommand{\SU}{\mathrm{SU}}
\newcommand{\U}{\mathrm U}
\newcommand{\HS}{\mathrm{HS}}
\newcommand{\Ric}{\operatorname{Ric}}
\newcommand{\Riem}{\operatorname{Riem}}
\newcommand{\Scal}{\operatorname{Scal}}
\newcommand{\vol}{\operatorname{vol}}
\newcommand{\norm}[1]{\lVert#1\rVert}
\newcommand{\abs}[1]{\lvert#1\rvert}
\newcommand{\ip}[2]{\langle#1,#2\rangle}
\newcommand{\Dtest}{\mathcal D}
\newcommand{\A}{\mathcal A}
\newcommand{\M}{\mathcal M}
\newcommand{\Sact}{\mathscr S}
\newcommand{\Vtest}{\mathscr V}
\newcommand{\Fcurv}{\mathsf F}
\newcommand{\Eins}{\mathsf G}
\newcommand{\Evar}{\mathcal E}
\newcommand{\bPsi}{\overline\Psi}
\newcommand{\dV}{\dd V}
\newcommand{\wt}{\widetilde}
\newcommand{\Iface}{\mathfrak I}
\newcolumntype{L}[1]{>{\raggedright\arraybackslash}p{#1}}

\title{\textbf{Finite-Action Closure and Classical Einstein--Standard-Model Limits\\in Renewal Geometry}}
\author{Aur\'elien P\'elissier$^{1,\dagger}$\\
{\small $^{1}$\'Ecole normale sup\'erieure Paris-Saclay, 91190 Gif-sur-Yvette, France}\\[-0.25em]
{\small $^{\dagger}$Correspondence: \href{mailto:aurelien.pelissier.38@gmail.com}{aurelien.pelissier.38@gmail.com}}}
\date{}
\begin{document}
\maketitle

\begin{abstract}
The finite reconstruction of Lorentzian spacetime and Standard-Model internal structure does not by itself imply the emergence of their coupled classical dynamics. We isolate a minimal finite Lorentzian Standard-Model interface and prove an autonomous variational closure theorem from that interface to the classical Einstein--Standard-Model system. The headline theorem is deliberately subsequential: finite-cutoff compactness, first-variation consistency, and physically scaled stationarity produce an Einstein--Standard-Model limit without assuming summable Cauchy transport; adjacent-cutoff summability is used only as an optional uniqueness and route-independence upgrade. To separate genuine a-priori control from compactness packaged as a hypothesis, we also give a regular branch on which uniform Sobolev bounds imply the required bosonic compactness directly. In the fermionic sector, a Lorentzian symmetric-hyperbolic Dirac estimate replaces an elliptic graph-norm assumption: convergent initial data, uniform higher spatial control, and a vanishing Dirac residual yield strong $H^1$ compactness by energy stability and interpolation. First-variation consistency together with common-action stationarity then yields the Einstein--Yang--Mills--Higgs--Dirac--Yukawa equations in distributions. On weaker energy-level sequences we retain the Yang--Mills and Higgs stress defects produced by oscillation or concentration. Exact finite critical points and a nondegenerate shadowing theorem provide complementary stationary branches. The companion Renewal Geometry constructions are used only to realize the finite structural interface; the analytic closure arguments are proved independently here. We do not assert microscopic selection for arbitrary renewal histories or construct an interacting quantum continuum measure.
\end{abstract}

\medskip
\noindent\textbf{Keywords.} Renewal Geometry; classical Standard Model; Einstein equations; finite-action closure; continuum limit; native selection; variational consistency; stress defects; constraint propagation.

\section{Introduction}
\label{sec:intro}

A finite theory can reconstruct the correct spacetime signature, internal symmetry, and matter representations without yet producing a classical field theory. The remaining problem is dynamical and analytic: fields at different cutoffs must be compared on one spacetime, the nonlinear quantities entering the common action must have controlled limits, and finite stationarity must survive in the physical normalization of the variations. For gravity coupled to matter, convergence of the matter Euler equations alone is insufficient. The metric variation probes quadratic derivative products and quartic potentials, so oscillation and concentration can survive even when lower-order fields converge.

This paper addresses that final classical handoff in Renewal Geometry. The published spacetime construction reconstructs Lorentzian kinematics and gravitational common-action data from conditional renewal histories \cite{PelissierSpacetime2026}; the spectral construction supplies a compatible differential/spin realization and positive Hodge structures \cite{PelissierSpectral2026}; and the spacetime--gauge construction fixes the structural Standard-Model branch, including the physical gauge group, fermion representations, generation multiplicity, Higgs incidence, finite Dirac/Yukawa operator, and a common finite gravity--matter action \cite{PelissierDuality2026}. We use those results only to certify the finite common-action structure formalized in \Cref{def:finite-interface}. From that point onward, the continuum analysis is independent of the microscopic proofs in the companion papers. The present article asks whether such finite data close, in the continuum, into the coupled \emph{classical Einstein--Standard-Model dynamics}.

Almost-commutative spectral geometry produces gravitational and Standard-Model actions from a common spectral datum \cite{ChamseddineConnes1997,ChamseddineConnesMarcolli2007}. Our question is complementary: given the common finite action reconstructed from renewal data, which finite-cutoff estimates force a classical continuum limit and permit its complete first variation to pass to the limit? We prove a quantitative finite-action stability theorem on a regular coupled branch, a same-source selection theorem that derives the finite Euler-row, leakage and alignment budgets on one actually observed record from explicit source-level conditions, and a regulator-independent low-regularity closure theorem identifying every strong-packet accumulation point under compactness and physically normalized stationarity. The emphasis is therefore on analytic closure of the classical coupled equations and on finite, auditable entrances to that closure, not on rederiving the representation content and not on constructing a quantum functional integral.

The logical split is therefore
\begin{equation}
\begin{aligned}
\text{renewal histories}
&\longrightarrow
\text{finite Lorentzian--Standard-Model common action},\\
\text{finite action}+\text{analytic control}
&\longrightarrow
\text{classical Einstein--Standard-Model dynamics}.
\end{aligned}
\label{eq:two-layer-programme}
\end{equation}
The first arrow is supplied by the companion constructions and establishes realizability inside Renewal Geometry. The second is the mathematical content developed here. Its quantitative form is tied to an explicit local common action, while its low-regularity form is portable to any finite family satisfying the same structural and physical-variation assumptions.

Geometrically compatible discretizations provide important precedents. The curvature measures of piecewise-flat approximations were studied by Cheeger, M\"uller, and Schrader \cite{CheegerMullerSchrader1984}. Finite element exterior calculus organizes discrete differential forms through commuting complexes and stable projections \cite{ArnoldFalkWinther2006,ArnoldFalkWinther2010}. Christiansen and Halvorsen constructed a simplicial gauge action with exact discrete gauge invariance, consistency, and a discrete Noether theorem \cite{ChristiansenHalvorsen2012}. These results motivate the reconstruction interface used below. They do not imply that an arbitrary renewal complex has a shape-regular manifold realization. We therefore state that realization explicitly and separate it from action consistency.

Nonlinear compactness presents a second issue. Uhlenbeck's gauge compactness theory shows how curvature bounds and gauge choice interact \cite{Uhlenbeck1982}; it does not turn every bounded curvature sequence into a strongly convergent one. The critical loss of compactness for four-dimensional Higgs fields belongs to the concentration-compactness framework \cite{Lions1985}. In general relativity, high-frequency limits can also create effective sources. Burnett's formulation and the subsequent work of Huneau and Luk show why weakly convergent geometric data cannot simply be substituted into a nonlinear field equation \cite{Burnett1989,HuneauLuk2024}. Our argument excludes the relevant defects by explicit positive compactness conditions; it does not infer their absence from a formal Ward identity.

The variational geometry of spinors requires particular care because both the spinor bundle and the Dirac operator depend on the metric. We use a fixed local spin-frame identification, equivalently a local trivialization of the universal metric--spinor construction of M\"uller and Nowaczyk \cite{MullerNowaczyk2017}. The main theory is the torsion-free, second-order Einstein--Standard-Model system. An independently varied Palatini connection coupled to a nonzero spin current instead gives an Einstein--Cartan system; this alternative is not silently identified with torsion-free Einstein gravity \cite{Kibble1961,HehlEtAl1976}.

\paragraph{Main results.}
The first result is quantitative and is stated without using Renewal-specific algebraic terminology.  For the explicit finite gravity--Yang--Mills--Higgs--Dirac action written in \Cref{app:native-action}, \Cref{prop:native-consistency,thm:native-source,thm:native-closure} show that a small finite Euler row, compatible initial and harmonic-gauge data, and measured high-frequency leakage imply a differentiated physical forcing estimate and then cutoff-uniform control of the coupled actual jets.  In particular,
\begin{equation}
 \norm{\mathcal U_h-\mathcal U_*}_{C_tH_x^k}
 +\norm{\Riem(g_h)-\Riem(g_*)}_{L^2_tH_x^{k-1}}
 +\norm{T^{\mathrm{SM}}(z_h)-T^{\mathrm{SM}}(z_*)}_{L^2_tH_x^k}
 \le C D_{h,k}^{\rm full}.
 \label{eq:intro-quantitative}
\end{equation}
The all-time Sobolev bound and the stress convergence in this estimate are conclusions, not entrance assumptions.  The finite action itself has the off-shell consistency error $O(hK^3)$ on bounded-amplitude families with derivative scale $K$, and the proof-only Fourier splitting never filters the state entering the action or the conclusion.  The explicit refinement in \Cref{cor:native-rate} yields summable dyadic state, curvature and stress differences, so cofinal transport is derived on this branch.

The finite error budgets are not only treated as abstract admissibility conditions.  \Cref{thm:native-selected-closure} works on one literal accepted finite source family and selects one actually observed record on which the full finite Euler row, native reader energy and initial/gauge mismatch are simultaneously small.  The exact reader criterion of \Cref{prop:native-reader,cor:native-reader-sobolev} then converts the native positive energy into the same measured spacetime leakage used by the PDE theorem.  The resulting selected-record closure therefore derives the complete quantitative entrance on one record from source-level occupation, action-drift and reader conditions.  The strongly convex prefix-balance estimate of \Cref{thm:prefix-stationarity} remains as an alternative sufficient stationarity mechanism.  A fully finite spectral--midpoint branch is retained only to establish nonemptiness of the controlled class; its role is not to replace the native-source theorem by a numerical scheme.

The second result is the portable low-regularity closure theorem.  The compact-screen lemma, \Cref{lem:screen}, and its spacetime version, \Cref{prop:time-compactness}, give one route to strong convergence of the quadratic bosonic packets.  The reduced certificate of \Cref{thm:reduced-closure} removes two unnecessary endpoint assumptions: compactness of $D_{A_h}H_h$ together with $A_h\to A$ in $L^4$ forces the strong critical Higgs endpoint, while bounded weak $H^1$ spinors suffice for the complete bilinear fermionic first variation, including the metric--spin-connection contribution.  On a regular branch, \Cref{prop:sobolev-bosonic,prop:dirac-stability,thm:regular-branch} derive the stronger action topology directly from Sobolev and Lorentzian evolution estimates.

The general variational closure theorem, \Cref{thm:main-limit}, assumes no summable cutoff transport.  From compactness, first-variation consistency, and physically scaled common-action stationarity, every cutoff sequence has a subsequence whose reconstructed fields solve
\begin{equation}
 \Eins_{\mu\nu}(g)+\Lambda g_{\mu\nu}
 =\kappa T^{\mathrm{SM}}_{\mu\nu},
 \qquad \Evar_A=\Evar_H=\Evar_\Psi=\Evar_{\bPsi}=0,
 \label{eq:intro-system}
\end{equation}
in distributions. Operational stationarity is measured in physical test directions through
\begin{equation}
 \norm{\mathcal F_h}_{(\Vtest_K^r)^*}
 \le L_h\sqrt{R_h}+e_h,
 \label{eq:intro-source}
\end{equation}
so the relevant finite condition is $L_h\sqrt{R_h}+e_h\to0$, not an unscaled probability-space residual.  Summable adjacent-cutoff transport is retained only as the optional uniqueness upgrade of \Cref{cor:cofinal-unique}; on the quantitative finite-action branch it follows instead from the explicit rates.

When strong matter compactness fails, \Cref{thm:higgs-defect,prop:YM-defect,thm:joint-defect} retain the kinetic and concentration defects rather than normalizing them away.  Finally, \Cref{thm:hyperbolic,prop:subsidiary} separate variational convergence from reduced evolution and constraint propagation, while \Cref{prop:critical-shadowing} gives a complementary exact finite-stationary branch near nondegenerate smooth critical points.

\paragraph{Logical scope and analytic self-containedness.}
All analytic hypotheses used below are stated in this article. The companion papers are invoked only to establish the finite structural data and common-action normalization realized by Renewal Geometry; no continuum compactness, stress convergence, cofinal field convergence, or Euler equation is imported from them. The quantitative finite-action theorem derives the all-time regularity and transport conclusions from finite physical error budgets on its regular branch. The native selection theorem goes one step upstream: for one actual accepted source family it combines signed action drift, occupation control, a verified physical reader and initial/gauge alignment to produce one selected record to which the quantitative theorem applies. A gauge-covariant weak-sector reader criterion shows how the required tail control can be certified without imposing high derivatives on an arbitrary original gauge frame. The low-regularity theorem remains deliberately broader and therefore retains explicit no-concentration hypotheses. What remains outside the results is verification that a particular independently specified primitive Renewal source satisfies all of these source-level conditions on the full coupled gravity--gauge--chiral-matter packet. Universal microscopic attraction, universal manifold formation, and the interacting quantum Standard Model coupled to gravity are not claimed.

\begin{center}
\small
\begin{tabularx}{0.97\textwidth}{@{}L{0.31\textwidth}L{0.27\textwidth}X@{}}
\toprule
Ingredient & Status & Role in this paper\\
\midrule
Lorentzian carrier and gravitational normalization & established upstream & finite common-action structure\\
Standard-Model group, representations, generations, Higgs/Yukawa incidence & established upstream & finite common-action structure\\
Finite differential/spin realization and Hodge structure & established upstream & finite structural data and possible screen realization\\
Common finite gravity--matter action and source normalization & established upstream & finite common-action structure\\
First-variation consistency & abstract input in the general theorem; derived quantitatively on the local common-action branch & finite-to-continuum Euler control\\
Native stationarity & direct signed action-drift selection; strongly convex prefix balance as an alternative & physical Euler-row control on an observed record\\
Reader/leakage transfer & exact positive-form reader criterion; gauge-covariant weak-sector realization & source energy to physical spacetime tails\\
Finite record generation & finite constraint solver and spectral--midpoint branch & nonemptiness of the controlled class\\
Subsequential nonlinear compactness & screens or a-priori estimates; derived from finite forcing on the quantitative branch & continuum accumulation points\\
Cofinal uniqueness & optional transport in the general theorem; derived rates on the quantitative branch & full-sequence / route independence\\
Stationarity passage and stress defects & proved here & classical field equations / weak limits\\
Hyperbolic and constraint stability & proved here & evolution-based handoff\\
\bottomrule
\end{tabularx}
\end{center}

\Cref{sec:interface} specifies the common field, action, and test spaces. \Cref{sec:limit} states the quantitative finite-action theorem first and then the general low-regularity variational closure theorem. \Cref{sec:defects} treats concentration, hyperbolic stability, constraint propagation, and controlled realizations. \Cref{sec:discussion} gives the interpretation and remaining analytic tasks. The appendices contain reconstruction, Palatini, and nonminimal-coupling details.

\section{Finite common-action structure and physical variations}
\label{sec:interface}

The continuum analysis begins from finite Lorentzian--Standard-Model data carrying one common gravity--matter action rather than from the microscopic predictive construction itself. For portability we package these data in the formal finite-interface definition below. The companion papers establish one Renewal Geometry realization, while the analytic arguments use only the stated common-action, reconstruction and physical-variation properties.

\begin{definition}[Finite classical Einstein--Standard-Model interface]
\label{def:finite-interface}
At cutoff $h$, a finite classical Einstein--Standard-Model interface $\Iface_h$ consists of the following data.
\begin{enumerate}[label=\textup{(I\arabic*)},leftmargin=3.2em]
\item A finite configuration space $\mathcal Q_h$ with a reconstructed oriented, time-oriented Lorentzian carrier and local coframe variables, together with a fixed structural realization on a smooth comparison cylinder $M$ whenever the reconstruction map below is applied.
\item A fixed finite-rank Standard-Model bundle type with
\begin{equation}
 G_{\mathrm{SM}}=S(\U(3)\times\U(2))
 \cong[\SU(3)\times\SU(2)\times\U(1)]/\Z_6,
 \label{eq:gauge-group}
\end{equation}
the retained chiral fermion representations, one rank-three generation factor, the Higgs doublet, and any optional neutral/Majorana blocks declared as part of the branch.
\item A gauge-covariant finite Dirac/Yukawa incidence operator $D_{F,h}$ on that fixed structural carrier and a coefficient bank $\theta_h=(\kappa_h,\Lambda_h,g_{1,h},g_{2,h},g_{3,h},\lambda_{H,h},v_{H,h},\boldsymbol Y_h)$.
\item A differentiable common finite action
\[
 \Sact_h=\Sact_{\mathrm g,h}+\Sact_{\mathrm{SM},h}
\]
with one physical relative normalization between gravity and matter and with the corresponding finite gauge/relabeling covariance on the declared branch.
\item A positive source geometry for physical variations, including a faithful source quotient or equivalent positive Gram on the retained variation directions. This is used only through the physical test-lift estimates stated explicitly below; no algebraic commutant gap is treated as an analytic compactness estimate.
\end{enumerate}
No reconstruction-to-continuum map, compactness, stationarity, stress convergence, or field equation is included in this structural definition. Reconstruction maps and physical test lifts are introduced separately in \Cref{def:regulator}.
\end{definition}

\begin{proposition}[Renewal realization of the finite interface]
\label{prop:renewal-interface}
The three companion constructions realize \Cref{def:finite-interface} on the selected source-native Standard-Model branch \cite{PelissierSpacetime2026,PelissierSpectral2026,PelissierDuality2026}. In particular, they provide a nonempty family of finite interfaces with the structural bundle content and common-action normalization required here. This proposition supplies realizability only; none of the analytic hypotheses of \Cref{sec:limit} is inferred from it.
\end{proposition}
\begin{proof}
The Lorentzian carrier and gravitational action normalization are supplied by Ref.~\cite{PelissierSpacetime2026}; the compatible differential/spin realization is supplied by Ref.~\cite{PelissierSpectral2026}; and the structural Standard-Model bundle, generation carrier, finite Dirac/Yukawa incidence, common source quotient, and finite gravity--matter action are supplied by Ref.~\cite{PelissierDuality2026}. These are exactly the structural items \textup{(I1)}--\textup{(I5)}. The compactness, transport, consistency, and stationarity estimates used later are separate statements of the present paper.
\end{proof}

\begin{remark}[Structural provenance versus analytic input]
\label{rem:structural-provenance}
One realization of \textup{(I2)}--\textup{(I5)} comes from the upstream double-centralizer relation
\begin{equation}
 \A_{\mathrm{act},h}'=\M_{\mathrm{type},h},
 \qquad \M_{\mathrm{type},h}'=\A_{\mathrm{act},h},
 \label{eq:imported-duality}
\end{equation}
and from the common-source Schur quotient
\begin{equation}
 G_{\mathrm{i}\mid\mathrm g,h}
 =G_{\mathrm i,h}-C_hG_{\mathrm g,h}^{-1}C_h^*
 =S_{\mathrm i,h}^*(I-P_{\mathrm g,h})S_{\mathrm i,h}.
 \label{eq:source-schur}
\end{equation}
These identities explain why the finite matter packet and the physical source normalization are not appended by hand. They are not used to prove spatial or temporal compactness, convergence of $F_{A_h}$ or $D_{A_h}H_h$, spinor $H^1$ convergence, or any nonlinear stress limit. Those are analytic questions addressed independently below.
\end{remark}

\subsection{Spacetime, bundles, and test fields}

We work on an oriented, time-oriented spin cylinder
\begin{equation}
 M=(0,T)\times\Sigma,
 \qquad \partial\Sigma=\varnothing,
 \label{eq:cylinder}
\end{equation}
where $\Sigma$ is a compact smooth three-manifold. All conclusions are local on compact subcylinders. There is no physical timelike boundary in the main theorem, and test variations vanish near the initial and final time slices. A boundary extension requires its own flux and trace estimates; it is not implicit in \eqref{eq:cylinder}.

Choose a smooth auxiliary Riemannian metric $r$ and density $\dV_0$ on $M$. All $L^p$, Sobolev, and Hilbert-space norms refer to these positive comparison data and fixed smooth bundle metrics. Lorentzian contractions occur only in the signed action. On a compact nondegenerate coframe chart the different positive comparison norms are uniformly equivalent. No positivity of the Lorentzian action is assumed.

Let $P\to M$ be a fixed principal $G_{\mathrm{SM}}$ bundle of the type specified by \Cref{def:finite-interface}. The chiral fermion bundle, generation factor, and Higgs bundle $E_H$ are the corresponding associated bundles. We write
\begin{equation}
 z=(e,A,H,\Psi,\bPsi),\qquad
 g_{\mu\nu}=\eta_{ab}e^a_\mu e^b_\nu,
 \quad \eta=\diag(-1,1,1,1).
 \label{eq:fields}
\end{equation}
We use the curvature convention $[\nabla_\mu,\nabla_\nu]v^\rho=R^\rho{}_{\sigma\mu\nu}v^\sigma$ and $\Ric_{\sigma\nu}=R^\rho{}_{\sigma\rho\nu}$. The coframe and spinor expressions are local in a fixed finite atlas, with compatible frame changes on overlaps. A global coframe is not required. The fields $\Psi,\bPsi$ are classical realified spinor variables, varied independently before imposing the physical duality relation. They are not a Grassmann integration measure.

Put
\begin{equation}
 \Fcurv_A=\dd A+A\wedge A,\qquad
 K=D_AH=\dd H+\rho_H(A)H.
 \label{eq:curvature-higgs}
\end{equation}
The spin derivative $\nabla^{e,A}$ uses the Levi--Civita spin connection of $e$ and the retained gauge representation. The map $\mathscr M_{\boldsymbol Y}(H)$ is the finite Dirac/Yukawa coupling on the fixed fermion carrier. In real Higgs coordinates it is affine, with a possible constant Majorana block, and satisfies
\begin{equation}
 \norm{\mathscr M_{\boldsymbol Y}(H)}\le C_{\boldsymbol Y}(1+\abs H),
 \qquad
 \norm{D_H\mathscr M_{\boldsymbol Y}(H)[\eta_H]}
 \le C_{\boldsymbol Y}\abs{\eta_H}.
 \label{eq:yukawa-bound}
\end{equation}
Its covariance is part of the finite incidence interface. If the optional Majorana sector is retained, the fermion carrier includes its charge-conjugate blocks and the realified bilinear is understood on that doubled carrier with its fixed normalization.

\begin{definition}[Physical variation space]
\label{def:tests}
For a compact region $K\Subset M$, let $\Vtest_K$ consist of smooth compactly supported tuples
\[
 v=(k,a,\eta_H,\eta_\Psi,\eta_{\bPsi}),
 \qquad k\in\Dtest(M;\Sym^2 TM),\quad
 a\in\Dtest(M;T^*M\otimes\operatorname{ad}P),
\]
with the remaining entries in the indicated matter bundles. Thus $k^{\mu\nu}=\delta g^{\mu\nu}$. Fix an integer $r_0\ge4$ and give this space the sum of the $C^{r_0}$ norms in the chosen atlas. Its completion is denoted $\Vtest_K^{r_0}$. A determining core means a subspace dense in this test space, with the stated uniform bounds extending to its completion.
\end{definition}

The symmetric coframe lift of $k$ is
\begin{equation}
 \dot e^a_\mu(k)=-\tfrac12e^a_\alpha g_{\mu\beta}k^{\alpha\beta}.
 \label{eq:metric-lift}
\end{equation}
It gives $\dot g_{\mu\nu}=-g_{\mu\alpha}g_{\nu\beta}k^{\alpha\beta}$. Spinor components are transported in the same local spin-frame identification. Antisymmetric coframe variations are local Lorentz gauge directions, not additional metric tests.

\subsection{The classical action and the torsion convention}

Let $\theta=(\kappa,\Lambda,g_1,g_2,g_3,\lambda_H,v_H,\boldsymbol Y)$ be a fixed finite coefficient bank, with $\kappa,g_j,\lambda_H>0$. The matter action is
\begin{align}
 \Sact_{\mathrm{SM},\theta}(z)
 &=\int_M\mathcal L_{\mathrm{SM},\theta}(z)\,\dV_g,
 \label{eq:SM-action}\\
 \mathcal L_{\mathrm{SM},\theta}
 &=-\tfrac14\ip{\Fcurv_{\mu\nu}}{\Fcurv^{\mu\nu}}_{\boldsymbol g}
   -g^{\mu\nu}\Rea\ip{D_\mu H}{D_\nu H}
   -V(H)+\mathcal L_D,\nonumber\\
 V(H)&=\lambda_H(\abs H^2-v_H^2)^2,\nonumber\\
 \mathcal L_D
 &=\Rea\left\{\tfrac{\ii}{2}
 \bigl[\bPsi\gamma^\mu\nabla^{e,A}_\mu\Psi
       -(\nabla^{e,A}_\mu\bPsi)\gamma^\mu\Psi\bigr]
       -\bPsi\mathscr M_{\boldsymbol Y}(H)\Psi\right\}.
 \label{eq:dirac-density}
\end{align}
The positive invariant Lie-algebra metric in the first term has the three coefficients $g_j^{-2}$. Pairings in the fermionic expression include all retained chiral and generation indices. The gravitational action is
\begin{equation}
 \Sact_{\mathrm g,\theta}(g)
 =\frac{1}{2\kappa}\int_M(\Scal(g)-2\Lambda)\,\dV_g,
 \qquad \Sact_\theta=\Sact_{\mathrm g,\theta}+\Sact_{\mathrm{SM},\theta}.
 \label{eq:total-action}
\end{equation}
Only compactly supported variations are used. At limited regularity the gravitational first variation is defined by the first-derivative bulk expression obtained by integrating the torsion-free Palatini action by parts; \Cref{app:palatini} fixes this definition. It agrees with \eqref{eq:total-action} on smooth fields up to the appropriate boundary term.

\begin{remark}[Metric gravity versus independent Cartan gravity]
\label{rem:torsion}
In \eqref{eq:total-action} the spin connection is the torsion-free function $\omega(e)$, including in the matter action. This is a declared second-order branch. If instead $\omega$ is varied independently in the total Palatini--fermion action, its Euler equation contains the spin current. Eliminating that connection generally produces Einstein--Cartan interactions, not \eqref{eq:total-action}. The spinless Palatini elimination from the gravitational companion does not establish torsionlessness for a newly coupled spin-current sector. \Cref{app:palatini} retains the two alternatives explicitly.
\end{remark}

The Hilbert stress convention is
\begin{equation}
 D_g\Sact_{\mathrm{SM},\theta}(z)[k]
 =-\tfrac12\int_M T^{\mathrm{SM}}_{\mu\nu}(z)k^{\mu\nu}\,\dV_g.
 \label{eq:stress-definition}
\end{equation}
For smooth fields the gauge and Higgs contributions are
\begin{align}
 T^{\mathrm{YM}}_{\mu\nu}
 &=\ip{\Fcurv_{\mu\alpha}}{\Fcurv_\nu{}^\alpha}_{\boldsymbol g}
   -\tfrac14g_{\mu\nu}\ip{\Fcurv_{\alpha\beta}}{\Fcurv^{\alpha\beta}}_{\boldsymbol g},
 \label{eq:YM-stress}\\
 T^H_{\mu\nu}
 &=2\Rea\ip{D_\mu H}{D_\nu H}
   -g_{\mu\nu}\bigl(\abs{D_AH}_g^2+V(H)\bigr).
 \label{eq:H-stress}
\end{align}
The Dirac--Yukawa contribution is always taken from the complete metric variation of \eqref{eq:dirac-density}. In particular, the spin-connection variation is not dropped by holding a metric-dependent Dirac operator fixed.

\subsection{Finite reconstruction and first-variation consistency}

\begin{definition}[A finite-interface regulator sequence]
\label{def:regulator}
A finite-interface regulator sequence is a cofinal family $(\Iface_h)_{h\downarrow0}$ of structural interfaces from \Cref{def:finite-interface}, with finite configurations $z_h^{\mathrm d}\in\mathcal Q_h$, together with analytic reconstruction maps giving smooth local fields
\[
 z_h=\mathcal R_h(z_h^{\mathrm d})=(e_h,A_h,H_h,\Psi_h,\bPsi_h)
\]
on the same bundles over $M$, with the coefficient bank $\theta_h$ carried by $\Iface_h$. For every $K\Subset M$ there are linear physical test lifts $\mathcal I_h:\Vtest_K\to T_{z_h^{\mathrm d}}\mathcal Q_h$, supported in a fixed slightly larger region. Define
\begin{align}
 c_h(K)&=\sum_{b\in\{\mathrm g,\mathrm{SM}\}}\;\sup_{\norm v_{C^{r_0}}\le1}
 \abs{D\Sact_{b,h}(z_h^{\mathrm d})[\mathcal I_hv]
       -D\Sact_{b,\theta_h}(z_h)[v]},
 \label{eq:C1-consistency}\\
 \varepsilon_h(K)&=\sup_{\norm v_{C^{r_0}}\le1}
 \abs{D\Sact_h(z_h^{\mathrm d})[\mathcal I_hv]}.
 \label{eq:stationarity}
\end{align}
The supremum is over $\Vtest_K$. First-variation consistency means $c_h(K)\to0$. Physical common-action stationarity means $\varepsilon_h(K)\to0$.
\end{definition}

The comparison in \eqref{eq:C1-consistency} is with the action evaluated on the \emph{finite reconstructed fields}, not on their prospective limit. It can be checked by local interpolation, quadrature, and differentiated stencil estimates. It therefore does not assume the continuum Euler equation. \Cref{app:reconstruction} supplies a concrete smooth comparison realization. The existence of such estimates on a prescribed mesh must not be confused with their validity on every microscopic renewal complex.

Convergence of action values alone cannot replace \eqref{eq:C1-consistency}: $s_n(x)=n^{-1}\sin(nx)$ converges uniformly to zero, whereas $s_n'(0)=1$. Likewise, matching only normalized probabilities loses the affine action normalization needed for \eqref{eq:stationarity}.

\begin{proposition}[Scaled source control of physical stationarity]
\label{prop:source-bound}
Let $a_h$ be a finite action covector on a positive source space with Gram $G_h$. Suppose
\[
 \norm{a_h}_{G_h^{-1}}^2\le R_h,
 \qquad \norm{v_h(v)}_{G_h}\le L_h\norm v_{C^{r_0}},
\]
and
\[
 \abs{D\Sact_h[\mathcal I_hv]-a_h[v_h(v)]}
 \le e_h\norm v_{C^{r_0}}.
\]
Then $\varepsilon_h(K)\le L_h\sqrt{R_h}+e_h$. In particular, the sufficient vanishing condition is $L_h\sqrt{R_h}+e_h\to0$, not merely $R_h\to0$.
\end{proposition}
\begin{proof}
Write the source pairing in the $G_h$ metric and apply Cauchy--Schwarz:
\[
 \abs{a_h[v_h(v)]}
 \le\norm{a_h}_{G_h^{-1}}\norm{v_h(v)}_{G_h}
 \le L_h\sqrt{R_h}\norm v_{C^{r_0}}.
\]
Adding the identification error and taking the supremum proves the claim.
\end{proof}

The same estimate holds for an averaged source after Cauchy--Schwarz in the probability variable. An averaged Euler identity, however, is not a fieldwise equation for one limiting metric. The main theorem is deterministic; \Cref{cor:random-realization} states the additional requirements for a realized probabilistic version.

\section{Quantitative finite-action closure and general variational limits}
\label{sec:limit}

\subsection{Quantitative finite-action stability and continuum closure}
\label{sec:native-branch}

We begin with the quantitative theorem for an explicit finite common action.  Its hypotheses are finite-record quantities, and its principal compactness and cofinality conclusions are obtained from those error budgets rather than assumed.  The finite action used here is the local link/plaquette representative written explicitly in \Cref{app:native-action}; it is a finite sum and is not defined by sampling a continuum solution.  The proof uses a Fourier decomposition only as an estimate.  The action is always evaluated on the complete record, including all of its represented frequencies.

\begin{remark}[Intrinsic finite action and analytical comparison devices]
The finite action in this subsection is not introduced as a numerical surrogate chosen for its convergence properties.  It is the local link/plaquette realization of the same gravity--matter common-action structure used in the finite model: Cartan and gauge plaquettes, covariant Higgs links, spin--gauge transport, and the calibrated relative gravity/matter normalization are part of the finite action itself.  The auxiliary periodic box, trigonometric reconstruction and the projection $P_{\le K}$ used in the leakage proof are analytical comparison devices only; they neither alter the finite Euler covector nor replace the physical record entering the theorem.
\end{remark}

Fix buffered slabs $Q\Subset Q'$ inside a periodic auxiliary coordinate box $\mathscr Q$ of period $2\pi$ in each coordinate, with the auxiliary temporal seam outside $Q'$.  Let $n$ be odd, $h=2\pi/n$, and
\[
 \Lambda_h=\{\ell\in\Z^4:|\ell_\mu|\le(n-1)/2\}.
\]
Let $u_h$ be the calibrated nodal physical-coordinate record $(e,A,H,\Psi,\bPsi)$ in the open nondegenerate chart of \Cref{app:native-action}, and let $z_h=\mathcal I_h^{\rm trig}u_h$ be its trigonometric reconstruction.  For $1\le K<(n-1)/2$ define
\begin{equation}
 \tau_{h,j}(K)=\sum_{\substack{\ell\in\Lambda_h\\|\ell|_\infty>K}}
 \left(1+\frac{|\ell|_1}{K}\right)^j|\widehat u_h(\ell)|,
 \qquad j\ge0.
 \label{eq:native-tail}
\end{equation}
These are measured tails of the existing record; no Fourier coefficient is deleted.

For the finite action $S_h^{\rm loc}$ of \eqref{eq:native-local-action}, define its mass-normalized interior physical Euler row $E_h^{\rm raw}$ by
\begin{equation}
 D S_h^{\rm loc}(u_h)[v_h]
 =h^4\sum_{x\in Q'_h}\ip{E_h^{\rm raw}(u_h)(x)}{v_h(x)}
 \label{eq:raw-euler-row}
\end{equation}
for every interior nodal physical variation, with the fixed finite coefficient and representation metrics.  Write
\begin{equation}
 \norm{E_h^{\rm raw}(u_h)}_{0,h;Q'}\le \sigma_h.
 \label{eq:native-sigma}
\end{equation}
The quantity $\sigma_h$ is an actual finite-action covector norm, not a continuum residual.

\begin{proposition}[Growing-band consistency of the unchanged finite action]
\label{prop:native-consistency}
Let $y$ be a smooth physical field tuple on $Q'$ with uniformly bounded head amplitudes and inverse metric, and assume
\begin{equation}
 \norm{\partial^\alpha y}_\infty\le C_qK^{|\alpha|},
 \qquad |\alpha|\le q,\quad q\ge5,
 \qquad hK\le c_{\rm res}.
 \label{eq:growing-derivatives}
\end{equation}
Then the local finite action of \Cref{app:native-action} satisfies
\begin{equation}
 \max_{x\in Q'_h}\left|E_h^{\rm raw}(\mathsf S_hy)(x)
 -\mathsf S_h\mathcal E_0(y)(x)\right|\le ChK^3,
 \label{eq:native-consistency}
\end{equation}
where $\mathcal E_0$ is the complete minimally coupled torsion-free Einstein--Standard-Model Euler density in the same physical coordinates.  The literal Cartan plaquette curvature also obeys
\begin{equation}
 \norm{I_h^0R_h^{\rm raw}(\mathsf S_hy)-\Riem(g)}_{L^2(Q),+}
 \le ChK^3.
 \label{eq:native-Cartan}
\end{equation}
The constants depend on the bounded chart but not on $h$ or $K$.
\end{proposition}
\begin{proof}
The proof is given in \Cref{app:native-scaling}.  Its central point is an amplitude-preserving rescaling $\xi=Kx$, $\rho=hK$, $\nu=K^{-1}$.  The exact local density splits as $K^2$ times a uniformly regular bosonic stencil plus $K$ times a uniformly regular Dirac stencil.  First-variation consistency at normalized mesh $\rho$ is $O(\rho)$, giving $O(hK^3)$ after restoring the physical derivative factors.  No field amplitude is divided by $K$ and the metric--spinor chain is included.
\end{proof}

Set $m=k+2$ for an integer $k\ge4$ and define
\begin{align}
 \varepsilon_{0,h}&=\sigma_h+hK^3,
 &\varepsilon_{c,h}&=\varepsilon_{0,h}+K^2\tau_{h,2}(K),\notag\\
 L_{h,k}&=1+\log\left(2+\frac{C_kK^{k+5}}{\varepsilon_{c,h}}\right),
 &F_{h,k}&=\varepsilon_{c,h}K^{k+1}L_{h,k}^{k+1}
             +K^{k+2}\tau_{h,m}(K).
 \label{eq:native-forcing-budget}
\end{align}

\begin{theorem}[Leakage-tolerant physical forcing estimate]
\label{thm:native-source}
Assume $hK\le c_{\rm res}$, $\tau_{h,m}(K)\le\tau_*$, and $\varepsilon_{c,h}\le1$ on the fixed calibrated chart.  For the complete, unfiltered reconstruction $z_h=\mathcal I_h^{\rm trig}u_h$,
\begin{align}
 \epsilon^{\rm phys}(z_h)
 &:={\norm{\mathcal R_{\rm B}(z_h)}_{L^2(Q)}}
   +{\norm{\mathcal R_{\rm D}(z_h)}_{L^2(Q)}}
 \le C\varepsilon_{0,h},\label{eq:native-zero-source}\\
 \mathcal Y_k(z_h)
 &:={\norm{\mathcal R_{\rm B}(z_h)}_{L^2_tH_x^k(Q)}}
   +{\norm{\mathcal R_{\rm D}(z_h)}_{L^2_tH_x^{k+1}(Q)}}
 \le C F_{h,k},\label{eq:native-strong-source}\\
 \norm{I_h^0R_h^{\rm raw}(u_h)-\Riem(g_h)}_{L^2(Q),+}
 &\le ChK^3.\label{eq:native-full-Cartan}
\end{align}
Here $\mathcal R_{\rm B}$ collects the Einstein, Yang--Mills and Higgs physical Euler rows, while $\mathcal R_{\rm D}$ contains the Dirac and dual-Dirac rows.  No filtering is applied to the field entering these residuals or to the finite action.
\end{theorem}
\begin{proof}
\Cref{app:native-tails} proves the finite-order tail transfer.  The low-frequency proof field has the same chart margins and a complex spatial strip of width $a/K$; \Cref{lem:log-source-upgrade} upgrades its small $L^2$ source to the differentiated source norm with only the logarithmic factor in \eqref{eq:native-forcing-budget}.  The full measured tail is then restored at cost $K^{k+2}\tau_{h,m}$.  The zero-order estimate and Cartan estimate follow directly from \Cref{prop:native-consistency} applied to the full reconstructed field, whose first five derivatives satisfy \eqref{eq:growing-derivatives} by the tail bounds.
\end{proof}

Let $z_*$ be a regular torsion-free Einstein--Standard-Model solution on $Q=[0,T]\times\mathbb T^3$, written in harmonic coordinates and temporal internal gauge, and let $\mathcal U_*$ denote its actual-jet state from \Cref{app:coupled-stability}.  For an actual reconstructed record put
\begin{equation}
 i_{h,k}=\norm{\mathcal U_h(0)-\mathcal U_*(0)}_{H^k},
 \qquad
 \gamma_{h,k}=\norm{C(g_h)}_{L^2_tH_x^{k+1}}
              +\norm{\partial_tC(g_h)}_{L^2_tH_x^k},
 \label{eq:native-initial-gauge}
\end{equation}
where $C^\mu(g)=g^{\alpha\beta}\Gamma^\mu_{\alpha\beta}(g)$.  The state $\mathcal U_h$ is built from the actual derivatives and curvatures of $z_h$; no independent jet variables are supplied.

\begin{theorem}[Quantitative finite-action stability and Einstein--Standard-Model closure]
\label{thm:native-closure}
Let $k\ge4$.  Suppose the complete finite records exist on $Q$, stay in the fixed open field chart at the initial time, and satisfy
\[
 K_h\to\infty,\qquad hK_h\to0,\qquad
 \tau_{h,k+2}(K_h)\le\tau_*,
\]
together with the finite-action row bound \eqref{eq:native-sigma}.  Define
\begin{equation}
 D_{h,k}^{\rm full}=i_{h,k}+\gamma_{h,k}+F_{h,k}.
 \label{eq:native-total}
\end{equation}
If $D_{h,k}^{\rm full}\to0$, then the all-time $C_tH_x^k$ bound for the actual state is a conclusion, and
\begin{align}
 &\norm{\mathcal U_h-\mathcal U_*}_{C_tH_x^k}
 +\norm{\Riem(g_h)-\Riem(g_*)}_{L^2_tH_x^{k-1}}\notag\\
 &\hspace{4em}
 +\norm{T^{\rm SM}(z_h)-T^{\rm SM}(z_*)}_{L^2_tH_x^k}
 \le C D_{h,k}^{\rm full}.
 \label{eq:native-geometry}
\end{align}
In particular the compactness, stress convergence and zero-defect conclusions needed for the classical limit are derived on this branch from the finite physical row, initial/gauge mismatch and measured leakage.  Existence of the finite record on the slab is a premise; no existence or selection theorem for an unspecified microscopic rule is inferred.
\end{theorem}
\begin{proof}
\Cref{thm:native-source} provides the physical forcing of the complete record before any comparison with $z_*$.  Apply the one-sided coupled actual-jet stability estimate \Cref{prop:coupled-bootstrap} with
$d=i_{h,k}+\gamma_{h,k}+\mathcal Y_k(z_h)$.  A first-exit argument makes the temporary bounded-state hypothesis self-improving and yields \eqref{eq:native-geometry}.  The stress estimate follows because the complete torsion-free Standard-Model stress is a smooth algebraic function of the controlled actual jets, with the normal spinor derivative recovered from the Dirac equation.  Thus no separately convergent stress tensor or curvature sequence is assumed.
\end{proof}

\begin{corollary}[An explicit cofinal rate]
\label{cor:native-rate}
Suppose $\sigma_h=O(h)$, $K_h\asymp h^{-\beta}$ with $0<\beta<1/(k+4)$, and
\begin{equation}
 \tau_{h,2}=O(hK_h),\qquad
 \tau_{h,k+2}=O(hK_h^2).
 \label{eq:native-tail-rate}
\end{equation}
Then
\[
 F_{h,k}=O\!\left(h^{1-\beta(k+4)}(1+|\log h|)^{k+1}\right),
 \qquad
 \epsilon^{\rm phys}(z_h)=O(h^{1-3\beta}).
\]
For example,
\begin{equation}
 k=4,\qquad K_h\asymp h^{-1/18},\qquad
 i_{h,4}+\gamma_{h,4}=O(h^{1/2})
 \label{eq:native-example-rate}
\end{equation}
with \eqref{eq:native-tail-rate} gives $O(h^{1/2})$ state, curvature and Standard-Model stress convergence; the strong forcing is $O(h^{5/9}(1+|\log h|)^5)$ and the zero-order physical and Cartan residuals are $O(h^{5/6})$.  Along a dyadic sequence the adjacent state, curvature and stress differences are summable, so cofinal transport is a conclusion on this branch.
\end{corollary}
\begin{proof}
Here $\varepsilon_{c,h}=O(hK_h^3)$ and $L_{h,k}\le C(1+|\log h|)$.  Both terms in \eqref{eq:native-forcing-budget} are $O(hK_h^{k+4}(1+|\log h|)^{k+1})$.  Substitution of $k=4$ and $\beta=1/18$ gives the displayed powers.  Apply \Cref{thm:native-closure}; dyadic summability follows from the geometric series for any positive power of $h$ times a fixed logarithmic factor.
\end{proof}



\subsection{Native source selection, reader control, and same-record closure}
\label{sec:native-selection}

The finite-action theorem above is deterministic and applies to any complete record satisfying the displayed physical budgets.  We now give a finite source-level entrance that produces those budgets on one \emph{actually observed} record.  The construction does not replace a record by its expectation, by a projected state, or by a trajectory of an inserted evolution equation.

At spacing $h$, let $X_0,\ldots,X_J$ be the records of an actual finite accepted process.  The state may retain the entire finite past, so no Markov assumption is required.  Let $K^{\rm leg}_{h,j}(x,y)$ denote the transition probabilities restricted by a predeclared legal-edge mask for the raw-action chart and action identification; an attempted illegal transition is counted as failure and the diagnostic stopped process retains its last legal action value.  Write $\mu_j$ for the resulting unconditioned surviving mass and $p_h^{\rm exit}=1-\sum_x\mu_J(x)$.  Every legal state has a complete physical record $u_h(x)$ and define
\begin{equation}
 g_h(x)=\norm{E_h^{\rm raw}(u_h(x))}_{0,h},\qquad
 W_h(x)=\bigl(i_{h,k}(x)+\gamma_{h,k}(x)\bigr)^2,
 \label{eq:source-gW}
\end{equation}
together with a nonnegative native energy readout $V_h(x)$.  Put
\begin{equation}
 M_{h,J}=\frac1J\sum_{j<J,x}\mu_j(x)V_h(x),\qquad
 A_{h,J}=\frac1J\sum_{j<J,x}\mu_j(x)W_h(x).
 \label{eq:source-occupations}
\end{equation}
Let $a_h(x)=S_h^{\rm loc}(u_h(x))$ and suppose $A_h^-\le a_h(x)\le A_h^+$ on the legal chart, with action span $D_h^A=A_h^+-A_h^-$.  The stopped conditional action increment is
\begin{equation}
 \mathfrak d_{h,j}(x)=\sum_yK^{\rm leg}_{h,j}(x,y)\,[a_h(y)-a_h(x)].
 \label{eq:source-action-drift}
\end{equation}

\begin{theorem}[Same-record selection from the actual signed action drift]
\label{thm:source-selection}
Assume there are $\delta_h>0$, $\kappa_A>0$ independent of $h$, and nonnegative finite remainders $r_{h,j}$ such that every active legal row satisfies
\begin{equation}
 \mathfrak d_{h,j}(x)
 \le-\kappa_A\delta_h g_h(x)^2+\delta_h r_{h,j}(x).
 \label{eq:source-drift-ineq}
\end{equation}
Set
\begin{equation}
 \mathcal R_{h,J}=\frac1J\sum_{j<J,x}\mu_j(x)r_{h,j}(x).
 \label{eq:source-work-occupation}
\end{equation}
Then
\begin{equation}
 \frac1J\sum_{j<J,x}\mu_j(x)g_h(x)^2
 \le \frac{D_h^A}{\kappa_A\delta_hJ}
       +\frac{\mathcal R_{h,J}}{\kappa_A}.
 \label{eq:source-slope-occupation}
\end{equation}
For thresholds $s_h,R_h,\alpha_h>0$, select on every surviving prefix an observed index minimizing
\begin{equation}
 \frac{g_h(X_j)^2}{s_h^2}
 +\frac{V_h(X_j)}{R_h^2}
 +\frac{W_h(X_j)}{\alpha_h^2}.
 \label{eq:source-score}
\end{equation}
The probability of failing to output one observed record with
\begin{equation}
 g_h\le s_h,\qquad V_h\le R_h^2,\qquad
 i_{h,k}+\gamma_{h,k}\le\alpha_h
 \label{eq:source-selected}
\end{equation}
is at most
\begin{equation}
 \boxed{
 p_h^{\rm exit}
 +\frac{D_h^A}{\kappa_A\delta_hJ s_h^2}
 +\frac{\mathcal R_{h,J}}{\kappa_A s_h^2}
 +\frac{M_{h,J}}{R_h^2}
 +\frac{A_{h,J}}{\alpha_h^2}.}
 \label{eq:source-selection-failure}
\end{equation}
No invariant distribution, mixing estimate, independence between records, or uniform high-order bound on every member of the prefix is required.
\end{theorem}
\begin{proof}
The stopped-action telescope and the deterministic score argument are proved in \Cref{app:native-source-selection}.  The key point is that killing is recorded as failure rather than being credited as a decrease of the signed action.
\end{proof}

The native energy must control the \emph{physical reader actually entering} \eqref{eq:native-tail}.  This relation has an exact finite-dimensional criterion.

\begin{proposition}[Exact native-energy to physical-tail certificate]
\label{prop:native-reader}
Let $\mathsf H_h^{\rm src}$ be a finite real source-coordinate Hilbert space, let $\mathsf Q_h\succeq0$ be a positive native energy form, and let $\xi_h(x)\in\mathsf H_h^{\rm src}$ be the source coordinate of a legal record.  Suppose the physical reader has the affine form
\begin{equation}
 u_h(x)=u_h^{\rm soft}+\mathsf R_h\xi_h(x).
 \label{eq:source-reader}
\end{equation}
For the weighted high-frequency coefficient map $\mathsf T_{h,K,m}$ associated with \eqref{eq:native-tail}, a finite constant $c_{h,K,m}$ satisfying
\begin{equation}
 \tau_{h,m}(K;u_h(x))
 \le \tau_{h,m}(K;u_h^{\rm soft})
      +c_{h,K,m}\,\ip{\xi_h(x)}{\mathsf Q_h\xi_h(x)}^{1/2}
 \label{eq:source-reader-tail}
\end{equation}
exists if and only if
\begin{equation}
 \mathsf T_{h,K,m}\ker\mathsf Q_h=0.
 \label{eq:source-reader-kernel}
\end{equation}
When this holds, the optimal constant is
\begin{equation}
 c_{h,K,m}
 =\norm{\mathsf T_{h,K,m}\mathsf Q_h^{\dagger/2}}_{2\to\ell^1(w)}.
 \label{eq:source-reader-constant}
\end{equation}
Thus the leakage assumption is a finite, directly auditable incidence condition between the native energy and the declared physical reader; finite-dimensional norm equivalence alone is not used.
\end{proposition}
\begin{proof}
See \Cref{app:native-reader-proof}.
\end{proof}

\begin{corollary}[Sobolev positive-form reader criterion]
\label{cor:native-reader-sobolev}
Let $\mathsf L_h$ be the diagonal multiplier $(1+|\ell|_2^2)^{1/2}$ on the all-mode Fourier reader.  If, on the retained source coordinates,
\begin{equation}
 \mathsf R_h^*\mathsf L_h^{2q}\mathsf R_h\preceq C_R\mathsf Q_h,
 \qquad \norm{u_h^{\rm soft}}_{H^q}\le C_R,
 \label{eq:source-reader-LMI}
\end{equation}
then for every $0\le j\le m<q-2$,
\begin{equation}
 \norm{u_h(x)}_{H^q}\le C\bigl(1+\sqrt{V_h(x)}\bigr),\qquad
 \tau_{h,j}(K;u_h(x))
 \le C K^{2-q}\bigl(1+\sqrt{V_h(x)}\bigr),
 \label{eq:source-reader-Sobolev}
\end{equation}
whenever $\ip{\xi_h}{\mathsf Q_h\xi_h}\le V_h$.  The power $K^{2-q}$ is sharp for this four-dimensional Sobolev-to-weighted-tail route; in particular the composition with the $k$th-order forcing estimate requires $q>k+5$ if no additional cancellation is used.
\end{corollary}
\begin{proof}
The positive-form inequality gives the $H^q$ estimate.  The sharp four-dimensional weighted Fourier bound is proved in \Cref{app:native-reader-proof}; inserting $m=k+2$ into the forcing budget yields the stated threshold.
\end{proof}

\begin{theorem}[Same-source selected finite-action Einstein--Standard-Model closure]
\label{thm:native-selected-closure}
Assume the hypotheses of \Cref{thm:source-selection}.  Suppose the same legal records and the same native energy satisfy, for $j=2,m=k+2$,
\begin{equation}
 \tau_{h,j}(K_h;u_h(x))\le t_{h,j}+c_{h,j}\sqrt{V_h(x)},
 \label{eq:selected-reader-general}
\end{equation}
with certified finite reader constants.  Put $T_{h,j}=t_{h,j}+c_{h,j}R_h$ and define
\begin{align}
 \bar\varepsilon_{c,h}&=s_h+hK_h^3+K_h^2T_{h,2},\notag\\
 \bar L_{h,k}&=1+\log\left(2+\frac{C_kK_h^{k+5}}{\bar\varepsilon_{c,h}}\right),\notag\\
 \bar F_{h,k}&=\bar\varepsilon_{c,h}K_h^{k+1}\bar L_{h,k}^{k+1}
                   +K_h^{k+2}T_{h,m},\qquad
 \bar D_{h,k}&=\alpha_h+\bar F_{h,k}.
 \label{eq:selected-composed-budget}
\end{align}
Assume $K_h\to\infty$, $hK_h\to0$, $T_{h,m}\to0$, $\bar\varepsilon_{c,h}\to0$ and $\bar D_{h,k}\to0$.  Outside the event in \eqref{eq:source-selection-failure}, the selected actually observed record satisfies
\begin{align}
 &\norm{\mathcal U_{h,j_*}-\mathcal U_*}_{C_tH_x^k}
 +\norm{\Riem(g_{h,j_*})-\Riem(g_*)}_{L^2_tH_x^{k-1}}\notag\\
 &\hspace{4em}
 +\norm{T^{\rm SM}(z_{h,j_*})-T^{\rm SM}(z_*)}_{L^2_tH_x^k}
 \le C\bar D_{h,k}.
 \label{eq:selected-closure}
\end{align}
Its zero-order physical residual is at most $C(s_h+hK_h^3)$ and its literal Cartan discrepancy is at most $ChK_h^3$.  If the failure bounds and $\bar D_{h_n,k}$ are summable along a common cofinal chain, then under any common coupling the selected successful records eventually have summable adjacent state, curvature and stress differences almost surely.  No independence between cutoffs is required.
\end{theorem}
\begin{proof}
The selector supplies on the same observed record $\sigma_h\le s_h$, $i_{h,k}+\gamma_{h,k}\le\alpha_h$, and $V_h\le R_h^2$.  Equation \eqref{eq:selected-reader-general} then gives the certified tail bounds $\tau_{h,j}\le T_{h,j}$.  These are precisely admissible upper budgets in \Cref{thm:native-source,thm:native-closure}, proving \eqref{eq:selected-closure}.  Summability follows from the first Borel--Cantelli lemma and comparison of adjacent successful records through the same regular physical solution.
\end{proof}

\begin{corollary}[An explicit polynomial selected-source regime]
\label{cor:selected-polynomial}
Fix $k\ge4$, put $m=k+2$, and suppose \Cref{cor:native-reader-sobolev} holds for some $q>k+5$.  Assume, for positive $a,b,c,\epsilon$, $d,d_A\ge0$,
\begin{align}
 \delta_h&\ge c_\delta h^d,& D_h^A&\le Ch^{-d_A},\notag\\
 J_h&\ge c_Jh^{-(d+d_A+2b+2\epsilon)},&
 \mathcal R_{h,J_h}&\le Ch^{2b+2\epsilon},\notag\\
 M_{h,J_h}&\le C,& A_{h,J_h}&\le Ch^{2c+2\epsilon},\notag\\
 p_h^{\rm exit}&\le Ch^{2\epsilon}.&&
 \label{eq:selected-polynomial-source}
\end{align}
Choose
\begin{equation}
 s_h=h^b,\qquad R_h=h^{-a},\qquad \alpha_h=h^c,
 \qquad K_h\asymp h^{-\beta},
 \label{eq:selected-polynomial-thresholds}
\end{equation}
where
\begin{equation}
 0<\beta<\min\left\{\frac{b}{k+1},\frac1{k+4}\right\},
 \qquad 0<a<\beta(q-k-5).
 \label{eq:selected-polynomial-window}
\end{equation}
Then the selection failure probability is $O(h^{2a}+h^{2\epsilon})$ and, on success,
\begin{equation}
 \text{left side of \eqref{eq:selected-closure}}
 \le C h^\rho(1+|\log h|)^{k+1},
 \quad
 \rho=\min\{c,b-\beta(k+1),1-\beta(k+4),\beta(q-k-5)-a\}>0.
 \label{eq:selected-polynomial-rate}
\end{equation}
Along $h_n\asymp2^{-n}$ the successful selected branch is eventually almost sure and has summable adjacent differences whenever the displayed probability bound is summable.
\end{corollary}
\begin{proof}
See \Cref{app:native-polynomial-proof}.
\end{proof}

\begin{corollary}[Gauge-robust weak-sector reader entrance]
\label{cor:gauge-robust-reader}
On a nonvanishing Higgs chart of the weak $SU(2)$ sector, a cutoff-uniform bound on the finite covariant-jet energy of \Cref{app:gauge-reader} yields, after algebraic Higgs normalization and a causal finite temporal gauge, a bounded all-mode $H^q$ reader and hence the tail estimate
\begin{equation}
 \tau_{h,m}(K)\le C K^{2-q},\qquad q>m+2,
 \label{eq:gauge-robust-tail}
\end{equation}
without requiring high ordinary derivatives of the original site frame or small original links.  Exact finite gauge invariance preserves the norm of the \emph{full} physical action covector under this normalization, and logarithmic link coordinates are uniformly equivalent to the physical link-tangent norm on the normalized chart.  This result supplies one source-auditable realization of the reader hypothesis for the weak sector; color, determinant-line, coframe and other independent reader components require their own corresponding bounds.
\end{corollary}
\begin{proof}
The finite covariant-jet, temporal-normalization and full-covector identities are proved in \Cref{app:gauge-reader}.
\end{proof}

\paragraph{Alternative same-action prefix balance.}
The direct signed action-drift entrance above is not the only way to obtain a small native Euler row.  A strongly convex operational potential gives the following independent sufficient criterion.  It is useful when its force-balance discrepancy is directly available, but is not imposed on source laws for which action drift can be verified more economically.

Give the free physical nodal variables the mass norm
\begin{equation}
 \norm v_{0,h}^2=h^4\sum_{x\in Q'_h}|v(x)|^2.
 \label{eq:prefix-mass}
\end{equation}
Let $\mathcal A_h=S_h^{\rm loc}$ on a compact convex field chart $\Theta_h$.  Suppose a positive operational potential $\Psi_h$ has
\begin{equation}
 g_*I\preceq \nabla_h^2\Psi_h(y)\preceq g^*I,
 \qquad 0<g_*\le g^*<\infty,
 \label{eq:prefix-Fisher}
\end{equation}
and the local action obeys
\begin{equation}
 \norm{\nabla_h^2\mathcal A_h(y)}_{0,h\to0,h}\le K_0h^{-2},
 \qquad \operatorname{osc}_{\Theta_h}\mathcal A_h\le D_0.
 \label{eq:prefix-Hessian}
\end{equation}

\begin{theorem}[Alternative finite-prefix stationarity criterion]
\label{thm:prefix-stationarity}
Let $\theta_0,\ldots,\theta_J\in\Theta_h$ be an accepted finite sequence whose connecting segments stay in the chart.  Put $E_h=\nabla_h\mathcal A_h$, choose $\upsilon_h=\upsilon_0h^2$ with $\upsilon_0K_0<g_*/2$, and define
\begin{equation}
 b_j=E_h(\theta_{j+1})
 +\upsilon_h^{-1}\{\nabla_h\Psi_h(\theta_{j+1})-\nabla_h\Psi_h(\theta_j)\}.
 \label{eq:prefix-balance}
\end{equation}
Then
\begin{equation}
 \frac1J\sum_{j=0}^{J-1}\norm{E_h(\theta_j)}_{0,h}^2
 \le C_0h^{-2}\frac{D_0}{J}
 +D_1\frac1J\sum_{j=0}^{J-1}\norm{b_j}_{0,h}^2.
 \label{eq:prefix-average}
\end{equation}
Consequently the smallest actually observed full native slope satisfies
\begin{equation}
 \norm{E_h(\theta_{j_*})}_{0,h}
 \le C\left[h^{-1}\sqrt{D_0/J}
 +\left(J^{-1}\sum_{j<J}\norm{b_j}_{0,h}^2\right)^{1/2}\right].
 \label{eq:prefix-selected}
\end{equation}
In particular $D_0/J=O(h^4)$ and $J^{-1}\sum\norm b_j\norm^2=O(h^2)$ give $\sigma_h=O(h)$ without an invariant probability distribution or an assumed small selected action covector.
\end{theorem}
\begin{proof}
The same-action descent calculation is given in \Cref{app:prefix-stationarity}.
\end{proof}

\begin{corollary}[Finite-table certification for the prefix-balance entrance]
\label{cor:prefix-table}
Suppose the accepted transitions are represented by finite tables, a safe set carries the chart and balance hypotheses above, and $p_h^{\rm exit}$ is the probability of leaving that safe set before $J$ accepted comparisons.  Put
\begin{equation}
 \mathcal Q_{h,J}^{\rm bal}
 =\frac1J\sum_{j<J}\mathbb E\!\left[
   \mathbf 1_{\rm safe}\norm{b_j}_{0,h}^2\right].
 \label{eq:prefix-table-balance}
\end{equation}
For any $0<\alpha\le1$, the probability that the best safe observed record fails a fixed-constant $O(h^\alpha)$ full native-row bound is at most
\begin{equation}
 p_h^{\rm exit}
 +C\frac{D_0}{J h^{2+2\alpha}}
 +C h^{-2\alpha}\mathcal Q_{h,J}^{\rm bal}.
 \label{eq:prefix-table-probability}
\end{equation}
If this bound is summable along a dyadic cutoff family, eventual certification holds almost surely without an independence assumption between cutoffs.
\end{corollary}
\begin{proof}
See \Cref{app:prefix-stationarity}.
\end{proof}

\begin{corollary}[Nonemptiness of the controlled finite class]
\label{cor:generated-nonempty}
The quantitative hypotheses above are nonempty.  For the constraint-compatible initial class constructed in \Cref{app:generated-dynamics}, the fully finite spectral--midpoint evolution there exists on a common local slab and produces records $z_{N,\tau}$ satisfying
\begin{equation}
 \norm{z_{N,\tau}-z}_{C_tH^{s+2}}
 +\norm{\Riem(g_{N,\tau})-\Riem(g)}_{L^\infty_tH^s}
 +\norm{\mathcal R_{\rm B}(z_{N,\tau})}_{L^\infty_tH^s}
 \le C(N^{-p}+\tau),
 \label{eq:generated-nonempty}
\end{equation}
with the Dirac residual $O(N^{-p}+\tau^2)$ and summable dyadic transport.  The construction starts from finite constrained initial data and does not sample a prescribed future solution.  Its purpose here is to establish nonemptiness of the controlled class; the source-native theorem is \Cref{thm:native-selected-closure}.
\end{corollary}
\begin{proof}
This is the content of \Cref{app:generated-dynamics}.
\end{proof}

\begin{remark}[What remains outside the quantitative branch]
\label{rem:native-boundary}
For one actual accepted finite law, \Cref{thm:native-selected-closure} derives the selected Euler-row, leakage and initial/gauge budgets jointly from explicit source-level conditions.  What is not proved is that an arbitrary independently specified primitive Renewal source satisfies the required signed action-drift, reader-energy, legal-domain and full coupled gravity--gauge--chiral-matter conditions.  The later source programme reduces several of those questions to finite range and phase-space tests, but the complete primitive-source Einstein--Standard-Model implication is not assumed here.
\end{remark}

\subsection{Compact screens for the quantities that enter the action}

\begin{lemma}[Common compact screen]
\label{lem:screen}
Let $Y_h\rightharpoonup Y$ in a fixed Hilbert space $\mathcal Y$. Suppose that, for every $R$, there are bounded $S_{h,R}$ and compact $S_R$ with
\begin{equation}
 \norm{S_{h,R}-S_R}_{\mathrm{op}}\longrightarrow0,
 \qquad
 \lim_{R\to\infty}\sup_h\norm{(I-S_{h,R})Y_h}=0,
 \qquad S_RY\longrightarrow Y.
 \label{eq:screen}
\end{equation}
Then $Y_h\to Y$ strongly. The same conclusion holds for a whitened packet $Z_h=\mathbb A_h^{1/2}Y_h$ provided $Z_h\rightharpoonup\mathbb A^{1/2}Y$, the screen conditions hold for $Z_h$, and
\[
 mI\preceq\mathbb A_h\preceq MI,
 \qquad \mathbb A_h^{-1/2}\to\mathbb A^{-1/2}\text{ strongly},
 \quad 0<m\le M<\infty.
\]
\end{lemma}
\begin{proof}
For fixed $R$, compactness gives $S_RY_h\to S_RY$, and norm convergence of $S_{h,R}$ gives $(S_{h,R}-S_R)Y_h\to0$. Decompose
\[
 Y_h-Y=(I-S_{h,R})Y_h+(S_{h,R}Y_h-S_RY)+(S_R-I)Y.
\]
Take the upper limit in $h$ and then let $R\to\infty$. This proves the first assertion. For the second, apply the first assertion to $Z_h$ and use
\[
 \norm{Y_h-Y}\le m^{-1/2}\norm{Z_h-Z}
   +\norm{(\mathbb A_h^{-1/2}-\mathbb A^{-1/2})Z}.
\]
\end{proof}

A screen for the action must include the full quadratic packet. For example, $\Fcurv_{A_h}$, $D_{A_h}H_h$, and the first spinor jets are different components, even when all are represented on the same finite history. A gap in the multiplicity commutant does not control their spatial or temporal frequencies.

\begin{proposition}[From spatial screens to spacetime compactness]
\label{prop:time-compactness}
Let $\mathcal H=L^2(\Sigma;E)$ and let $P_R$ be finite-rank orthogonal projections with smooth ranges and $P_R\to I$ strongly. Suppose $Y_h$ is bounded in $L^2(0,T;\mathcal H)$,
\[
 \lim_{R\to\infty}\sup_h\norm{(I-P_R)Y_h}_{L^2_t\mathcal H}=0,
\]
and, for each fixed $R$, $P_RY_h$ is bounded in $H^1(0,T;P_R\mathcal H)$. Then $(Y_h)$ is precompact in $L^2(0,T;\mathcal H)$. The last condition follows, for example, from a bound on $\partial_tY_h$ in $L^2(0,T;H^{-m}(\Sigma;E))$ for a fixed $m$.
\end{proposition}
\begin{proof}
For fixed $R$, the range is finite dimensional, so one-dimensional compact Sobolev embedding makes $(P_RY_h)$ precompact in $L^2_t\mathcal H$. The tail estimate gives a finite $\epsilon$-net for $(Y_h)$ by first choosing $R$ and then a finite net for the projected sequence. For the final assertion, pair $\partial_tY_h$ with a smooth orthonormal basis of $P_R\mathcal H$; each coefficient has a bounded $L^2$ time derivative.
\end{proof}

This elementary special case of spacetime compactness is consistent with the general theory of Simon \cite{Simon1986}. The time condition is indispensable: $Y_h(t,x)=\sin(t/h)v(x)$ remains in a single spatial mode but has no strongly convergent subsequence in spacetime $L^2$. Thus compact spatial Hodge screens alone do not establish the required spacetime result.

The Higgs potential requires a different endpoint norm. On a four-dimensional bounded region, an $H^1$ bound controls $L^4$ but does not make its unit ball compact in $L^4$.

\begin{proposition}[A local quartic compactness certificate]
\label{prop:orlicz}
Let $H_h\rightharpoonup H$ in $H^1(K)$ on a bounded smooth four-dimensional region. Suppose $\Phi:[0,\infty)\to[0,\infty)$ is increasing, satisfies $\Phi(t)/t\to\infty$, and
\begin{equation}
 \sup_h\int_K\Phi(\abs{H_h}^4)\,\dV_0<\infty.
 \label{eq:orlicz}
\end{equation}
Then $H_h\to H$ strongly in $L^4(K)$. A uniform $L^{4+\delta}$ bound, $\delta>0$, is sufficient.
\end{proposition}
\begin{proof}
Rellich compactness gives strong $L^2$ convergence and convergence in measure. For any $B>0$,
\[
 \int_{\{\abs{H_h}^4>B\}}\abs{H_h}^4
 \le\sup_{t>B}\frac{t}{\Phi(t)}\int_K\Phi(\abs{H_h}^4).
\]
The right side tends to zero uniformly in $h$. Hence the fourth powers are uniformly integrable. Convergence in measure and this uniform integrability give $\norm{H_h-H}_4\to0$ by the elementary Vitali argument, applied after using $\abs{H_h-H}^4\le8(\abs{H_h}^4+\abs H^4)$. Take $\Phi(t)=t^{1+\delta/4}$ for the last assertion.
\end{proof}

\subsection{A-priori routes to the action topology}

The compact-screen formulation is useful at low regularity, but it should not be read as the only way to obtain the nonlinear action topology. On a regular branch, ordinary a-priori Sobolev estimates imply the required bosonic compactness directly.

\begin{proposition}[A-priori Sobolev bounds imply bosonic strong compactness]
\label{prop:sobolev-bosonic}
Let $K\Subset M$ be a bounded smooth four-dimensional region. Suppose that for some $\sigma>0$, after fixed local gauge and frame choices,
\begin{equation}
 \sup_h\Bigl(
 \norm{e_h}_{H^{2+\sigma}(K)}
 +\norm{A_h}_{H^{1+\sigma}(K)}
 +\norm{H_h}_{H^{1+\sigma}(K)}
 \Bigr)<\infty,
 \qquad
 \sup_h\norm{e_h^{-1}}_{L^\infty(K)}<\infty.
 \label{eq:apriori-bosonic}
\end{equation}
Assume also that the coefficient banks remain in a compact subset of the physical parameter region. Then every cutoff sequence has a subsequence such that
\begin{align}
 e_h&\to e &&\text{strongly in }H^1(K)\cap L^\infty(K),\notag\\
 A_h&\to A &&\text{strongly in }H^1(K)\text{ and hence in }L^4(K),\notag\\
 H_h&\to H &&\text{strongly in }H^1(K)\cap L^4(K),\notag\\
 \Fcurv_{A_h}&\to\Fcurv_A &&\text{strongly in }L^2(K),\notag\\
 D_{A_h}H_h&\to D_AH &&\text{strongly in }L^2(K).
 \label{eq:apriori-bosonic-convergence}
\end{align}
Moreover $H_h$ is uniformly bounded in $L^{4+\delta}(K)$ for some $\delta>0$ depending only on $\sigma$ and the dimension, so quartic uniform integrability follows automatically. Thus, on this regular branch, the bosonic part of the strong packet is obtained from a-priori boundedness rather than assumed precompactness or screen tails.
\end{proposition}
\begin{proof}
Rellich--Kondrachov gives compact embeddings $H^{2+\sigma}(K)\Subset H^1(K)$ and $H^{1+\sigma}(K)\Subset H^1(K)$. Since $2+\sigma>2=d/2$, the first family is also precompact in $C^0(K)$ after lowering the H\"older exponent, hence in $L^\infty(K)$. The uniform inverse bound preserves the nondegenerate signature chart. Strong $H^1$ convergence of $A_h$ gives strong $L^4$ convergence in four dimensions. Therefore
\[
 \dd A_h\to\dd A\quad\text{in }L^2,
 \qquad
 A_h\wedge A_h\to A\wedge A\quad\text{in }L^2,
\]
which proves the curvature convergence. The same argument, using
\[
 D_{A_h}H_h=\dd H_h+\rho_H(A_h)H_h,
\]
gives strong $L^2$ convergence of the Higgs covariant derivative. Finally $H^{1+\sigma}$ embeds continuously into some $L^{4+\delta}$ with $\delta>0$, yielding quartic uniform integrability. Finite-dimensional compactness handles the coefficients.
\end{proof}

\begin{proposition}[Lorentzian Dirac stability produces spinor $H^1$ compactness]
\label{prop:dirac-stability}
Let $I\times\Sigma$ be a compact subslab with $\Sigma$ closed. In a common local spin trivialization suppose the reconstructed spinors satisfy symmetric-hyperbolic systems
\begin{equation}
 A_h^0\partial_t\Psi_h+A_h^j\partial_j\Psi_h+B_h\Psi_h=r_h,
 \label{eq:dirac-hyperbolic}
\end{equation}
with the analogous dual system for $\bPsi_h$. Assume the matrices $A_h^0$ are uniformly positive definite, the systems are uniformly symmetric hyperbolic, and
\begin{equation}
 \sup_h\bigl(
 \norm{A_h^\mu}_{W^{1,\infty}(I\times\Sigma)}
 +\norm{B_h}_{L^\infty_{t,x}}
 +\norm{\Psi_h}_{L^\infty_tH^2_x}
 +\norm{\bPsi_h}_{L^\infty_tH^2_x}
 \bigr)<\infty.
 \label{eq:dirac-uniform}
\end{equation}
Suppose the coefficients are Cauchy in $L^\infty(I\times\Sigma)$, the spinor and dual-spinor initial data are Cauchy in $L^2(\Sigma)$, and the residuals $r_h,\bar r_h$ are Cauchy in $L^2_tL^2_x$. Then $\Psi_h$ and $\bPsi_h$ are Cauchy in $H^1(I\times\Sigma)$. In particular, if $r_h,\bar r_h\to0$, the limiting fields satisfy the Lorentzian Dirac--Yukawa and dual equations determined by the limiting coefficients.
\end{proposition}
\begin{proof}
For $w=\Psi_h-\Psi_j$, subtract the two first-order systems and use the $h$-system as the principal symmetric-hyperbolic operator. Uniform positivity of $A_h^0$ and the coefficient bounds in \eqref{eq:dirac-uniform} give the standard $L^2$ energy inequality
\begin{equation}
 \norm{w(t)}_{L^2_x}
 \le C_T\left[
 \norm{w(0)}_{L^2_x}
 +\norm{r_h-r_j}_{L^1_tL^2_x}
 +\delta_{h,j}
 \right],
 \label{eq:dirac-L2-stability}
\end{equation}
where $\delta_{h,j}\to0$ is the coefficient distance multiplied by uniform $L^2$ bounds for the first derivatives of $\Psi_j$. The spatial derivative bounds follow from the uniform $H^2_x$ estimate; the normal derivative is bounded by solving \eqref{eq:dirac-hyperbolic} for $\partial_t\Psi_j$, using uniform invertibility of $A_j^0$ and boundedness of the residuals. Hence $\Psi_h$ is Cauchy in $C_tL^2_x$.

Interpolation on the closed three-manifold gives, uniformly in $t$,
\[
 \norm{w(t)}_{H^1_x}
 \le C\norm{w(t)}_{L^2_x}^{1/2}
        \norm{w(t)}_{H^2_x}^{1/2},
\]
so $\Psi_h$ is Cauchy in $C_tH^1_x$. Finally solve \eqref{eq:dirac-hyperbolic} for $\partial_t\Psi_h$. Strong convergence of the spatial $H^1$ terms and coefficients, together with $r_h$ Cauchy in $L^2_tL^2_x$, gives strong convergence of $\partial_t\Psi_h$ in $L^2_tL^2_x$. This is strong $H^1(I\times\Sigma)$ convergence. The dual system is identical. Passing to the limit in \eqref{eq:dirac-hyperbolic} proves the final assertion.
\end{proof}

\subsection{Primitive compactness certificates}

The strong packet used below is a conclusion of a concrete compactness mechanism once the quantities entering the nonlinear first variation are controlled in their own topologies. The low-regularity formulation is deliberately modular: direct spacetime compact screens may be used, or spatial screens may be combined with the time regularity of \Cref{prop:time-compactness}. For spinors one may instead use the Lorentzian evolution route of \Cref{prop:dirac-stability}; the regular Sobolev branch of \Cref{thm:regular-branch} avoids compact screens altogether.

\begin{definition}[Classical compactness certificate]
\label{def:compactness-certificate}
A finite-interface regulator sequence satisfies the classical compactness certificate on $K\Subset M$ if, after fixed local gauge and spin-frame identifications,
\begin{enumerate}[label=\textup{(C\arabic*)},leftmargin=3.2em]
\item the coframes are precompact in $L^\infty(K)$, uniformly nondegenerate, and bounded in $H^1(K)$; their first-derivative packet has a common compact screen in $L^2(K)$;
\item the connections are precompact in $L^4(K)$, the curvatures are bounded in $L^2(K)$, and the curvature packet has a common compact screen in $L^2(K)$;
\item the Higgs fields are bounded in $H^1(K)$, the covariant gradients are bounded in $L^2(K)$ and have a common compact screen, and there is an increasing $\Phi$ with $\Phi(t)/t\to\infty$ such that
\[
 \sup_h\int_K\Phi(\abs{H_h}^4)\,\dV_0<\infty;
\]
\item for the spinors, either
      \begin{enumerate}[label=\textup{(C4\alph*)},leftmargin=3.0em]
      \item $\Psi_h,\bPsi_h$ are precompact in $L^2(K)$ and bounded in $H^1(K)$, while their ordinary first-jet packets in the fixed local frames have common compact screens in $L^2(K)$; or
      \item on a compact subslab containing $K$, the reconstructed Dirac and dual-Dirac systems satisfy the hypotheses of \Cref{prop:dirac-stability};
      \end{enumerate}
\item the coefficient banks remain in a compact subset of the region $\kappa,g_j,\lambda_H>0$.
\end{enumerate}
A ``common compact screen'' means the hypotheses of \Cref{lem:screen}; it may instead be verified by the spatial-tail and time-regularity conditions of \Cref{prop:time-compactness}.
\end{definition}

\subsection{The primitive convergence class}

\begin{definition}[Classical strong packet]
\label{def:strong-packet}
A finite-interface regulator sequence has a classical strong packet with limit $z=(e,A,H,\Psi,\bPsi)$ if, on every $K\Subset M$ and in the fixed local bundle identifications,
\begin{align}
 &e_h\to e\text{ in }L^\infty(K)\cap H^1(K),
 \quad \sup_h(\norm{e_h}_\infty+\norm{e_h^{-1}}_\infty)<\infty,
 \label{eq:strong-geometry}\\
 &A_h\to A\text{ in }L^4(K),
 \qquad \Fcurv_{A_h}\to\Fcurv_A\text{ in }L^2(K),
 \label{eq:strong-gauge}\\
 &H_h\to H\text{ in }L^4(K),
 \qquad D_{A_h}H_h\to D_AH\text{ in }L^2(K),
 \label{eq:strong-Higgs}\\
 &\Psi_h\to\Psi,\quad\bPsi_h\to\bPsi
       \text{ in }H^1(K),
 \label{eq:strong-spinors}\\
 &\theta_h\to\theta,
 \quad \kappa_h,g_{j,h},\lambda_{H,h}\text{ remain positive and bounded away from zero}.
 \label{eq:strong-parameters}
\end{align}
The limiting coframe is in the same orientation and signature chart. The curvatures in \eqref{eq:strong-gauge} are identified with the reconstructed connections, not independent weak tensor limits.
\end{definition}

These conditions are sufficient, not claimed optimal. They can be discharged componentwise by common screens and lower-order identifications, together with a geometric noncollapse certificate. For example, \Cref{lem:screen} applies to curvature and covariant Higgs gradients, and \Cref{prop:orlicz} handles a critical Higgs amplitude. Such a screen must be established for the actual packet; it cannot be transferred from unrelated operator modes by notation alone.

The Higgs conditions imply strong ordinary $H^1$ convergence because
\[
 \dd H_h=D_{A_h}H_h-\rho_H(A_h)H_h
\]
and the second term converges strongly in $L^2$ by H\"older. For spinors, route \textup{(C4a)} is a direct compactness certificate, whereas route \textup{(C4b)} derives the same $H^1$ compactness from the Lorentzian Dirac evolution. No elliptic coercivity of a Lorentzian Dirac graph norm is used.

\begin{theorem}[Finite certificates produce the classical strong packet]
\label{thm:certificate-packet}
If a finite-interface regulator sequence satisfies \Cref{def:compactness-certificate} on every $K\Subset M$, then every cutoff sequence has a subsequence and a field tuple $z=(e,A,H,\Psi,\bPsi)$ for which the convergences of \Cref{def:strong-packet} hold locally.  The limiting curvature is $\Fcurv_A$ and the limiting Higgs derivative is $D_AH$; they are not independent tensor limits.
\end{theorem}
\begin{proof}
Fix $K$ and use a diagonal extraction over a finite bundle atlas. Zeroth-order bosonic precompactness gives, after extraction, $e_h\to e$ in $L^\infty$ and $A_h\to A$ in $L^4$. Under route \textup{(C4a)}, also extract $\Psi_h\to\Psi$ and $\bPsi_h\to\bPsi$ in $L^2$; route \textup{(C4b)} instead obtains the spinor limit from \Cref{prop:dirac-stability}.  The screened first-coframe derivatives converge strongly in $L^2$ by \Cref{lem:screen}, hence $e_h\to e$ in $H^1$; uniform inverse bounds keep $e$ in the same nondegenerate signature chart.  The screened curvatures converge strongly in $L^2$.  Since
\[
 \dd A_h=\Fcurv_{A_h}-A_h\wedge A_h,
\]
the right side converges distributionally to the strong curvature limit minus $A\wedge A$, which identifies that limit with $\Fcurv_A$.

Rellich compactness gives, after a further extraction, $H_h\to H$ in measure and weakly in $H^1$.  The Orlicz hypothesis and \Cref{prop:orlicz} then give strong $L^4$ convergence.  The screened covariant gradients converge strongly in $L^2$, and
\[
 \dd H_h=D_{A_h}H_h-\rho_H(A_h)H_h
\]
identifies their limit with $D_AH$ and yields strong ordinary $H^1$ convergence.  For the spinors, route \textup{(C4a)} gives strong $H^1$ convergence from the screened first jets. Under route \textup{(C4b)}, after the bosonic subsequence has been fixed, \Cref{prop:dirac-stability} gives the same strong spacetime $H^1$ convergence directly from the Lorentzian evolution. Finite-dimensional compactness gives a convergent coefficient bank, with the positive lower bounds retained.  A diagonal argument over an exhaustion $K_1\Subset K_2\Subset\cdots$ proves the local statement.
\end{proof}

\begin{lemma}[Local product and coefficient continuity]
\label{lem:products}
If $u_h\to u$ in $L^p$, $v_h\to v$ in $L^q$, and $1/p+1/q=1/r\le1$, then $u_hv_h\to uv$ in $L^r$. In particular,
\begin{equation}
 \norm{B(Y_h,Y_h)-B(Y,Y)}_1
 \le C\norm{Y_h-Y}_2(\norm{Y_h}_2+\norm Y_2)
 \label{eq:quadratic-product}
\end{equation}
for a bounded bilinear contraction $B$. On the coframe class \eqref{eq:strong-geometry}, every smooth algebraic function of $e_h,e_h^{-1}$ converges in $L^\infty\cap H^1$. Moreover,
\begin{equation}
 \omega(e_h)\to\omega(e)\quad\text{in }L^2,
 \qquad
 \dot\omega(e_h)[\dot e_h(k)]\to\dot\omega(e)[\dot e(k)]
       \quad\text{in }L^2,
 \label{eq:spin-connection-convergence}
\end{equation}
uniformly for $k$ in bounded $C^1$ test sets.
\end{lemma}
\begin{proof}
Subtract the products one factor at a time and apply H\"older. The coefficient assertion follows from the chain rule on a compact nondegenerate matrix chart. In coordinates,
\[
 \Gamma^\alpha_{\mu\nu}(g)
 =\tfrac12g^{\alpha\beta}
  (\partial_\mu g_{\beta\nu}+\partial_\nu g_{\beta\mu}-\partial_\beta g_{\mu\nu}),
\]
and
\[
 \omega^a{}_{b\mu}
 =e^a_\nu\Gamma^\nu_{\mu\lambda}(g)(e^{-1})_b^\lambda
   -(\partial_\mu e^a_\nu)(e^{-1})_b^\nu.
\]
Each expression is a bounded algebraic coefficient times one first coframe derivative. Its variation has the same structure, with an additional term linear in $\partial k$. The asserted convergence and uniform test bound follow from \eqref{eq:strong-geometry} and \eqref{eq:metric-lift}.
\end{proof}

\subsection{Continuity of the complete first variation}

Let $d_K(z_h,z)$ denote the sum of the difference norms in \eqref{eq:strong-geometry}--\eqref{eq:strong-spinors}, together with $\abs{\theta_h-\theta}$.  The same formula defines $d_K(z_h,z_j)$ for two reconstructed cutoffs on the common local bundle realization.  When necessary, take these norms on a slightly larger compact region containing the supports of the test lifts.

\begin{proposition}[First-variation continuity from primitive fields]
\label{prop:variation-continuity}
On a uniformly bounded classical strong packet,
\begin{equation}
 \norm{D\Sact_{\theta_h}(z_h)-D\Sact_\theta(z)}_{(\Vtest_K^{r_0})^*}
 \le C_K d_K(z_h,z).
 \label{eq:variation-bound}
\end{equation}
The same estimate holds separately for the gravitational and matter terms. The matter metric variation has the limit prescribed by \eqref{eq:stress-definition}; it is not an independent convergence assumption.
\end{proposition}

\begin{proof}
All arguments take place in one compact region and a finite bundle atlas. A partition of unity introduces only fixed smooth coefficients. We give the separate sector estimates, including the terms specific to metric variation.

\emph{Yang--Mills sector.} Write its local density as a bounded coframe coefficient contracted with two curvatures. A gauge variation gives
\[
 \dot\Fcurv_{A_h}=\dd a+[A_h,a],
\]
which converges in $L^2$ uniformly for bounded smooth $a$, by $A_h\to A$ in $L^4$. The pairing of $\dot\Fcurv_{A_h}$ and $\Fcurv_{A_h}$ therefore converges in $L^1$. A metric variation changes the coefficient and has the form
\begin{equation}
 B'_h[k](\Fcurv_{A_h},\Fcurv_{A_h}).
 \label{eq:metric-gauge-quadratic}
\end{equation}
The coefficient converges in $L^\infty$ and both curvature factors converge in $L^2$. Equation \eqref{eq:quadratic-product} proves convergence with the stated norm bound. If curvature were only weakly convergent, \eqref{eq:metric-gauge-quadratic} would contain two weak factors; a strongly convergent Hodge coefficient would not repair that loss.

\emph{Higgs kinetic sector.} With $K_h=D_{A_h}H_h$,
\[
 \dot K_h=D_{A_h}\eta_H+\rho_H(a)H_h.
\]
This converges strongly in $L^2$. The field variation is a pairing of $\dot K_h$ and $K_h$; the metric variation is a bounded coefficient applied to $K_h\otimes K_h$. Both follow from \Cref{lem:products}.

\emph{Higgs potential.} On a bounded $L^4$ set,
\begin{align}
 \norm{\abs{H_h}^4-\abs H^4}_1
 &\le C(\norm{H_h}_4+\norm H_4)^3\norm{H_h-H}_4,\notag\\
 \norm{\abs{H_h}^2H_h-\abs H^2H}_{4/3}
 &\le C(\norm{H_h}_4+\norm H_4)^2\norm{H_h-H}_4.
 \label{eq:quartic-continuity}
\end{align}
The lower-degree terms satisfy the analogous estimates. These inequalities pass both the volume variation of the potential and its Higgs variation, including the finite parameter limits.

\emph{Dirac--Yukawa sector.} Sobolev embedding on the four-dimensional compact chart gives strong $L^4$ convergence of both spinor variables from \eqref{eq:strong-spinors}. The derivative terms have the form
\[
 C(e_h)\bPsi_h\partial\Psi_h,
 \qquad C(e_h)(\partial\bPsi_h)\Psi_h.
\]
They converge in $L^1$, using the $L^2$ bounds for spinors and first derivatives. Spin-connection terms contain one $L^2$ coefficient $\omega(e_h)$ and the $L^2$ product $\bPsi_h\Psi_h$; they converge in $L^1$ by \eqref{eq:spin-connection-convergence}. Gauge and Yukawa terms contain, respectively, $A_h\bPsi_h\Psi_h$ and an affine function of $H_h$ multiplying $\bPsi_h\Psi_h$. H\"older with three $L^4$ factors gives convergence in $L^{4/3}$, hence in $L^1$ on the compact chart.

Varying $\Psi,\bPsi,A,H$ replaces one of these factors by its corresponding smooth test or derivative. Varying the metric differentiates $C(e_h)$ and $\omega(e_h)$. The additional spin term is $\dot\omega(e_h)[\dot e_h(k)]\bPsi_h\Psi_h$, again an $L^2$--$L^2$ product. Thus the full metric variation, including the spin connection, converges uniformly on bounded $C^1$ test sets. This proves the matter assertion without omitting a fermionic stress term.

\emph{Gravity.} Let $B(e)$ be the coframe bivector density of \Cref{app:palatini}. On compactly supported variations, integration by parts changes the torsion-free Palatini density into
\begin{equation}
 -\dd B(e)\wedge\omega(e)+B(e)\wedge\omega(e)\wedge\omega(e),
 \label{eq:palatini-first-order}
\end{equation}
with the chosen Einstein--Hilbert normalization. The volume term is algebraic in $e$. In the first variation of \eqref{eq:palatini-first-order}, each differentiated factor is in $L^2$, each algebraic coefficient is in $L^\infty$, and the variation introduces at most $k$ and $\partial k$. The coefficient and connection estimates of \Cref{lem:products} therefore apply to every term. They also define the gravitational first variation at the limiting $H^1\cap L^\infty$ metric. On smooth metrics it is the ordinary Einstein first variation, and its distributional extension is fixed by the integration-by-parts formula, as explained in \Cref{app:palatini}.

Subtracting each finite multilinear expression one factor at a time gives a bound linear in $d_K$. The other factors are uniformly bounded by the packet hypotheses. A finite sum of these estimates proves \eqref{eq:variation-bound}.
\end{proof}

\begin{remark}[Weak--strong gravity is a different route]
The geometric condition \eqref{eq:strong-geometry} is a strong first-jet route, not a requirement of strong second derivatives. \Cref{prop:weak-palatini} gives an alternative Palatini passage from strong coframe products and weak, identified curvature. That alternative does not weaken the strong curvature condition in the \emph{Yang--Mills stress}. The gravitational Palatini term is linear in curvature; the Yang--Mills metric variation is quadratic in curvature.
\end{remark}

\Needspace{26\baselineskip}
\subsection{Subsequential variational closure and cofinal uniqueness}

\begin{theorem}[General subsequential Einstein--Standard-Model variational closure]
\label{thm:main-limit}
Let $(\Iface_h)_{h\downarrow0}$ be any finite-interface regulator sequence in the sense of \Cref{def:regulator}. Assume it satisfies the classical compactness certificate of \Cref{def:compactness-certificate} on every $K\Subset M$, the finite reconstruction is first-variation consistent,
\begin{equation}
 c_h(K)\longrightarrow0,
 \label{eq:main-consistency}
\end{equation}
and the physical common-action residual is controlled by \Cref{prop:source-bound} with
\begin{equation}
 L_h\sqrt{R_h}+e_h\longrightarrow0.
 \label{eq:main-source-vanishing}
\end{equation}
Then every cutoff sequence $h_n\downarrow0$ has a subsequence $h_{n_j}$ and a field tuple $z=(e,A,H,\Psi,\bPsi)$ such that $z_{h_{n_j}}$ converges to $z$ in the classical strong-packet topology on compact subsets. Along this subsequence the reconstructed first variations converge to $D\Sact_\theta(z)$ and
\begin{equation}
 D\Sact_\theta(z)[v]=0\qquad(v\in\Vtest_K).
 \label{eq:main-stationary}
\end{equation}
Consequently the limiting fields solve the classical Yang--Mills, Higgs, and Dirac--Yukawa equations, and
\begin{equation}
 \Eins_{\mu\nu}(g)+\Lambda g_{\mu\nu}
 =\kappa T^{\mathrm{SM}}_{\mu\nu}(z)
 \quad\text{in }\Dtest'(M).
 \label{eq:Einstein-SM}
\end{equation}
More precisely, for the selected subsequence,
\begin{equation}
 \norm{D\Sact_\theta(z)}_{(\Vtest_K^{r_0})^*}
 \le L_{h_{n_j}}\sqrt{R_{h_{n_j}}}+e_{h_{n_j}}
      +c_{h_{n_j}}(K)+C_Kd_K(z_{h_{n_j}},z).
 \label{eq:main-residual-bound}
\end{equation}
No summable adjacent-cutoff transport, continuum stress tensor, Euler equation, or candidate strong limit is assumed. Compactness produces accumulation points; consistency and stationarity identify every such strong-packet accumulation point as a classical Einstein--Standard-Model solution.
\end{theorem}

\begin{proof}
Apply \Cref{thm:certificate-packet} to an arbitrary cutoff sequence and extract a strong-packet subsequence with coefficient limit $\theta$. \Cref{prop:variation-continuity} and first-variation consistency give
\[
 \abs{D\Sact_\theta(z)[v]-D\Sact_{h_{n_j}}(z_{h_{n_j}}^{\mathrm d})[\mathcal I_{h_{n_j}}v]}
 \le\bigl(C_Kd_K(z_{h_{n_j}},z)+c_{h_{n_j}}(K)\bigr)\norm v_{C^{r_0}}.
\]
The scaled source estimate of \Cref{prop:source-bound} bounds the operational term by $\bigl(L_{h_{n_j}}\sqrt{R_{h_{n_j}}}+e_{h_{n_j}}\bigr)\norm v_{C^{r_0}}$. This proves \eqref{eq:main-residual-bound}, whose right side tends to zero. Pure gauge, Higgs, spinor, and dual-spinor variations give the matter Euler equations. For a pure inverse-metric variation,
\[
 D_g\Sact_{\mathrm g,\theta}[k]
 =\frac{1}{2\kappa}
 \left\langle (\Eins(g)+\Lambda g)\dV_g,k\right\rangle,
\]
while \eqref{eq:stress-definition} gives the complete matter contribution. Hence stationarity yields \eqref{eq:Einstein-SM}. At the stated metric regularity the Einstein tensor is the variationally defined first-order distribution of \Cref{app:palatini}.
\end{proof}

\begin{proposition}[Summable cofinal transport]
\label{prop:cofinal-transport}
Let $h_n\downarrow0$ be a cofinal sequence satisfying the classical compactness certificate on every compact region, with adjacent cutoffs represented on the same local bundles.  Suppose
\begin{equation}
 d_K(z_{h_{n+1}},z_{h_n})\le\eta_{n,K},
 \qquad \sum_{n=1}^\infty\eta_{n,K}<\infty
 \label{eq:cofinal-transport}
\end{equation}
for every $K\Subset M$.  Then $(z_{h_n})$ is Cauchy in each norm entering $d_K$ and converges to a unique local strong-packet limit.  If two admissible cofinal routes admit direct comparisons whose $d_K$ error tends to zero after both cutoffs pass a common scale, their limits agree.  In particular, if \Cref{thm:certificate-packet} supplies subsequential strong-packet limits and \eqref{eq:cofinal-transport} holds, no subsequence choice remains.
\end{proposition}
\begin{proof}
For $m>n$, the triangle inequality gives
\[
 d_K(z_{h_m},z_{h_n})\le\sum_{j=n}^{m-1}\eta_{j,K}.
\]
The tail of the convergent series tends to zero, so every component is Cauchy in its complete target space.  The algebraic identifications of curvature and covariant derivative pass from any subsequential strong-packet limit supplied by \Cref{thm:certificate-packet}; uniqueness of the norm limit identifies the whole sequence with that field.  A direct comparison between two routes bounds the distance between their limiting fields by the vanishing comparison tail.
\end{proof}

\begin{corollary}[Cofinal uniqueness and route independence]
\label{cor:cofinal-unique}
Under the hypotheses of \Cref{thm:main-limit}, let a cofinal sequence additionally satisfy the summable transport estimate \eqref{eq:cofinal-transport}. Then the entire cofinal sequence converges in the strong-packet topology to a single classical Einstein--Standard-Model solution. If two admissible cofinal routes satisfy the direct-comparison hypothesis of \Cref{prop:cofinal-transport}, their limits agree.
\end{corollary}
\begin{proof}
\Cref{prop:cofinal-transport} makes the sequence Cauchy and removes subsequence ambiguity. Its unique limit is a strong-packet accumulation point, so \Cref{thm:main-limit} identifies it as a classical Einstein--Standard-Model solution. The route-independence assertion is the final statement of \Cref{prop:cofinal-transport}.
\end{proof}

\begin{theorem}[Regularity criterion for variational closure]
\label{thm:regular-branch}
This is a sufficient regularity criterion, not the constructive finite-dynamics theorem of \Cref{thm:generated-dynamics} or the quantitative finite-action theorem \Cref{thm:native-closure}.
Let $(\Iface_h)$ be a finite-interface regulator sequence and let $Q=I\times\Sigma\Subset M$ be a compact subslab. Suppose that for some $\sigma>0$, in fixed local gauges and frames,
\begin{equation}
 \sup_h\left(
 \norm{e_h}_{H^{3+\sigma}(Q)}
 +\norm{A_h}_{H^{3+\sigma}(Q)}
 +\norm{H_h}_{H^{3+\sigma}(Q)}
 \right)<\infty,
 \qquad
 \sup_h\norm{e_h^{-1}}_{L^\infty(Q)}<\infty,
 \label{eq:regular-bosonic-bound}
\end{equation}
and the coefficient banks remain in a compact physical set. Suppose the reconstructed spinors and dual spinors obey the Lorentzian systems of \Cref{prop:dirac-stability}, with uniformly bounded $L^\infty_tH^2_x$ norms, precompact initial data in $L^2(\Sigma)$, and residuals tending to zero in $L^2_tL^2_x$. Finally assume first-variation consistency and scaled stationarity, \eqref{eq:main-consistency} and \eqref{eq:main-source-vanishing}.

Then every cutoff sequence has a subsequence converging in the full classical strong packet on $Q$, and every such limit solves the classical Einstein--Standard-Model equations there. No compact-screen hypothesis and no a-priori spinor first-jet compactness is required on this regular branch.
\end{theorem}
\begin{proof}
The bounds \eqref{eq:regular-bosonic-bound} are stronger than those of \Cref{prop:sobolev-bosonic}; after extraction they give strong bosonic convergence and, by Sobolev embedding in four dimensions, convergence of the coefficients of the reconstructed Dirac systems in the $L^\infty$ topology required by \Cref{prop:dirac-stability}. Precompactness of the initial spinor data permits a further subsequence on which those data are Cauchy; the residuals already converge strongly to zero. \Cref{prop:dirac-stability} therefore gives strong spinor and dual-spinor $H^1(Q)$ convergence. Thus the full strong packet holds without invoking \Cref{def:compactness-certificate}. The proof of \Cref{thm:main-limit} from first-variation continuity, consistency, and scaled stationarity then applies verbatim.
\end{proof}


\subsection{A reduced low-regularity variational route}
\label{subsec:reduced-route}

The preceding strong-packet theorem is convenient when one wants norm convergence of every first jet, but the minimally coupled action needs less in two sectors.  The Higgs quartic endpoint is forced by a compact covariant gradient once the connections are compact in $L^4$, and the bilinear first-order fermionic variation is continuous under weak $H^1$ spinor convergence.  These observations remove an independent Higgs Orlicz condition and all spinor first-jet compactness from the reduced closure theorem below.

\begin{proposition}[Covariant-gradient compactness closes the Higgs endpoint]
\label{prop:covariant-higgs-endpoint}
Let $K\Subset M$ be a bounded smooth chart.  Suppose
\[
 A_h\to A\quad\text{in }L^4(K),\qquad
 \sup_h\norm{H_h}_{L^2(K)}<\infty,
 \qquad D_{A_h}H_h\to Y\quad\text{in }L^2(K).
\]
Then $(H_h)$ is precompact in $H^1(K)$.  Every subsequential limit $H$ satisfies $Y=D_AH$, and the convergence along that subsequence is strong in $H^1(K)$ and hence in $L^4(K)$.  If $H_h\to H$ in $L^2(K)$ is already known, the whole sequence converges strongly in these spaces.  In particular the quartic potential and the complete Higgs stress converge in $L^1$ without an independent Orlicz or $L^{4+\delta}$ hypothesis.
\end{proposition}
\begin{proof}
The closure of $\{A_h\}\cup\{A\}$ is norm compact in $L^4(K)$.  Multiplication by a fixed $L^4$ coefficient is compact $H^1\to L^2$: approximate the coefficient in $L^4$ by bounded functions and compose multiplication by the bounded approximation with the compact embedding $H^1\Subset L^2$.  Uniformity on the compact coefficient set therefore gives, for every $\epsilon>0$,
\begin{equation}
 \norm{BH}_2\le \epsilon\norm H_{H^1}+C_\epsilon\norm H_2
 \qquad(B\in\overline{\{A_h\}\cup\{A\}}).
 \label{eq:relative-multiplier}
\end{equation}
Since $\dd H_h=D_{A_h}H_h-\rho_H(A_h)H_h$, the ordinary $H^1$ norm and \eqref{eq:relative-multiplier} give a uniform $H^1$ bound after the first term on the right is absorbed.  Extract $H_h\rightharpoonup H$ in $H^1$ and $H_h\to H$ in $L^2$.

The product identity
\[
 A_hH_h-AH=(A_h-A)H_h+A(H_h-H)
\]
converges to zero in $L^2$: the first term uses boundedness of $H_h$ in $L^4$, and the second uses compact multiplication by the fixed $A\in L^4$.  Passing to distributions in the formula for $\dd H_h$ identifies $Y=D_AH$.  Applying the same graph estimate to $H_h-H$ gives strong $H^1$ convergence.  The continuous embedding $H^1(K)\hookrightarrow L^4(K)$ in four dimensions gives strong $L^4$ convergence, and the quadratic and quartic stress terms then converge by the product estimates of \Cref{lem:products}.
\end{proof}

\begin{proposition}[Weak-$H^1$ continuity of the complete fermionic first variation]
\label{prop:weak-fermion}
On $K\Subset M$, suppose
\begin{align*}
 e_h&\to e\quad\text{in }L^\infty\cap H^1,
 &\sup_h\norm{e_h^{-1}}_\infty&<\infty,\\
 A_h&\to A\quad\text{in }L^4,
 &H_h&\to H\quad\text{in }L^4,\\
 \Psi_h&\rightharpoonup\Psi\quad\text{in }H^1,
 &\bPsi_h&\rightharpoonup\bPsi\quad\text{in }H^1,
\end{align*}
and let the finite Yukawa/Majorana coefficients converge.  Then the complete first variation of the Dirac--Yukawa action \eqref{eq:dirac-density}, including its metric and spin-connection variation, converges against every physical test tuple.  For every integer $r\ge2$,
\begin{equation}
 \norm{D\Sact_{D,h}-D\Sact_D}_{(\Vtest_K^r)^*}\longrightarrow0.
 \label{eq:weak-fermion-variation}
\end{equation}
No strong convergence of the spinor first derivatives is required.  In particular the fermionic Hilbert stress has no independent defect in this class.
\end{proposition}
\begin{proof}
Weak $H^1$ convergence and Rellich compactness give strong $L^2$ convergence of both undifferentiated spinors and uniform $L^4$ bounds.  Their bilinears converge strongly in $L^1$ and remain bounded in $L^2$, hence converge weakly in $L^2$.  A derivative term pairs one weakly convergent $L^2$ derivative with one strongly convergent $L^2$ undifferentiated spinor.  The spin connection has the form $P(e)\partial e$, while its metric variation has the form $P_1(e)k\,\partial e+P_2(e)\partial k$; the assumed coframe convergence therefore makes both coefficients converge strongly in $L^2$, uniformly on bounded $C^1$ metric tests.  Gauge and Yukawa coefficients converge in $L^4$ and multiply weakly convergent $L^2$ spinor bilinears.  Thus every term in the complete first variation converges for a fixed smooth test.

The same expansions give a cutoff-independent bound by $C\norm v_{C^1}$ on bounded $H^1$ spinor sets.  The $C^r$ unit ball supported in a fixed compact set is totally bounded in $C^1$ for $r\ge2$.  A finite-net argument upgrades pointwise convergence to the dual norm \eqref{eq:weak-fermion-variation}.
\end{proof}

\begin{definition}[Reduced compactness certificate]
\label{def:reduced-certificate}
A regulator sequence satisfies the reduced compactness certificate on $K\Subset M$ when, after fixed local gauge and spin-frame identifications,
\begin{enumerate}[label=\textup{(R\arabic*)},leftmargin=3.2em]
\item the coframes are precompact in $L^\infty$, uniformly nondegenerate and bounded in $H^1$, and their first-derivative packet has a common compact screen in $L^2$;
\item the connections are precompact in $L^4$, and the curvature packet has a common compact screen in $L^2$;
\item the Higgs fields are bounded in $L^2$, and the covariant-gradient packet has a common compact screen in $L^2$;
\item $\Psi_h$ and $\bPsi_h$ are bounded in $H^1$; no spinor first-jet screen is required;
\item the coefficient banks remain in a compact subset of the physical parameter region.
\end{enumerate}
\end{definition}

\begin{theorem}[Reduced-certificate Einstein--Standard-Model closure]
\label{thm:reduced-closure}
Let $(\Iface_h)$ be a finite-interface regulator sequence satisfying \Cref{def:reduced-certificate} on every compact region, first-variation consistency $c_h(K)\to0$, and scaled common-action stationarity $L_h\sqrt{R_h}+e_h\to0$.  Then every cutoff sequence has a subsequence for which
\begin{align}
 e_h&\to e &&\text{in }L^\infty\cap H^1,\notag\\
 A_h&\to A &&\text{in }L^4,\qquad
 \Fcurv_{A_h}\to\Fcurv_A &&\text{in }L^2,\notag\\
 H_h&\to H &&\text{in }H^1\cap L^4,\qquad
 D_{A_h}H_h\to D_AH &&\text{in }L^2,\notag\\
 \Psi_h&\rightharpoonup\Psi &&\text{in }H^1,\qquad
 \bPsi_h\rightharpoonup\bPsi &&\text{in }H^1.
 \label{eq:reduced-limits}
\end{align}
Along this subsequence the complete reconstructed first variation converges in $(\Vtest_K^{r_0})^*$ and the limit satisfies all classical Einstein--Standard-Model Euler equations.  Thus neither an independent Higgs quartic compactness condition nor strong spinor first-jet compactness is part of this closure criterion.
\end{theorem}
\begin{proof}
The screen lemma gives strong convergence of the coframe first derivatives, curvature and covariant Higgs-gradient packets.  Zeroth-order precompactness identifies $e$ and $A$, and $\dd A_h=\Fcurv_{A_h}-A_h\wedge A_h$ identifies the curvature limit with $\Fcurv_A$.  \Cref{prop:covariant-higgs-endpoint} gives the Higgs limits in \eqref{eq:reduced-limits}.  Boundedness of the spinors in $H^1$ gives weakly convergent subsequences and, by Rellich, strong $L^2$ convergence.  The bosonic sectorwise first variations converge by the strong product estimates already used in \Cref{prop:variation-continuity}; the complete fermionic variation converges by \Cref{prop:weak-fermion}.  First-variation consistency and \Cref{prop:source-bound} therefore imply $D\Sact_\theta(z)=0$ exactly as in the proof of \Cref{thm:main-limit}.
\end{proof}

\begin{corollary}[Renewal Geometry realization]
\label{cor:renewal-realization}
On the source-native Standard-Model branch of Refs.~\cite{PelissierSpacetime2026,PelissierSpectral2026,PelissierDuality2026}, the structural hypotheses of \Cref{thm:main-limit} are realized by \Cref{prop:renewal-interface}. Hence every cutoff sequence on that branch satisfying the compactness, first-variation consistency, and scaled stationarity hypotheses has an Einstein--Standard-Model subsequential limit. If the additional transport hypotheses of \Cref{cor:cofinal-unique} hold, the full cofinal limit is unique and route independent.  On the explicit local common-action realization of \Cref{sec:native-branch}, \Cref{thm:native-closure} gives a stronger quantitative route: finite physical-row, initial/gauge and leakage budgets imply the required regularity and stress convergence, and \Cref{cor:native-rate} derives cofinal summability for the displayed refinement family.
\end{corollary}
\begin{proof}
The companion constructions provide the finite interface by \Cref{prop:renewal-interface}; the remaining hypotheses are analytic. Apply \Cref{thm:main-limit} for subsequential closure and \Cref{cor:cofinal-unique} when the optional transport condition is available.
\end{proof}

\begin{corollary}[Exact finite critical branch]
\label{cor:exact-critical}
Under the compactness and consistency hypotheses of \Cref{thm:main-limit}, suppose the finite configurations are exact critical points in all lifted physical directions,
\[
 D\Sact_h(z_h^{\mathrm d})[\mathcal I_hv]=0
 \qquad(v\in\Vtest_K).
\]
Then every strong-packet accumulation point satisfies \eqref{eq:Einstein-SM} and all classical matter Euler equations without a separate source-residual hypothesis.
\end{corollary}
\begin{proof}
Exact criticality gives $\varepsilon_h(K)=0$ in \eqref{eq:stationarity}.  The proof of \Cref{thm:main-limit} then uses only $c_h(K)\to0$ and $d_K(z_h,z)\to0$.
\end{proof}

\begin{corollary}[Distributional and negative-Sobolev stress convergence]
\label{cor:stress-topology}
Assume the sectorwise consistency in \eqref{eq:C1-consistency}. Define the operational stress density by
\[
 \langle\mathsf T_h,k\rangle
 =-2D\Sact_{\mathrm{SM},h}[\mathcal I_h(k,0,0,0,0)].
\]
Along every strong-packet subsequence furnished by \Cref{thm:main-limit}, $\mathsf T_h\to T^{\mathrm{SM}}\dV_g$ in distributions. On each fixed compact test region it also converges in $H^{-m}$ for every integer $m>r_0+2$, after using the fixed reference density and bundle frames.
\end{corollary}
\begin{proof}
The separate matter estimate in \Cref{prop:variation-continuity} gives convergence in the dual $C^{r_0}$ test norm. In four dimensions the continuous embedding $H_0^m\hookrightarrow C^{r_0}$ holds for $m>r_0+2$. Restricting the functional estimate to the unit ball of $H_0^m$ proves the stated norm convergence. This is a conservative negative-Sobolev conclusion; no pointwise or smooth stress convergence is inferred.
\end{proof}

\begin{corollary}[A realized random branch]
\label{cor:random-realization}
On a common probability space suppose the reconstructed fields satisfy the compactness and consistency hypotheses of \Cref{thm:main-limit} and the cofinal-transport hypotheses of \Cref{cor:cofinal-unique} almost surely, with limiting field $z(\omega)$.  Let the source estimate of \Cref{prop:source-bound} hold pathwise with deterministic $L_h$ and $e_h\to0$, and assume
\[
 \sum_h L_h^2\,\mathbb E\norm{a_h}_{G_h^{-1}}^2<\infty.
\]
Then the limiting fields satisfy \eqref{eq:Einstein-SM} and the matter Euler equations almost surely. It is enough to formulate the hypotheses on a countable dense test core with the displayed uniform bounds.
\end{corollary}
\begin{proof}
Tonelli's theorem implies $\sum_hL_h^2\norm{a_h}_{G_h^{-1}}^2<\infty$ almost surely. Hence the source bound makes $\varepsilon_h\to0$ on the realized sequence. Apply \Cref{thm:main-limit,cor:cofinal-unique} on the intersection of the full-measure events for a countable determining core and extend by continuity. No independence between cutoffs is required.
\end{proof}

\section{Defects, constraint propagation, and controlled realizations}
\label{sec:defects}

\subsection{What survives at critical Higgs regularity}

We first weaken the matter compactness rather than silently deleting the resulting source. Work on a compact chart with $g_h\to g$ in $C^1$, uniform inverse bounds, $A_h\to A$ in $L^4$, and
\begin{equation}
 H_h\rightharpoonup H\text{ in }H^1,
 \qquad H_h\to H\text{ almost everywhere and in }L^q\ (q<4).
 \label{eq:critical-Higgs}
\end{equation}
The latter convergence follows along a subsequence from Rellich compactness. These are positive comparison norms even though the action is Lorentzian.

\begin{lemma}[Cubic Euler terms and quartic concentration]
\label{lem:critical-cubic}
Under \eqref{eq:critical-Higgs},
\begin{equation}
 \abs{H_h}^2H_h\rightharpoonup\abs H^2H\text{ in }L^{4/3},
 \qquad D_{A_h}H_h\rightharpoonup D_AH\text{ in }L^2.
 \label{eq:critical-weak}
\end{equation}
After extraction on a fixed compact domain, there is a nonnegative Radon measure $\nu_H$ with
\begin{equation}
 \abs{H_h}^4\dV_{g_h}\rightharpoonup^*
 \abs H^4\dV_g+\nu_H.
 \label{eq:quartic-defect}
\end{equation}
On that compact domain, including its boundary, $\nu_H=0$ is equivalent to strong $L^4$ convergence along the extracted sequence. The local version applies on smaller compact regions, without excluding escape through the edge of a larger chart.
\end{lemma}
\begin{proof}
The cubic sequence is bounded in the reflexive space $L^{4/3}$ and converges almost everywhere; boundedness and convergence in measure identify every weak subsequential limit with the displayed cubic. For the covariant derivative, write
\[
 A_hH_h-AH=(A_h-A)H_h+A(H_h-H).
\]
The first term tends to zero in $L^2$. For any $L^2$ test $v$, the second term pairs with $A^*v\in L^{4/3}$ and tends to zero by weak $L^4$ convergence of $H_h$. The ordinary derivative converges weakly in $L^2$.

The fourth-power densities have bounded mass. Weak-star compactness and Fatou's lemma against nonnegative continuous tests give \eqref{eq:quartic-defect}. On a compact domain, testing with one gives convergence of fourth-power integrals if $\nu_H=0$. Weak $L^4$ convergence and uniform convexity then imply strong convergence. The converse follows from the product estimate \eqref{eq:quartic-continuity} and uniform convergence of the volume density.
\end{proof}

Let $\mathfrak K_H$ be the tensor-valued measure defect in the complete kinetic stress:
\begin{align}
 &\left(2\Rea\ip{D_\mu H_h}{D_\nu H_h}
      -g_{h,\mu\nu}\abs{D_{A_h}H_h}_{g_h}^2\right)\dV_{g_h}
 \rightharpoonup^*\nonumber\\[-0.3em]
 &\hspace{2em}
 \left(2\Rea\ip{D_\mu H}{D_\nu H}
      -g_{\mu\nu}\abs{D_AH}_g^2\right)\dV_g+\mathfrak K_{H,\mu\nu}.
 \label{eq:kinetic-defect}
\end{align}
Such a subsequence exists by the positive $L^2$ bound on the covariant derivatives. No positivity of the Lorentzian tensor $\mathfrak K_H$ is claimed.

\begin{theorem}[Critical Higgs stress-defect identity]
\label{thm:higgs-defect}
For the minimal action \eqref{eq:SM-action}, under the preceding hypotheses and convergent parameters,
\begin{equation}
 T^H_h\dV_{g_h}\rightharpoonup^*
 T^H\dV_g+\mathfrak S_H,
 \qquad
 \mathfrak S_{H,\mu\nu}
 =\mathfrak K_{H,\mu\nu}-\lambda_Hg_{\mu\nu}\nu_H.
 \label{eq:H-defect}
\end{equation}
If the remaining stress contributions and the gravitational first variation have their intended limits, and the operational metric residual and consistency error vanish, then
\begin{equation}
 \bigl(\Eins(g)+\Lambda g\bigr)\dV_g
 =\kappa\bigl(T^{\mathrm{SM}}\dV_g+\mathfrak S_H\bigr).
 \label{eq:defect-Einstein}
\end{equation}
If the total finite stress is conserved up to a vanishing distributional residual, then its limit is conserved. If the regular limiting fields satisfy the complete matter equations, $\nabla^\mu\mathfrak S_{H,\mu\nu}=0$. Conservation does not imply $\mathfrak S_H=0$.
\end{theorem}
\begin{proof}
The kinetic term is \eqref{eq:kinetic-defect}. In the potential, the quadratic density $\abs{H_h}^2$ converges in $L^1$, and the constant term converges with the volume density. The only additional potential contribution is the quartic measure \eqref{eq:quartic-defect}, with sign fixed by \eqref{eq:H-stress}. This proves \eqref{eq:H-defect}. Test the metric equation with a compactly supported smooth $k$ and pass each term to obtain \eqref{eq:defect-Einstein}.

For conservation, pair the finite stress with the symmetrized covariant derivative of a compactly supported vector field. The latter converges uniformly because $g_h\to g$ in $C^1$. Bounded stress measures and their weak-star convergence pass the identity. Subtract the regular limiting Noether identity when it is defined and holds. A nonzero constant symmetric tensor on a flat chart is already a counterexample to the implication from divergence zero to tensor zero.
\end{proof}

\begin{proposition}[Yang--Mills quadratic stress defect]
\label{prop:YM-defect}
Work on a compact chart with $g_h\to g$ in $C^1$, uniform inverse bounds, $A_h\to A$ in $L^4$, and
\[
 \Fcurv_{A_h}\rightharpoonup\Fcurv_A\quad\text{in }L^2.
\]
After extraction there is a symmetric tensor-valued finite Radon measure $\mathfrak S_{\mathrm{YM}}$ such that
\begin{equation}
 T^{\mathrm{YM}}_h\dV_{g_h}
 \rightharpoonup^*
 T^{\mathrm{YM}}\dV_g+\mathfrak S_{\mathrm{YM}}.
 \label{eq:YM-defect}
\end{equation}
Strong $L^2$ convergence of $\Fcurv_{A_h}$ is sufficient for $\mathfrak S_{\mathrm{YM}}=0$.
\end{proposition}
\begin{proof}
The positive reference $L^2$ bound on the curvature gives a uniform $L^1$ bound for every component of the quadratic stress density because the metric coefficients and volume densities converge uniformly.  Weak-star compactness of finite measures therefore gives a subsequential tensor-valued limit.  Subtracting the regular quadratic tensor built from $\Fcurv_A$ defines $\mathfrak S_{\mathrm{YM}}$.  If the curvatures converge strongly in $L^2$, the quadratic product estimate \eqref{eq:quadratic-product} makes the difference of the stress densities tend to zero in $L^1$.
\end{proof}

\begin{theorem}[Joint matter-defect Einstein limit]
\label{thm:joint-defect}
Assume the critical Higgs hypotheses of \Cref{thm:higgs-defect}, the Yang--Mills hypotheses of \Cref{prop:YM-defect}, convergence of the remaining fermionic stress contribution, and convergence of the gravitational first variation.  If the operational metric residual and its reconstruction-consistency error vanish, then
\begin{equation}
 \bigl(\Eins(g)+\Lambda g\bigr)\dV_g
 =\kappa\Bigl(T^{\mathrm{SM}}\dV_g
       +\mathfrak S_{\mathrm{YM}}+\mathfrak S_H\Bigr).
 \label{eq:joint-defect-Einstein}
\end{equation}
If the total finite stress is conserved up to a vanishing distributional residual and the regular limiting matter fields satisfy their Euler equations, the total defect $\mathfrak S_{\mathrm{YM}}+\mathfrak S_H$ is distributionally conserved.  In particular, the classical Einstein--Standard-Model equation is the zero-defect branch of this weak-limit identity.
\end{theorem}
\begin{proof}
Test the finite metric equation against a compactly supported smooth symmetric tensor and pass the gravitational, regular fermionic, Yang--Mills, and Higgs terms separately using \Cref{prop:YM-defect,thm:higgs-defect}.  This gives \eqref{eq:joint-defect-Einstein}.  Passing the finite conservation identity and subtracting the regular limiting Noether identity gives conservation of the sum of the defects.  Neither conservation nor the Ward identity forces that sum to vanish.
\end{proof}

For a nonminimal $\xi\abs H^2\Scal(g)$ term, the weak-limit source must additionally retain the assembled improvement defect of \Cref{app:improvement}.  The strong curvature and quartic/gradient compactness conditions in \Cref{thm:main-limit} are therefore best interpreted as explicit zero-defect certificates, not as consequences of conservation.

\begin{example}[A critical bubble]
\label{ex:bubble}
On a four-dimensional chart choose $\varphi\in C_c^\infty(\R^4)$, a point $x_0$, and a unit Higgs vector $v$. Set
\[
 H_\epsilon(x)=\epsilon^{-1}\varphi((x-x_0)/\epsilon)v.
\]
Then $\norm{H_\epsilon}_2=O(\epsilon)$, while $\norm{\dd H_\epsilon}_2$ and $\norm{H_\epsilon}_4$ are independent of $\epsilon$. The fields converge weakly in $H^1$ to zero, but
\[
 \abs{H_\epsilon}^4\dd x\rightharpoonup^*
 \left(\int_{\R^4}\abs\varphi^4\dd x\right)\delta_{x_0}.
\]
These assertions follow directly by $x=x_0+\epsilon y$. This example disproves a compactness inference from bounded critical energy; it does not assert that this bubble is a coupled stationary solution.
\end{example}

\begin{example}[Vanishing matter residual with conserved kinetic defect]
\label{ex:null-wave}
On flat $(0,T)\times\mathbb T^3$, take $v_H=0$, $A=0$, $\Psi=0$, a fixed unit Higgs vector $v$, and
\[
 H_n(t,x)=n^{-1}\sin(n(t+x^1))v,
 \qquad \ell=\dd t+\dd x^1.
\]
Since $\ell$ is null, $\Box H_n=0$. The Higgs Euler residual is the cubic potential force and is $O(n^{-3})$ uniformly. The gauge current vanishes because the amplitude is real and has a fixed internal direction. Although $H_n\to0$ strongly in $L^4$,
\[
 T^H_n
 =2\cos^2(n(t+x^1))\ell\otimes\ell
   -\lambda_Hn^{-4}\sin^4(n(t+x^1))g
 \rightharpoonup\ell\otimes\ell.
\]
Thus $\nu_H=0$ but the kinetic defect is nonzero and conserved. The example concerns the matter system on a fixed background: its nonzero metric residual is exactly what prevents it from being a vacuum Einstein limit. It isolates the error in trying to derive stress convergence from matter stationarity or Bianchi conservation alone.
\end{example}

\begin{proposition}[A sufficient Euclidean defect-removal mechanism]
\label{prop:monotonicity}
Let $V\hookrightarrow L^4$ be a real Hilbert space, and let the complete gauge-fixed Euclidean Higgs operators $\mathcal A_h:V\to V'$ satisfy
\[
 \ip{\mathcal A_h(u)-\mathcal A_h(v)}{u-v}
 \ge c_1\norm{u-v}_V^2+c_2\norm{u-v}_4^4,
 \qquad c_1,c_2>0.
\]
If $H_h\rightharpoonup H$ in $V$, $\mathcal A_h(H_h)=J_h\to J$ in $V'$, and $\mathcal A_h(H)\to J$ in $V'$, then $H_h\to H$ in $V\cap L^4$. If $V$ controls the covariant $H^1$ norm, both Higgs defects vanish.
\end{proposition}
\begin{proof}
Use $(u,v)=(H_h,H)$. The right pairing is
$\ip{J_h-\mathcal A_h(H)}{H_h-H}$, which tends to zero by strong dual convergence and weak boundedness. Each nonnegative term on the left tends to zero.
\end{proof}

This is an alternative sufficient branch, not an assertion about the Lorentzian evolution operator. A positive quartic term alone does not establish the complete monotonicity estimate, particularly with a symmetry-breaking mass term or unfixed gauge directions.

\subsection{A derivative-counted hyperbolic stability criterion}

We now describe a way of obtaining geometric convergence from a reduced evolution. This result is intentionally local in physical time. It assumes a common hyperbolic slab and a strongly controlled \emph{complete} source; it does not manufacture that source from the algebraic duality. The energy argument belongs to the standard symmetric-hyperbolic framework \cite{Kato1975,FriedrichRendall2000}.

Let $s\ge5$ be an integer and let smooth metrics $g_h$ on $[0,T]\times\Sigma$ solve, in one background-harmonic coordinate/frame realization,
\begin{equation}
 g_h^{\alpha\beta}\partial_\alpha\partial_\beta g_h
 =\mathcal N(g_h,\partial g_h)+\mathcal S_h.
 \label{eq:reduced-wave}
\end{equation}
Fixed background-connection terms are included in $\mathcal N$. Its dependence on first derivatives is the usual smooth quadratic one on a compact nondegenerate metric chart. Here $\mathcal S_h$ is a tensor source, not the action $\Sact_h$.

\begin{theorem}[Common-slab metric stability]
\label{thm:hyperbolic}
Assume uniform hyperbolicity with a common time function, uniform inverse-metric bounds, and
\begin{equation}
 \sup_h\left(\norm{g_h}_{L^\infty_tH^s_x}
       +\norm{\partial_tg_h}_{L^\infty_tH^{s-1}_x}\right)<\infty.
 \label{eq:hyperbolic-bound}
\end{equation}
Suppose the initial data are Cauchy in $H^{s-1}\times H^{s-2}$ and the sources are Cauchy in $L^1_tH^{s-2}_x$. Then
\begin{align}
 g_h&\longrightarrow g\quad\text{in }C_tH^{s-1}_x,\notag\\
 \partial_tg_h&\longrightarrow\partial_tg\quad\text{in }C_tH^{s-2}_x,\notag\\
 \Riem(g_h)&\longrightarrow\Riem(g)
       \quad\text{in }L^1_tH^{s-3}_x.
 \label{eq:hyperbolic-limit}
\end{align}
In particular, compatible coframes chosen on one compact frame chart converge in the geometric topology of \eqref{eq:strong-geometry}. A vanishing reduced-equation residual can be included in $\mathcal S_h$ with the same source norm.
\end{theorem}

\begin{proof}
Subtract the equations for $h$ and $j$, and put $w=g_h-g_j$. Use the principal wave operator with coefficient $g_h$ on the left. The forcing consists of $\mathcal S_h-\mathcal S_j$, the difference of the quadratic first-derivative terms, and
\[
 (g_j^{\alpha\beta}-g_h^{\alpha\beta})\partial_\alpha\partial_\beta g_j.
\]
At one derivative below the a priori bound, $H^{s-2}(\Sigma)$ is an algebra. Spatial second derivatives of $g_j$ and mixed time--space derivatives are uniformly bounded in $H^{s-2}$. Solving \eqref{eq:reduced-wave} for $\partial_t^2g_j$ gives an $L^1_tH^{s-2}$ bound with size controlled by $1+\norm{\mathcal S_j}_{L^1_tH^{s-2}}$. Consequently the forcing, apart from $\mathcal S_h-\mathcal S_j$, has $H^{s-2}$ norm bounded by
\[
 a_{h,j}(t)\bigl(\norm w_{H^{s-1}}+\norm{\partial_tw}_{H^{s-2}}\bigr),
 \qquad \sup_{h,j}\int_0^T a_{h,j}(t)\dd t<\infty.
\]
The differentiated wave energy $E_{h,j}$ is uniformly equivalent to the square of the displayed difference norm. Its square root obeys the integrated estimate
\begin{equation}
 E_{h,j}(t)^{1/2}
 \le C\left(E_{h,j}(0)^{1/2}
   +\int_0^t\norm{\mathcal S_h-\mathcal S_j}_{H^{s-2}}\dd u\right)
   \exp\!\left(C\int_0^t(1+a_{h,j}(u))\dd u\right).
 \label{eq:hyperbolic-energy}
\end{equation}
For completeness, commute spatial derivatives through the principal operator, use the differentiated normal derivative as energy multiplier, and integrate on the closed $\Sigma$. The uniformly spacelike slices and bounded lapse and shift give a positive wave energy; commutators are bounded by the coefficient norms in \eqref{eq:hyperbolic-bound} and the Sobolev product estimate above. Apply the scalar differential inequality to $(E_{h,j}+\delta)^{1/2}$ and let $\delta\downarrow0$ to obtain \eqref{eq:hyperbolic-energy}. This form uses the $L^1$ source norm directly; it does not replace it by an unjustified squared-source bound.

Cauchy initial and source data prove the first two limits. The reduced equation then gives convergence of $\partial_t^2g_h$ in $L^1_tH^{s-3}$. All other second derivatives converge in $C_tH^{s-3}$, and the inverse metrics and quadratic first-derivative terms converge by Sobolev product estimates. The coordinate formula for curvature gives the last limit. Smooth dependence of a locally fixed coframe on the metric gives the coframe assertion.
\end{proof}

The loss of one derivative in the difference estimate avoids treating a quasilinear coefficient difference times a second derivative as if it had top-order Lipschitz regularity. Likewise, $L^1$ control of the source yields the stated $L^1$-in-time curvature conclusion, not continuity in time of every curvature component. Stronger conclusions require stronger source and multiplier estimates.

\subsection{The Ward residual forces the harmonic constraint}

For a harmonic gauge covector $C_\mu$ define
\begin{equation}
 \widehat\Eins_{\mu\nu}
 =\Eins_{\mu\nu}-\nabla_{(\mu}C_{\nu)}
        +\tfrac12g_{\mu\nu}\nabla^\alpha C_\alpha.
 \label{eq:reduced-Einstein}
\end{equation}
Suppose a smooth reconstructed field satisfies
\[
 \widehat\Eins_h+\Lambda_hg_h=\kappa_h T_h+r_h,
 \qquad
 \mathcal W_{h,\nu}=\kappa_h\nabla_h^\mu T_{h,\mu\nu}
                         +\nabla_h^\mu r_{h,\mu\nu},
\]
with constant parameters on spacetime. The source $\mathcal W_h$ includes every retained sector and the reduced-equation residual.

\begin{proposition}[Subsidiary equation and constraint stability]
\label{prop:subsidiary}
With the curvature convention fixed by \eqref{eq:reduced-Einstein},
\begin{equation}
 \Box_{g_h}C_{h,\nu}+\Ric(g_h)_\nu{}^\mu C_{h,\mu}
 =-2\mathcal W_{h,\nu}.
 \label{eq:subsidiary}
\end{equation}
On a common closed spatial slab with uniform wave-energy coefficients and a uniform $L^1_tL^\infty_x$ bound for the curvature potential,
\begin{equation}
 \norm{(C_h,\nabla C_h)}_{L^\infty_tL^2_x}
 \le C_T\left(
 \norm{(C_h,\nabla C_h)(0)}_{L^2_x}
 +\norm{\mathcal W_h}_{L^1_tL^2_x}\right).
 \label{eq:constraint-bound}
\end{equation}
Vanishing initial constraint data and Ward residual therefore imply $C_h\to0$ in the subsidiary energy norm.
\end{proposition}
\begin{proof}
Take the covariant divergence of the reduced Einstein equation. Contracted Bianchi removes the Einstein tensor and the constant cosmological term. The identity
\[
 2\nabla^\mu\nabla_{(\mu}C_{\nu)}
   -\nabla_\nu\nabla^\alpha C_\alpha
 =\Box C_\nu+\Ric_\nu{}^\mu C_\mu
\]
gives \eqref{eq:subsidiary}. For the estimate, use the positive wave energy of a one-form relative to the common time foliation, including its $L^2$ term. Integration by parts has no spatial boundary flux. Coefficient derivatives and the curvature potential multiply the energy by an integrable function, and the forcing pairs with the time derivative. The square-root energy inequality and Gronwall give \eqref{eq:constraint-bound}.
\end{proof}

This proposition is not a consequence of the finite commutant identity. Exact classical gauge and relabeling invariance provide Noether identities, but their continuum use requires the whole source and the physical test realization. In an evolution-based construction, the subsidiary residual must be monitored along with the reduced equation. On the variational branch of \Cref{thm:main-limit}, by contrast, independent physical variations already impose the full Euler system. Conservation is asserted at the regularity at which the corresponding covariant distributional identity is defined, not by multiplying rough distributions formally.

A physical timelike boundary adds a boundary energy flux to \eqref{eq:constraint-bound}. Maximal dissipation for the main system does not automatically preserve the subsidiary constraint, and dissipation alone does not identify a boundary domain of the same closed action. These are separate extensions; the present closed-spatial-cylinder result avoids, rather than solves, that boundary problem.

\subsection{A controlled comparison branch}

\begin{proposition}[Smooth sampling is a consistency realization]
\label{prop:smooth-sampling}
Let $z_*$ be a smooth classical configuration on a compact slab with a fixed nondegenerate coframe chart. On the prescribed shape-regular reconstruction and local action of \Cref{app:reconstruction}, smooth sampling gives $z_h\to z_*$ in the strong packet and
\begin{equation}
 c_h(K)\le C_Kh.
 \label{eq:sampling-consistency}
\end{equation}
If $z_*$ solves the classical Euler equations, then on the same smooth family
\begin{equation}
 \varepsilon_h(K)\le C'_Kh.
 \label{eq:sampling-stationarity}
\end{equation}
Under the nondegeneracy hypotheses of \Cref{prop:critical-shadowing}, these approximate critical points are shadowed by exact finite critical points.
\end{proposition}
\begin{proof}
The interpolation, holonomy, and differentiated quadrature estimates are proved in \Cref{prop:mesh-consistency}. They imply the packet and \eqref{eq:sampling-consistency}. If $D\Sact_\theta(z_*)=0$, the triangle inequality and \Cref{prop:variation-continuity} give
\[
 \varepsilon_h(K)\le c_h(K)+C_Kd_K(z_h,z_*).
\]
The smooth reconstruction estimates of \Cref{prop:mesh-consistency} give $d_K(z_h,z_*)=O(h)$ on the prescribed smooth family, proving \eqref{eq:sampling-stationarity}.
\end{proof}

\begin{proposition}[Nondegenerate critical-point shadowing]
\label{prop:critical-shadowing}
Let $z_*$ be a smooth classical critical configuration and $x_h^0$ its sampled finite configuration in a gauge-fixed finite-dimensional physical slice $X_h$.  Represent the finite Euler covector on that slice by a map $F_h:X_h\to X_h$ using a fixed positive slice inner product.  Suppose, on a ball $B_r(x_h^0)$ and uniformly for small $h$,
\begin{enumerate}[label=\textup{(S\arabic*)},leftmargin=3.2em]
\item $\norm{F_h(x_h^0)}\le C_0h$;
\item $D F_h(x_h^0)$ is invertible and $\norm{D F_h(x_h^0)^{-1}}\le M$;
\item $\norm{D F_h(x)-D F_h(y)}\le L\norm{x-y}$;
\item the reconstruction is locally Lipschitz, $d_K(\mathcal R_hx,\mathcal R_hy)\le C_R\norm{x-y}$.
\end{enumerate}
Then, for sufficiently small $h$, there is a unique exact finite critical point $x_h^*\in B_r(x_h^0)$ in the Newton ball
\begin{equation}
 \norm{x_h^*-x_h^0}\le 2MC_0h,
 \label{eq:critical-shadowing}
\end{equation}
and its reconstruction satisfies $d_K(\mathcal R_hx_h^*,z_*)=O(h)$.  Hence a nondegenerate smooth Einstein--Standard-Model solution generates a nearby exact finite stationary branch whenever the finite linearized physical Euler operator has a cutoff-uniform inverse and the second variation is uniformly controlled.
\end{proposition}
\begin{proof}
Put $B_h=D F_h(x_h^0)^{-1}$ and define the frozen-Newton map
\[
 \mathcal T_h(x)=x-B_hF_h(x).
\]
Let $\rho_h=2MC_0h$.  Taylor's formula and (S3) give, for $x,y\in B_{\rho_h}(x_h^0)$,
\[
 \norm{\mathcal T_h(x)-\mathcal T_h(y)}
 \le ML\rho_h\norm{x-y}.
\]
For small $h$, $ML\rho_h<1/2$.  The same Taylor estimate gives
\[
 \norm{\mathcal T_h(x)-x_h^0}
 \le MC_0h+\tfrac12ML\rho_h^2\le\rho_h,
\]
again for small $h$.  Thus $\mathcal T_h$ is a contraction of the Newton ball into itself and has a unique fixed point there.  Since $B_h$ is invertible, the fixed-point equation is $F_h(x_h^*)=0$.  The reconstruction estimate follows from (S4), \eqref{eq:critical-shadowing}, and the $O(h)$ sampling convergence of \Cref{prop:smooth-sampling}.
\end{proof}

If the positive slice norm used to represent $F_h$ is uniformly equivalent to the physical dual norm induced by the test lifts, condition (S1) follows from \eqref{eq:sampling-stationarity}.  The substantive additional hypothesis is then the cutoff-uniform invertibility and local Lipschitz control of the linearized physical Euler map.

This establishes a nonempty classical comparison class once a smooth solution is specified, and it upgrades that class to exact finite stationarity on every locally nondegenerate branch satisfying \Cref{prop:critical-shadowing}.  It is not a proof that every renewal history selects this discretization or that every classical solution has a uniformly invertible physical linearization.  The following explicit reduction supplies a backreacting comparison field without appealing to a general coupled existence theorem.

\begin{proposition}[A nonvacuum Einstein--Higgs comparison family]
\label{prop:homogeneous}
Let $\Sigma=\mathbb T^3$, $\kappa,\lambda_H>0$, and take
\begin{equation}
 g=-\dd t^2+a(t)^2\delta_{ij}\dd x^i\dd x^j,
 \quad H=\frac{\phi(t)}{\sqrt2}h_0,
 \quad A=0,\quad\Psi=\bPsi=0,
 \label{eq:homogeneous-fields}
\end{equation}
where $h_0$ is a fixed unit Higgs vector and $\phi$ is real. Put
\[
 U(\phi)=\lambda_H(\phi^2/2-v_H^2)^2.
\]
For initial data $a_0>0$, $\phi_0$, $\pi_0$, and $\mathcal H_0$ satisfying
\begin{equation}
 3\mathcal H_0^2=\Lambda+\kappa\left(\tfrac12\pi_0^2+U(\phi_0)\right),
 \label{eq:Friedmann-initial}
\end{equation}
the system
\begin{equation}
 \dot a=a\mathcal H,\qquad
 \dot\phi=\pi,\qquad
 \dot\pi=-3\mathcal H\pi-U'(\phi),\qquad
 \dot{\mathcal H}=-\frac\kappa2\pi^2
 \label{eq:homogeneous-ODE}
\end{equation}
has a smooth local solution with $a>0$, and \eqref{eq:homogeneous-fields} solves the torsion-free Einstein--Standard-Model equations. Choosing $\pi_0\ne0$ gives a genuinely time-dependent matter stress and backreaction. Smooth sampling yields the finite residual conclusion of \Cref{prop:smooth-sampling}.
\end{proposition}
\begin{proof}
The right side of \eqref{eq:homogeneous-ODE} is smooth, so local ordinary differential equation existence and uniqueness apply. Since $a(t)=a_0\exp(\int_0^t\mathcal H)$, the scale factor stays positive on its existence interval. The Higgs kinetic term is $-\tfrac12g^{\mu\nu}\partial_\mu\phi\partial_\nu\phi$, giving energy density $\rho=\tfrac12\pi^2+U$ and pressure $p=\tfrac12\pi^2-U$. Its Euler equation is the third equation of \eqref{eq:homogeneous-ODE}.

The gauge current vanishes: $H$ and $\dot H$ have the same fixed internal direction and a real relative coefficient, whereas all infinitesimal unitary generators are skew-adjoint. All fermionic equations are satisfied by zero spinors. Differentiating the Friedmann residual gives
\begin{align*}
 \frac{\dd}{\dd t}\left(3\mathcal H^2-\Lambda
             -\kappa(\tfrac12\pi^2+U)\right)
 &=6\mathcal H\left(-\tfrac\kappa2\pi^2\right)
   -\kappa\left(\pi(-3\mathcal H\pi-U')+U'\pi\right)=0.
\end{align*}
Thus the constraint persists. Together with $\dot{\mathcal H}=-(\kappa/2)(\rho+p)$ it gives both independent Einstein equations for the flat homogeneous metric. If $\pi_0\ne0$, then $\dot{\mathcal H}(0)<0$, so the example is not merely a vacuum metric relabeled by an internal carrier. On a smaller compact existence slab all required smoothness and noncollapse bounds hold, and \Cref{prop:smooth-sampling} applies.
\end{proof}

\section{Discussion and outlook}
\label{sec:discussion}

The paper has two complementary analytic outputs.  The quantitative result, \Cref{thm:native-closure}, is a finite-action stability theorem for the coupled Einstein--Yang--Mills--Higgs--Dirac system: a small physical finite Euler row, compatible initial and harmonic-gauge data, and measured spectral leakage yield the all-time Sobolev estimate and quantitative convergence of the fields, curvature and complete Standard-Model stress.  The local action used there is written explicitly in \Cref{app:native-action}; the spectral decomposition enters only the estimates and never filters the physical record.  Renewal Geometry supplies one intrinsic realization of this common finite action through the companion constructions.

The regulator-independent result, \Cref{thm:main-limit}, gives a broader low-regularity variational closure principle.  It identifies classical Einstein--Standard-Model accumulation points whenever the quadratic bosonic packets satisfy explicit no-concentration conditions and the physical common-action residual vanishes.  The reduced certificate of \Cref{thm:reduced-closure} removes an independent Higgs quartic assumption and strong spinor first-jet compactness: compact covariant Higgs gradients force the critical Higgs endpoint, while bounded weak $H^1$ spinors suffice for the complete bilinear fermionic first variation.  The defect theorems describe the complementary weak-limit regime when those bosonic compactness mechanisms fail.

The distinction between the results is useful.  The general theorem is portable across finite realizations but necessarily exposes its compactness assumptions.  On the explicit finite-action branch the same nonlinear regularity, complete stress convergence and cofinal transport are outputs of finite physical budgets.  The same-source theorem \Cref{thm:native-selected-closure} now produces those budgets jointly on one actually observed record from signed action drift, occupation control, a verified physical reader and initial/gauge alignment.  The exact reader criterion makes leakage a finite incidence test between the native positive energy and the declared spacetime reader, while \Cref{cor:gauge-robust-reader} gives a gauge-covariant realization for the weak sector.  What remains source-specific is verification that the independently specified primitive coupled source satisfies these finite tests simultaneously across gravity, gauge and chiral matter.

The second improvement is variational.  The quantity $L_h\sqrt{R_h}+e_h$ records the cost of lifting a physical variation into the finite source geometry, so stationarity is tested in the same normalization as the limiting field equation.  Exact finite critical points are the zero-residual case.  Conversely, smooth sampling gives an $O(h)$ physical residual on the prescribed regular reconstruction, and \Cref{prop:critical-shadowing} shows that, when this residual controls a uniformly compatible physical-slice norm, a uniform inverse linearization and second-variation bound promote the sampled approximate critical point to a nearby exact one.  This separates the genuinely dynamical nondegeneracy question from the already proved discretization consistency.

The matter source is more delicate than the matter Euler equations.  Higgs stationarity probes a cubic force while the metric variation sees a quartic density and a quadratic gradient tensor.  Yang--Mills stationarity is linear in one curvature variation while the metric variation is quadratic in curvature.  The defect theorems therefore give a useful complementary formulation: an energy-level weak limit does not cease to exist when strong compactness fails; instead it solves Einstein's equation with the additional source $\mathfrak S_{\mathrm{YM}}+\mathfrak S_H$.  The ordinary classical Einstein--Standard-Model system is the defect-free branch selected by the compactness certificates.  Conservation constrains the effective source but cannot remove it by itself.

The fermionic sector illustrates this distinction sharply. A Lorentzian Dirac operator is hyperbolic, so an elliptic $H^1$ graph-norm estimate would be inappropriate. The low-regularity certificate retains a direct jet-screen option, but \Cref{prop:dirac-stability} supplies an independent evolution route: symmetric-hyperbolic $L^2$ stability plus a uniform $H^2_x$ bound and interpolation yield strong $H^1$ compactness, and the equation itself controls the normal derivative. This addresses spinor compactness without importing the Riemannian norm-resolvent theorem of the spectral companion and without pretending that Lorentzian Dirac coercivity is elliptic.

The hyperbolic results address a logically different route.  If the continuum configuration is generated by reduced evolution rather than obtained directly from variational compactness, one must control a common physical-time slab, the complete source, and the harmonic constraint.  The derivative-counted estimate records the actual loss of one derivative in comparing quasilinear solutions, while the subsidiary equation shows that the assembled Ward residual is precisely the forcing of the harmonic constraint.  These statements prevent reduced evolution, variational stationarity, and constraint propagation from being conflated.

The controlled class is not empty.  The homogeneous Einstein--Higgs family gives a simple backreacting comparison field, and \Cref{cor:generated-nonempty} supplies a fully finite approximation generated from constrained finite initial data.  More importantly for the native interpretation, \Cref{thm:native-selected-closure} acts on an actual accepted source family rather than on that comparison scheme: it keeps the selected record and its stochastic fluctuations intact and requires only the displayed source-level action, occupation, reader and legal-domain certificates.  What remains outside the theorem is the end-to-end verification of those certificates for an arbitrary independently specified primitive Renewal source.

The quantum problem is also separate.  A positive operational history law and a classical spinor action do not by themselves construct a chiral fermionic measure, reflection-positive gravity, or an interacting renormalized continuum source.  None of those quantum requirements is used as an unspoken hypothesis in the classical theorem.  The present endpoint is the quantitative and variational continuum closure of one finite common gravity--matter action on the controlled branches described above.

The logical architecture therefore has a quantitative finite-action route and a general low-regularity route:
{\small
\begin{equation}
\begin{aligned}
\text{accepted source}+\text{action drift/occupation control}
&\longrightarrow \text{one selected record with small full Euler row},\\
\text{native energy}+\text{verified physical reader}
&\longrightarrow \text{measured spacetime leakage},\\
\text{selected row/leakage}+\text{initial/gauge control}
&\longrightarrow \text{field/curvature/stress convergence},\\
\text{summable source and refinement budgets}
&\longrightarrow \text{unique cofinal route almost surely},\\[0.2em]
\text{finite common-action data}+\text{low-regularity control}
&\longrightarrow \text{Einstein--Standard-Model limits}.
\end{aligned}
\label{eq:final-chain}
\end{equation}
}
The first four lines describe the same-source native selection and quantitative finite-action route; the last line is the portable low-regularity variational closure theorem.  The companion constructions certify that the finite common-action structure is realized inside Renewal Geometry.  The same-source theorem keeps one observed record throughout the selection, reader and PDE handoff, while the finite approximation branch serves only as an independent nonemptiness check.  On the quantitative branch the compactness, stress and transport conclusions are outputs of finite error budgets rather than separate hypotheses.  The remaining limitation is primitive-source verification for the full coupled packet, not the existence of a closure theorem once the stated finite certificates hold.  Failure of nonlinear compactness outside the controlled branches remains visible through the effective stress defects rather than being normalized away.

\paragraph{Data and formalization.}
The results use no empirical dataset. The finite consistency estimates and continuum arguments are given in ordinary mathematical form; no machine-checked formalization of the present analytic theorems is asserted.

\begin{thebibliography}{99}
\small
\setlength{\itemsep}{2pt}

\bibitem{PelissierRenewalSpacetime2026}
A.~P\'elissier,
\newblock Renewal Geometry and the emergence of Lorentzian spacetime,
\newblock HAL preprint hal-05744772, 2026.
\newblock \url{https://hal.science/hal-05744772}

\bibitem{PelissierSpectral2026}
A.~P\'elissier,
\emph{From predictive dynamics to spectral geometry}, 2026.
% Companion paper II. Insert its final public bibliographic identifier before submission.

\bibitem{PelissierDuality2026}
A.~P\'elissier,
\emph{Spacetime--Gauge Duality in Renewal Geometry}, 2026.
% Companion paper III. Insert its final public bibliographic identifier before submission.

\bibitem{ChamseddineConnes1997}
A.~H.~Chamseddine and A.~Connes,
The spectral action principle,
\emph{Commun. Math. Phys.} \textbf{186} (1997), 731--750.
\href{https://doi.org/10.1007/s002200050126}{doi:10.1007/s002200050126}.

\bibitem{ChamseddineConnesMarcolli2007}
A.~H.~Chamseddine, A.~Connes, and M.~Marcolli,
Gravity and the standard model with neutrino mixing,
\emph{Adv. Theor. Math. Phys.} \textbf{11} (2007), 991--1089.
\href{https://doi.org/10.4310/ATMP.2007.v11.n6.a3}{doi:10.4310/ATMP.2007.v11.n6.a3}.

\bibitem{CheegerMullerSchrader1984}
J.~Cheeger, W.~M\"uller, and R.~Schrader,
On the curvature of piecewise flat spaces,
\emph{Commun. Math. Phys.} \textbf{92} (1984), 405--454.
\href{https://doi.org/10.1007/BF01210729}{doi:10.1007/BF01210729}.

\bibitem{ArnoldFalkWinther2006}
D.~N.~Arnold, R.~S.~Falk, and R.~Winther,
Finite element exterior calculus, homological techniques, and applications,
\emph{Acta Numer.} \textbf{15} (2006), 1--155.
\href{https://doi.org/10.1017/S0962492906210018}{doi:10.1017/S0962492906210018}.

\bibitem{ArnoldFalkWinther2010}
D.~N.~Arnold, R.~S.~Falk, and R.~Winther,
Finite element exterior calculus: from Hodge theory to numerical stability,
\emph{Bull. Amer. Math. Soc.} \textbf{47} (2010), 281--354.
\href{https://doi.org/10.1090/S0273-0979-10-01278-4}{doi:10.1090/S0273-0979-10-01278-4}.

\bibitem{ChristiansenHalvorsen2012}
S.~H.~Christiansen and T.~G.~Halvorsen,
A simplicial gauge theory,
\emph{J. Math. Phys.} \textbf{53} (2012), 033501.
\href{https://doi.org/10.1063/1.3692167}{doi:10.1063/1.3692167}.

\bibitem{Uhlenbeck1982}
K.~K.~Uhlenbeck,
Connections with $L^p$ bounds on curvature,
\emph{Commun. Math. Phys.} \textbf{83} (1982), 31--42.
\href{https://doi.org/10.1007/BF01947069}{doi:10.1007/BF01947069}.

\bibitem{Lions1985}
P.-L.~Lions,
The concentration-compactness principle in the calculus of variations.
The limit case, Part~1,
\emph{Rev. Mat. Iberoam.} \textbf{1} (1985), no.~1, 145--201.
\href{https://doi.org/10.4171/RMI/6}{doi:10.4171/RMI/6}.

\bibitem{Burnett1989}
G.~A.~Burnett,
The high-frequency limit in general relativity,
\emph{J. Math. Phys.} \textbf{30} (1989), 90--96.
\href{https://doi.org/10.1063/1.528594}{doi:10.1063/1.528594}.

\bibitem{HuneauLuk2024}
C.~Huneau and J.~Luk,
Burnett's conjecture in generalized wave coordinates,
\href{https://arxiv.org/abs/2403.03470}{arXiv:2403.03470 [math.AP]}, 2024.

\bibitem{MullerNowaczyk2017}
O.~M\"uller and N.~Nowaczyk,
A universal spinor bundle and the Einstein--Dirac--Maxwell equation as a variational theory,
\emph{Lett. Math. Phys.} \textbf{107} (2017), 933--961.
\href{https://doi.org/10.1007/s11005-016-0929-4}{doi:10.1007/s11005-016-0929-4}.

\bibitem{Kibble1961}
T.~W.~B.~Kibble,
Lorentz invariance and the gravitational field,
\emph{J. Math. Phys.} \textbf{2} (1961), 212--221.
\href{https://doi.org/10.1063/1.1703702}{doi:10.1063/1.1703702}.

\bibitem{HehlEtAl1976}
F.~W.~Hehl, P.~von der Heyde, G.~D.~Kerlick, and J.~M.~Nester,
General relativity with spin and torsion: Foundations and prospects,
\emph{Rev. Mod. Phys.} \textbf{48} (1976), 393--416.
\href{https://doi.org/10.1103/RevModPhys.48.393}{doi:10.1103/RevModPhys.48.393}.

\bibitem{Simon1986}
J.~Simon,
Compact sets in the space $L^p(0,T;B)$,
\emph{Ann. Mat. Pura Appl.} \textbf{146} (1986), 65--96.
\href{https://doi.org/10.1007/BF01762360}{doi:10.1007/BF01762360}.
% The publisher's issue metadata gives December 1986; the work is often cited as 1987.

\bibitem{Kato1975}
T.~Kato,
The Cauchy problem for quasi-linear symmetric hyperbolic systems,
\emph{Arch. Ration. Mech. Anal.} \textbf{58} (1975), 181--205.
\href{https://doi.org/10.1007/BF00280740}{doi:10.1007/BF00280740}.

\bibitem{FriedrichRendall2000}
H.~Friedrich and A.~D.~Rendall,
The Cauchy problem for the Einstein equations,
in B.~G.~Schmidt (ed.), \emph{Einstein's Field Equations and Their Physical Implications},
Lecture Notes in Physics \textbf{540}, Springer, Berlin, 2000, pp.~127--223.
\href{https://arxiv.org/abs/gr-qc/0002074}{arXiv:gr-qc/0002074}.

\bibitem{LeFlochMardare2007}
P.~G.~LeFloch and C.~Mardare,
Definition and stability of Lorentzian manifolds with distributional curvature,
\emph{Portugaliae Math.} \textbf{64} (2007), 535--573.
\href{https://doi.org/10.4171/PM/1794}{doi:10.4171/PM/1794}.

\bibitem{ChoquetBruhatIsenbergPollack2007}
Y.~Choquet-Bruhat, J.~Isenberg, and D.~Pollack,
The constraint equations for the Einstein--scalar field system on compact manifolds,
\emph{Class. Quantum Grav.} \textbf{24} (2007), 809--828.
\href{https://arxiv.org/abs/gr-qc/0610045}{arXiv:gr-qc/0610045}.

\bibitem{IsenbergMaxwell2021}
J.~Isenberg and D.~Maxwell,
A phase space approach to the conformal construction of non-vacuum initial data sets in general relativity,
\href{https://arxiv.org/abs/2106.15027}{arXiv:2106.15027}, 2021.

\end{thebibliography}

\appendix

\section{Concrete reconstruction and smooth first-variation consistency}
\label{app:reconstruction}

This appendix is not a second derivation of the Lorentzian, spectral, or structural Standard-Model constructions. Its sole purpose is analytic: it exhibits one conventional finite regulator for which the reconstruction and first-variation consistency hypotheses of \Cref{def:regulator} can be verified directly, with $c_h(K)=O(h)$ on smooth bounded families. The argument is therefore a model realization of the continuum interface, not an alternative proof of the companion papers.

\subsection{Whitney maps and a signed Hodge realization}

On a contractible coordinate region choose a quasiuniform shape-regular simplicial mesh of size $h$. Write $C_h^k$ for oriented $k$-cochains. For an oriented simplex $\sigma=[i_0,\ldots,i_k]$ with barycentric coordinates $\lambda_i$, define the Whitney form
\begin{equation}
 w_\sigma=k!\sum_{j=0}^k(-1)^j\lambda_{i_j}
  \dd\lambda_{i_0}\wedge\cdots\wedge
  \widehat{\dd\lambda_{i_j}}\wedge\cdots\wedge\dd\lambda_{i_k},
 \qquad W_hc=\sum_\sigma c_\sigma w_\sigma.
 \label{eq:Whitney}
\end{equation}
The de Rham map $R_h\alpha=(\int_\sigma\alpha)_\sigma$ is used for smooth test forms. It is not asserted to be a bounded point- or trace-evaluation map on arbitrary $L^2$ forms.

\begin{lemma}[Commuting reconstruction and positive comparison norm]
\label{lem:Whitney}
The maps satisfy $R_hW_h=I$ and $\dd W_h=W_h\dd_h$. If
\[
 \norm c_h^2:=\int_K\abs{W_hc}_r^2\,\dV_0,
\]
then $W_h$ is an isometry onto its finite-element range. The reference mass matrix is positive definite. Shape regularity gives the usual cellwise scaling constants uniformly in $h$. Orthogonal projections $P_h^k$ onto these ranges converge strongly to the identity on $L^2$ $k$-forms.
\end{lemma}
\begin{proof}
Integration of \eqref{eq:Whitney} over oriented simplices gives the Kronecker pairing with the cochain basis. Differentiating the barycentric formula and using $\sum_i\lambda_i=1$ gives the coboundary identity; equivalently, integrate it and use Stokes on each simplex. The norm statement is its definition, and injectivity follows from $R_hW_h=I$. Scaling to a reference simplex and finite-dimensional norm equivalence give the uniform local mass bounds. Smooth forms are approximated by their Whitney interpolants in $L^2$; this follows by Taylor expansion on each cell and shape-regular scaling. Density of smooth forms and $\norm{P_h^k}\le1$ give strong convergence of the projections.
\end{proof}

This is the standard untwisted de Rham structure \cite{ArnoldFalkWinther2006,ArnoldFalkWinther2010}. Gauge curvature is treated separately: for a nonflat connection $D_A^2$ is curvature, and the ordinary commuting-complex identity is not a claim that $D_A^2=0$.

Let $g_h\to g$ uniformly with fixed signature and uniform inverse bounds. Using the fixed positive reference projections, define a discrete signed Hodge map by
\begin{equation}
 W_h^{4-k}\star_h^k
 =P_h^{4-k}\star_{g_h}W_h^k.
 \label{eq:Hodge-definition}
\end{equation}
For $W_h^kc_h\to\alpha$ in $L^2$,
\begin{equation}
 W_h^{4-k}\star_h^kc_h\longrightarrow\star_g\alpha\quad\text{in }L^2.
 \label{eq:Hodge-convergence}
\end{equation}
Indeed, subtract $\star_g\alpha$, use uniform boundedness of $\star_{g_h}$, uniform convergence of its algebraic coefficients, and strong convergence of $P_h^{4-k}$. The same proof applies to a fixed smooth metric variation because $\dot\star_g[k]$ is algebraic in $g,g^{-1},k$. The projections here are fixed by the positive reference metric, so no omitted derivative of a metric-dependent projection occurs.

This is a concrete signed-Hodge realization. It neither guarantees strong convergence of $W_hc_h$ nor excludes concentration of its quadratic products. Those are independent compactness requirements on the low-regularity branch; the regular a-priori branch of \Cref{thm:regular-branch} derives them from stronger Sobolev control. Exact mass integration can be replaced by quadrature only after its action and first-variation errors have been estimated.

\subsection{A finite smooth comparison reconstruction}

The following construction suffices to populate the consistency statement; it is not a claim of a canonical interpolation for every renewal graph. On a regular coordinate mesh let $I_hf$ be the continuous piecewise-affine interpolant of nodal samples. Choose a fixed smooth mollifier $\rho$ supported in the unit ball and put
\begin{equation}
 \mathcal J_hf=\rho_h*I_hf,
 \qquad \rho_h(x)=h^{-4}\rho(x/h).
 \label{eq:smooth-reconstruction}
\end{equation}
Work on a slightly larger region, so this local convolution creates no physical boundary condition. A fixed partition of unity handles a finite chart cover. The reconstruction is finite dimensional because $I_hf$ is determined by finitely many samples.

For $f$ in a bounded $C^3$ set, reference-cell interpolation and convolution give
\begin{equation}
 \norm{\mathcal J_hf-f}_{W^{1,\infty}}\le Ch\norm f_{C^2},
 \qquad \norm{\mathcal J_hf}_{W^{2,\infty}}\le C\norm f_{C^2}.
 \label{eq:smooth-reconstruction-estimate}
\end{equation}
For the second bound, differentiate the convolution of $\partial I_hf-\partial f$, whose $L^\infty$ norm is $O(h)$; the derivative of $\rho_h$ has $L^1$ norm $O(h^{-1})$. The remaining term is the convolution of $\partial^2f$. The same estimates hold for a smooth parameter derivative $\dot f$. Nondegeneracy of a smooth reference coframe is preserved for small $h$ by uniform approximation.

Use this map for coframe and matter comparison fields, with fixed local spin/gauge frames. For a smooth gauge connection let $U_\mu(x)$ denote backward parallel transport from $x+h e_\mu$ to $x$. Define
\begin{align}
 U_{\mu\nu}(x)
 &=U_\mu(x)U_\nu(x+h e_\mu)
          U_\mu(x+h e_\nu)^{-1}U_\nu(x)^{-1},\notag\\
 F^h_{\mu\nu}(x)&=h^{-2}\log U_{\mu\nu}(x),\notag\\
 K^h_\mu(x)&=h^{-1}\bigl(U_\mu(x)H(x+h e_\mu)-H(x)\bigr).
 \label{eq:lattice-comparison}
\end{align}
The logarithm is evaluated on the identity chart, valid uniformly for sufficiently small $h$ on the smooth bounded family. In \eqref{eq:lattice-comparison} the fields may be the reconstructed smooth fields $\mathcal J_hf$.

For spinors use backward spin--gauge transport $\mathcal U_\mu$ and the centered covariant difference
\begin{equation}
 \nabla^h_\mu\Psi(x)=\frac{
 \mathcal U_\mu(x)\Psi(x+h e_\mu)
 -\mathcal U_\mu(x-h e_\mu)^{-1}\Psi(x-h e_\mu)}{2h}.
 \label{eq:spin-difference}
\end{equation}
The dual-spinor difference uses the dual transport. Insert these quantities into the symmetric density \eqref{eq:dirac-density}, the Higgs density, and the Yang--Mills density, and sum with the cell-volume/Hodge coefficients computed from the same coframe. Define the gravitational comparison action by the integrated first-derivative torsion-free density \eqref{eq:palatini-first-order} on the same finite reconstruction. This is a finite-dimensional differentiable action; numerical integration, if used, must satisfy the stated derivative consistency.

\begin{proposition}[Smooth local action consistency]
\label{prop:mesh-consistency}
On a bounded smooth coframe/gauge/matter family with a common inverse-coframe bound, the preceding comparison action and its physical variations satisfy, uniformly for $\norm v_{C^{r_0}}\le1$,
\begin{equation}
 \abs{D\Sact_{b,h}[\mathcal I_hv]
       -D\Sact_{b,\theta_h}(z_h)[v]}\le C_Kh,
 \qquad b\in\{\mathrm g,\mathrm{SM}\}.
 \label{eq:mesh-C1}
\end{equation}
The reconstruction converges in the strong packet. The estimates are local smooth consistency estimates; they do not assert stability on arbitrary finite-energy link configurations.
\end{proposition}
\begin{proof}
The integral equation for parallel transport and Taylor expansion along an edge give
\[
 U_\mu(x)=I+hA_\mu(x)+O(h^2).
\]
Multiplying the four edge expansions cancels the terms linear in $h$ and leaves
\[
 U_{\mu\nu}(x)=I+h^2\Fcurv_{A,\mu\nu}(x)+O(h^3).
\]
The second-order contribution is $\partial_\mu A_\nu-\partial_\nu A_\mu+[A_\mu,A_\nu]$. Smoothness and the uniform identity chart therefore give $F^h=\Fcurv_A+O(h)$. Covariant Taylor expansion similarly gives $K^h=D_AH+O(h)$ and $\nabla^h\Psi=\nabla^{e,A}\Psi+O(h)$.

These expansions also hold after one field or metric variation. To see this without interchanging an uncontrolled limit and derivative, differentiate the transport equation itself. If $U(t,s)$ is transport along a path and $a$ is the connection variation, its derivative is a single integral of $U(t,u)a(u)U(u,s)$ with the fixed orientation sign. A second Taylor expansion gives the derivative of each displayed remainder, bounded by $Ch$ after the same normalizations. For a metric variation the spin-connection derivative contains only the first jet of the coframe variation; \eqref{eq:smooth-reconstruction-estimate} supplies the needed uniform bounds. The reconstruction of \eqref{eq:metric-lift} differs from the same lift evaluated at the reconstructed metric by $O(h)$ in the norms entering the density.

The pointwise densities and their first variations are finite sums of smooth coefficient functions and products of these quantities. Their difference is $O(h)\norm v_{C^{r_0}}$. A Riemann sum over cells of volume $O(h^4)$ has $O(h^{-4})$ terms on a fixed region, so the integrated error is $O(h)$, not an extensive error. The first-derivative gravitational expression obeys the same estimate directly from \eqref{eq:smooth-reconstruction-estimate}; exact integration eliminates its quadrature error. Differentiated consistent quadrature adds a further $O(h)$ term. This proves \eqref{eq:mesh-C1}.

The reconstructed fields converge in $W^{1,\infty}$ on the smooth family, and their curvature and covariant derivatives converge by the same expansions. This implies every norm in the strong packet on a compact region.
\end{proof}

The link-variable action has its usual finite gauge covariance. The comparison interpolation is only a means of comparing smooth fields and physical variations; arbitrary node-wise gauge transformations need not preserve that finite-dimensional interpolation chart. No exact lattice gauge identity is inferred merely from restriction to a non-invariant interpolation subspace. Equivariant reconstruction and nonlinear stability for arbitrary rough links are additional questions, distinct from the proved smooth estimate. Discretely gauge-invariant simplicial formulations provide another route, subject to their own consistency estimates \cite{ChristiansenHalvorsen2012}.

\section{Palatini products, metric variations, and the spin-current alternative}
\label{app:palatini}

\subsection{The first-order gravitational distribution}

Fix the orientation and curvature sign so that, for the torsion-free connection,
\begin{equation}
 B_{ab}(e)\wedge R^{ab}(\omega(e))=\Scal(g)\dV_g,
 \label{eq:B-normalization}
\end{equation}
where $B(e)$ is the appropriately normalized internal dual of $e\wedge e$. For a cutoff function $\chi$ equal to one near the support of a physical test,
\begin{align}
 \int\chi B\wedge R(\omega)
 =\int\bigl[-\dd(\chi B)\wedge\omega
              +\chi B\wedge\omega\wedge\omega\bigr].
 \label{eq:Palatini-IBP}
\end{align}
Exterior Stokes proves the identity for smooth compactly supported data. The terms involving $\dd\chi$ do not contribute to a variation supported where $\chi=1$.

For $e,e^{-1}\in L^\infty$ with $e\in H^1$, the Levi--Civita connection $\omega(e)$ is in $L^2$. Every term on the right of \eqref{eq:Palatini-IBP} is integrable, and its derivative in the lift \eqref{eq:metric-lift} is bounded by a constant times $\norm k_{C^1}$. This defines the local Einstein first-variation distribution. Different choices of $\chi$ give the same functional on tests supported in their common unit region. On smooth metrics the distribution is $(\Eins(g)+\Lambda g)\dV_g/(2\kappa)$.

This definition keeps the derivative of the connection in divergence form. It is consistent with the distributional-curvature framework for locally bounded metrics with locally square-integrable first derivatives and bounded inverse \cite{LeFlochMardare2007}. It does not assert a pointwise curvature tensor or an unqualified product of a rough coefficient with an arbitrary distribution.

\subsection{A weaker curvature passage}

\begin{proposition}[Weak-curvature Palatini passage]
\label{prop:weak-palatini}
On a compact four-dimensional region suppose
\begin{equation}
 e_h\to e\text{ in }L^2,\qquad
 \sup_h\norm{e_h}_6<\infty,\qquad
 R_h\rightharpoonup R\text{ in }L^p,
 \quad p>\tfrac32.
 \label{eq:Palatini-hypotheses}
\end{equation}
Assume $R$ is identified with the curvature of the limiting represented connection. Then $B(e_h)\to B(e)$ in $L^{p'}$, $p'=p/(p-1)$, and the Palatini pairings converge. The volume densities converge in $L^1$. If the represented coframe variations converge strongly in $L^2$, are uniformly bounded in $L^6$, and have compatible supports, their coframe first variations also converge. For a full metric variation, any connection-Euler pairing introduced by the lift must additionally tend to zero, or be retained as part of the limiting connection equation.
\end{proposition}
\begin{proof}
Interpolation between strong $L^2$ and bounded $L^6$ gives strong $L^q$ convergence for every $2\le q<6$; the lower exponents follow by finite volume. Since $p>3/2$, $2p'<6$. The product estimate gives
\[
 \norm{e_h\wedge e_h-e\wedge e}_{p'}
 \le C\norm{e_h-e}_{2p'}(\norm{e_h}_{2p'}+\norm e_{2p'})\to0.
\]
Pair with weak $L^p$ curvature. Taking $q=4$ gives convergence of the four-coframe volume product in $L^1$ by a four-factor telescoping estimate. The same argument applies to $\dot B_h$, a sum of one coframe and one coframe-variation factor, and to the volume variation. A connection variation produces its own Euler pairing after integration by parts. The stated extra condition, rather than coframe compactness alone, is what controls that term.
\end{proof}

The hypothesis identifying $R$ is substantive. A weak limit of matrix-valued two-forms is not automatically the curvature of the limiting connection. Nor does the Palatini lemma apply to a quadratic gauge stress by replacing one factor with a Hodge coefficient.

\subsection{Stationary elimination is not a universal torsion argument}

The independent connection Euler equation for a Palatini--Holst action has the form
\begin{equation}
 \mathcal C_e(T)=J_{\mathrm{spin}},
 \label{eq:Cartan-source}
\end{equation}
where $\mathcal C_e$ is the Cartan map determined by the coframe and the Palatini--Holst coefficients. On a nondegenerate chart with a uniformly bounded inverse of that map,
\[
 \norm T\le C\norm{J_{\mathrm{spin}}+r_{\mathrm{conn}}}.
\]
Thus zero spin current and vanishing connection residual imply torsionlessness. With nonzero current, the inverse instead reconstructs contorsion. Substitution into the total action retains the resulting spin-current interaction \cite{Kibble1961,HehlEtAl1976}. This is compatible with the connection alternative in the renewal gravitational construction, but differs from the metric action chosen in the main theorem.

There is also a separate normalization condition. A nonzero Holst coefficient alone does not imply a nonzero Einstein coefficient: on a torsion-free connection the Holst coframe equation reduces to the first Bianchi identity. The main action fixes a nonzero Palatini/Einstein coefficient $1/(2\kappa)$ before taking the limit. Pure Holst stationarity therefore cannot be used to infer \eqref{eq:Einstein-SM}.

\section{Nonminimal Higgs curvature and the assembled improvement}
\label{app:improvement}

The principal theorem concerns the minimally coupled classical action. A retained nonminimal source must be added before taking variations. For the convention
\begin{equation}
 \Sact_\xi=\xi\int f\Scal(g)\,\dV_g,
 \qquad f=H^\dagger H,
 \label{eq:xi-action}
\end{equation}
the compactly supported metric variation is
\begin{equation}
 D_g\Sact_\xi[k]
 =\xi\int\bigl[f\Eins_{\mu\nu}
       +(g_{\mu\nu}\Box_g-\nabla_\mu\nabla_\nu)f\bigr]
       k^{\mu\nu}\,\dV_g.
 \label{eq:improvement-variation}
\end{equation}
Consequently the stress contribution is minus twice the tensor in brackets multiplied by $\xi$. Formula \eqref{eq:improvement-variation} follows by varying $\sqrt{|g|}fR$ and integrating the two derivative terms in $\delta R$ twice. Compact support is essential to omitting the boundary contribution; with a physical boundary the corresponding scalar--GHY term must be included.

\begin{proposition}[An assembled strong--weak improvement criterion]
\label{prop:improvement}
Suppose $g_h\to g$ in $W^{1,\infty}$, $g_h\rightharpoonup g$ in $W^{2,p}$ for $1<p<\infty$, with uniform inverse bounds and fixed signature, and $H_h\to H$ in $L^{2p'}$. Then
\begin{equation}
 f_h\Eins(g_h)+(g_h\Box_{g_h}-\nabla^{h}\nabla^{h})f_h
 \longrightarrow
 f\Eins(g)+(g\Box_g-\nabla\nabla)f
 \label{eq:improvement-limit}
\end{equation}
in distributions, with the volume densities included in the pairing. If the minimal stress also converges and $\xi_h\to\xi$, the full improved metric variation converges.
\end{proposition}
\begin{proof}
Strong $L^{2p'}$ convergence gives $f_h\to f$ in $L^{p'}$. Curvature is linear in second metric derivatives with strongly convergent bounded coefficients, plus quadratic first-derivative terms. Hence $\Eins(g_h)\rightharpoonup\Eins(g)$ in $L^p$. This passes the first product. For the Hessian contribution, test the \emph{assembled} expression and integrate both derivatives onto the smooth test and metric-density coefficients. The resulting coefficient has a weak $L^p$ limit: its second metric derivatives occur linearly, and every lower-order coefficient converges strongly. Pairing with strong $f_h$ in $L^{p'}$ passes the entire expression. This argument does not separately multiply uncontrolled distributions.
\end{proof}

At the critical $p=2$ threshold the Higgs input is strong $L^4$, not merely a bounded $H^1$ sequence. If the assembled improvement has a subsequential distributional limit but not the value in \eqref{eq:improvement-limit}, write its defect as $\mathfrak I$. Then the Higgs stress defect becomes
\begin{equation}
 \mathfrak S_H=\mathfrak K_H-\lambda_Hg\nu_H-2\xi\mathfrak I.
 \label{eq:improved-defect}
\end{equation}
Existence of the assembled limit is an additional hypothesis; a critical Higgs norm alone does not supply it.

For completeness, if the scalar-curvature coefficient is
$F(H)=1/(2\kappa)+\xi f$ and is uniformly positive, the conformal change $\wt g=(F/\chi_*)g$ gives a constant Einstein coefficient $\chi_*>0$. For a real scalar-coordinate kinetic metric $G^0_{AB}$, the transformed kinetic metric is
\[
 \wt G_{AB}=\frac{\chi_*}{F}G^0_{AB}
             +\frac{3\chi_*}{F^2}\partial_AF\,\partial_BF.
\]
This identity follows from the four-dimensional conformal scalar-curvature formula and one integration by parts. With a boundary, the derivative term must be combined with the transformed GHY contribution. The change of frame reorganizes the principal coupling but supplies neither an effective-Planck lower bound nor hyperbolic smoothing. It therefore does not remove the compactness conditions of the main theorem.



\section{Coupled Einstein--matter actual-jet stability}
\label{app:coupled-stability}

This appendix proves the geometric-analysis stability estimate used in \Cref{thm:native-closure}.  Its variables are actual derivatives and curvatures of the reconstructed fields, so the argument is a coupled Einstein--matter estimate rather than an evolution of auxiliary jets followed by an identification step.

Choose an orthonormal frame $e_0,e_1,e_2,e_3$ for a regular metric chart, with $e_0$ future timelike, and Clifford multiplication $c_a$ satisfying $c_ac_b+c_bc_a=-2\eta_{ab}I$.  Let $\nabla$ be the torsion-free spin connection twisted by the retained gauge representation and write
\[
 \mathcal R^\nabla_{ab}=\mathcal R^S_{ab}+\rho(F_{ab}),\qquad
 \mathcal R^S_{ab}=\frac14\sum_{c,d}\epsilon_c\epsilon_dR_{ab cd}c_cc_d.
\]

\begin{lemma}[Twisted half-Ricci contraction]
\label{lem:half-ricci}
With the preceding conventions,
\begin{equation}
 \sum_b\epsilon_b c_b\mathcal R^\nabla_{ba}
 =\frac12\sum_b\epsilon_b\Ric_{ab}c_b
 +\sum_b\epsilon_bc_b\rho(F_{ba}).
 \label{eq:half-ricci}
\end{equation}
\end{lemma}
\begin{proof}
At one point choose a normal orthonormal frame.  Insert
$\mathcal R^S_{ba}=\frac14R_{ba cd}c^cc^d$ and use
$c_bc_cc_d=c_{[b}c_cc_{d]}-\eta_{bc}c_d+\eta_{bd}c_c-\eta_{cd}c_b$.
The fully alternating term vanishes by the first Bianchi identity, the last contraction vanishes by antisymmetry, and the remaining two contractions give one half of the Ricci contraction.  The gauge curvature acts on the twisting factor and contributes the last term directly.
\end{proof}

Put
\[
 r_D=\mathscr D_{g,A}\Psi-\mathscr M_{\boldsymbol Y}(H)\Psi,
 \qquad X_a=\nabla_{e_a}\Psi.
\]
\begin{proposition}[Residual-sensitive tangential spinor prolongation]
\label{prop:spinor-prolongation}
For each frame index $a$,
\begin{equation}
 (\mathscr D^{T^*}X)_a
 =\mathscr M X_a+\mathscr M_H[D_aH]\Psi+\nabla_a r_D
 +\frac12\Ric_{ab}c^b\Psi+c^b\rho(F_{ba})\Psi.
 \label{eq:spinor-prolongation}
\end{equation}
Moreover the normal jet is algebraic,
\begin{equation}
 X_0=c_0c_iX_i-c_0\mathscr M\Psi-c_0r_D,
 \label{eq:normal-spinor-jet}
\end{equation}
so the tangential prolonged equations require spatial derivatives of $r_D$ but no time derivative of $r_D$.
\end{proposition}
\begin{proof}
Commute $\nabla_a$ with the twisted Dirac operator and use \Cref{lem:half-ricci}.  The covariant product rule differentiates the affine mass map to $\mathscr M_H[D_aH]$.  The Dirac equation itself gives \eqref{eq:normal-spinor-jet}.  In a moving frame, the normal covector component enters the tangential covector connection only algebraically, so substituting \eqref{eq:normal-spinor-jet} introduces no normal derivative of the residual.
\end{proof}

Let
\[
 \mathcal E_{ab}=G_{ab}+\Lambda g_{ab}-\kappa T^{\rm SM}_{ab},
 \qquad
 \mathcal E^{\rm tr}_{ab}=\mathcal E_{ab}-\tfrac12g_{ab}\operatorname{tr}_g\mathcal E.
\]
Then
\begin{equation}
 \Ric_{ab}=\kappa\left(T^{\rm SM}_{ab}-\tfrac12g_{ab}\operatorname{tr}_gT^{\rm SM}\right)
 +\Lambda g_{ab}+\mathcal E^{\rm tr}_{ab}.
 \label{eq:trace-reversal-residual}
\end{equation}
The complete torsion-free Standard-Model stress is algebraic in the actual first jets
\[
 (g,F,H,DH,\Psi,\bPsi,X_a,\bar X_a)
\]
and contains no derivative of $F$, $DH$ or $X$.  After \eqref{eq:normal-spinor-jet} it is a smooth algebraic function of the tangential actual-jet state and is at most linear in the Dirac residual.

In harmonic coordinates and temporal internal gauge define, for every coordinate metric component $u=g_{\mu\nu}$,
$p_u=e_0u$, $q_{u,a}=e_au$, and
\[
 E_a=F(e_0,e_a),\quad B_a=-\tfrac12\varepsilon_{abc}F(e_b,e_c),
 \quad \Pi=D_{e_0}H,\quad Q_a=D_{e_a}H.
\]
The realified actual-jet state is
\begin{equation}
 \mathcal U=(g_{\mu\nu},p_u,q_{u,a},A_i,E_a,B_a,H,\Pi,Q_a,
              \Psi,X_1,X_2,X_3),
 \label{eq:actual-jet-state}
\end{equation}
with the physical dual spinor included through the fixed adjoint identification.  Let $r^A$, $r_H$ and $r_D$ denote the physical Yang--Mills, Higgs and Dirac residuals.

\begin{proposition}[Block-symmetric actual-jet equation]
\label{prop:actual-jet-writer}
On a compact nondegenerate chart, every smooth actual field tuple in harmonic coordinates and temporal internal gauge satisfies
\begin{equation}
 \partial_t\mathcal U+\sum_{j=1}^3\mathcal A^j(g)\partial_j\mathcal U
 =\mathcal F(\mathcal U)+\mathcal B(\mathcal U)
    (\mathcal E,r^A,r_H,r_D,\bar r_D)
 +\sum_{j=1}^3\mathcal Q^j(g)\partial_j(r_D,\bar r_D),
 \label{eq:actual-jet-writer}
\end{equation}
where the matrices $\mathcal A^j$ are symmetric after realification and all remaining coefficients are smooth on bounded jet sets.  No derivative of $\mathcal E$, $r^A$ or $r_H$ occurs.
\end{proposition}
\begin{proof}
Harmonic reduction gives a wave equation for each metric component with source \eqref{eq:trace-reversal-residual}; first-order reduction yields the standard symmetric wave blocks for $(p_u,q_u)$.  The Yang--Mills equation together with the Bianchi identity yields the symmetric Maxwell blocks for $(E,B)$.  The Higgs equation and the commutator identity for covariant derivatives give the same wave blocks for $(\Pi,Q)$.  The Dirac head equation has Hermitian principal matrices $-c_0c_a$.  For the tangential spinor jets use \Cref{prop:spinor-prolongation}, substitute \eqref{eq:trace-reversal-residual} and then \eqref{eq:normal-spinor-jet}.  The half-Ricci identity is exactly what prevents an uncontrolled Weyl-curvature derivative or differentiated matter stress from entering the spinor block.  Frame and connection coefficients are smooth functions of the actual state.  These blocks are symmetric in the fixed positive component norm.
\end{proof}

\begin{theorem}[Coupled actual-jet stability and first-exit closure]
\label{prop:coupled-bootstrap}
Let $z_*$ be a regular exact Einstein--Standard-Model solution on $Q=[0,T]\times\mathbb T^3$ in the preceding gauges with $\mathcal U_*\in C_tH_x^{k+1}$, $k\ge4$.  Let $\widehat z$ be any smooth actual field tuple on the same slab and chart, and define
\[
 d=\norm{\widehat{\mathcal U}(0)-\mathcal U_*(0)}_{H^k}
 +\norm{\mathcal R_{\rm B}(\widehat z)}_{L^2_tH_x^k}
 +\norm{\mathcal R_{\rm D}(\widehat z)}_{L^2_tH_x^{k+1}}
 +\norm{C(\widehat g)}_{L^2_tH_x^{k+1}}
 +\norm{\partial_tC(\widehat g)}_{L^2_tH_x^k}.
\]
There are constants $d_*,C_*>0$, depending only on the reference and fixed chart margins, such that $d\le d_*$ implies
\begin{align}
 \norm{\widehat{\mathcal U}-\mathcal U_*}_{C_tH_x^k}&\le C_*d,\label{eq:bootstrap-state}\\
 \norm{\Riem(\widehat g)-\Riem(g_*)}_{L^2_tH_x^{k-1}}
 +\norm{T^{\rm SM}(\widehat z)-T^{\rm SM}(z_*)}_{L^2_tH_x^k}
 &\le C_*d.\label{eq:bootstrap-curvature}
\end{align}
No all-time $H^k$ bound for $\widehat{\mathcal U}$ is assumed.  Existence of $\widehat z$ on the slab is assumed.
\end{theorem}
\begin{proof}
On the bootstrap interval $\|\widehat{\mathcal U}-\mathcal U_*\|_{H^k}<1$, subtract \eqref{eq:actual-jet-writer} for the two states and put the approximate-state principal matrices on the left.  Since $k>5/2$, $H^k(\mathbb T^3)$ is an algebra and the standard commutator estimate
\[
 \|[\partial^\alpha,a]\partial_jW\|_2
 \le C_k(\|\nabla a\|_\infty\|W\|_{H^k}
       +\|a\|_{H^k}\|\nabla W\|_\infty)
\]
controls all differentiated principal terms.  The reference has one additional derivative, so coefficient differences multiply only controlled reference derivatives.  Symmetry of the principal matrices and integration by parts therefore give a positive energy $E_k^{1/2}\simeq\|W\|_{H^k}$ satisfying
\[
 \frac{\dd}{\dd t}E_k^{1/2}
 \le C E_k^{1/2}
 +C\bigl(\|\mathcal R_{\rm B}\|_{H^k}
          +\|\mathcal R_{\rm D}\|_{H^{k+1}}
          +\|C\|_{H^{k+1}}+\|\partial_tC\|_{H^k}\bigr).
\]
Gronwall and Cauchy--Schwarz yield $\sup_t\|W(t)\|_{H^k}\le C_*d$ while the bootstrap holds.  Choosing $d_*$ so that $C_*d_*<1/2$ rules out the first exit and proves \eqref{eq:bootstrap-state}.  The equation controls $\partial_tW$ one order lower.  The metric second derivatives are recovered from the metric first-jet rows, so curvature differs by the same forcing-controlled amount in $L^2_tH^{k-1}$.  The complete stress is algebraic in the controlled actual jets, with the normal spinor derivative supplied by \eqref{eq:normal-spinor-jet}; this gives \eqref{eq:bootstrap-curvature}.
\end{proof}




\section{Native source selection and physical reader certificates}
\label{app:native-source-selection}

This appendix proves the finite source-level results used in \Cref{sec:native-selection}.  All probabilities refer to the actual accepted table; failures are retained, and no transition law is conditioned on success.

\subsection{Stopped signed-action drift and same-record selection}

Let $Z_j$ denote the diagnostic stopped action value: on a legal transition it equals $a_h(X_j)$, while after the first attempted illegal transition it retains the last legal value.  Shifting the action by the constant $A_h^-$ gives $0\le Z_j\le D_h^A$ without changing any increment or Euler covector.  Conditional on an active legal row,
\begin{equation}
 \mathbb E[Z_{j+1}-Z_j\mid X_j=x]=\mathfrak d_{h,j}(x).
 \label{eq:app-stopped-drift}
\end{equation}
Summing \eqref{eq:source-drift-ineq} against the surviving subprobability masses and telescoping gives
\begin{equation}
 \kappa_A\delta_h\sum_{j<J,x}\mu_j(x)g_h(x)^2
 \le D_h^A+\delta_h\sum_{j<J,x}\mu_j(x)r_{h,j}(x),
 \label{eq:app-slope-telescope}
\end{equation}
which is \eqref{eq:source-slope-occupation}.

For the selection step define
\[
 \mathcal S_j=\frac{g_h(X_j)^2}{s_h^2}
 +\frac{V_h(X_j)}{R_h^2}
 +\frac{W_h(X_j)}{\alpha_h^2}.
\]
On a surviving prefix, if $j_*$ minimizes $\mathcal S_j$ and $J^{-1}\sum_{j<J}\mathcal S_j\le1$, then every term in \eqref{eq:source-selected} is bounded by its threshold.  Hence selection failure on survival implies that the mean score exceeds one.  Markov's inequality, \eqref{eq:source-slope-occupation}, and the definitions \eqref{eq:source-occupations} give exactly \eqref{eq:source-selection-failure}.  Notice that no union bound over all candidates is used; the high-order energy and alignment conditions are required only of the selected record.

\subsection{Exact reader constant and sharp Fourier transfer}
\label{app:native-reader-proof}

For the affine reader \eqref{eq:source-reader}, collect the high-frequency coefficient maps in
\[
 \mathsf T_{h,K,m}\xi=(T_{h,\ell}\xi)_{|\ell|_\infty>K},
 \qquad w_\ell=(1+|\ell|_1/K)^m.
\]
If $\xi\in\ker\mathsf Q_h$, the energy side of \eqref{eq:source-reader-tail} vanishes; hence \eqref{eq:source-reader-kernel} is necessary.  Conversely restrict to $(\ker\mathsf Q_h)^\perp$ and write $\xi=\mathsf Q_h^{\dagger/2}y$.  Then
\[
 \norm y^2=\ip{\xi}{\mathsf Q_h\xi},
\]
and weighted block-$\ell^1$ duality gives
\begin{align}
 \sup_{\ip{\xi}{\mathsf Q_h\xi}\le1}
 \sum_{|\ell|_\infty>K}w_\ell\norm{T_{h,\ell}\xi}
 &=\norm{\mathsf T_{h,K,m}\mathsf Q_h^{\dagger/2}}_{2\to\ell^1(w)}.
 \label{eq:app-reader-dual}
\end{align}
This proves \Cref{prop:native-reader}; the affine offset is handled by the triangle inequality.

For completeness, on the normalized four-torus Parseval and Cauchy--Schwarz give
\begin{equation}
 \tau_{h,m}(K;u)
 \le A_{N,K,m,q}\norm u_{H^q},\qquad
 A_{N,K,m,q}^2=
 \sum_{K<|\ell|_\infty\le N}
 \frac{(1+|\ell|_1/K)^{2m}}{(1+|\ell|_2^2)^q}.
 \label{eq:app-exact-tail-sum}
\end{equation}
Dyadic shells in four dimensions contain $O(R^4)$ lattice points.  On a shell of radius $R\asymp2^aK$, the contribution to \eqref{eq:app-exact-tail-sum} is bounded by $CK^{-2m}R^{2m-2q+4}$.  Thus
\begin{equation}
 \tau_{h,m}(K;u)\le C_{q,m}K^{2-q}\norm u_{H^q}
 \qquad(q>m+2),
 \label{eq:app-sharp-tail}
\end{equation}
and the exponent is sharp on shell-supported examples.  At $q=m+2$ the dyadic contributions are comparable and the constant diverges logarithmically with the ultraviolet cutoff.  Combining \eqref{eq:source-reader-LMI} with \eqref{eq:app-sharp-tail} proves \Cref{cor:native-reader-sobolev}.

\subsection{A directly verifiable action-drift mechanism}

The drift hypothesis \eqref{eq:source-drift-ineq} is itself testable against the actual transition law and the same action.  Let $E=\nabla\mathcal A$ in the physical mass norm and $\Delta=Y-x$ for one legal transition.  Suppose
\begin{equation}
 m(x):=\mathbb E[\Delta\mid x]
 =-\delta M(x)E(x)+\delta e(x),
 \qquad \frac{M+M^*}{2}\succeq\lambda I,
 \label{eq:app-mean-incidence}
\end{equation}
with $\lambda>0$.  Taylor's integral identity gives
\begin{equation}
 \mathbb E[\mathcal A(Y)-\mathcal A(x)\mid x]
 =\ip{E(x)}{m(x)}+\mathcal W(x),
 \quad
 \mathcal W(x)=\mathbb E\!\left[\int_0^1(1-t)D^2\mathcal A(x+t\Delta)[\Delta,\Delta]\,dt\,\middle|\,x\right].
 \label{eq:app-signed-work}
\end{equation}
Therefore
\begin{equation}
 \mathfrak d(x)\le-\frac\lambda2\delta\norm E^2
 +\delta\left(\frac{\norm e^2}{2\lambda}+\frac{[\mathcal W(x)]_+}{\delta}\right).
 \label{eq:app-work-drift}
\end{equation}
Only the signed action work of the innovation is charged.  In particular, if the actual transition factors as $K(x,dy)=\int D(x,dz)N_x(z,dy)$ and the second kernel is action-neutral,
\begin{equation}
 \int\mathcal A(y)N_x(z,dy)=\mathcal A(z),
 \label{eq:app-action-neutral}
\end{equation}
then its stochastic variance contributes nothing to the action-drift remainder.  The selected record itself, including that retained noise, is still tested by $V_h$ and $W_h$.  These identities do not authorize inserting a gradient step or an action-neutral component into a source that does not actually possess them.

\subsection{Polynomial regime}
\label{app:native-polynomial-proof}

Under \eqref{eq:selected-polynomial-source}--\eqref{eq:selected-polynomial-window}, the five terms in \eqref{eq:source-selection-failure} are respectively $O(h^{2\epsilon})$, $O(h^{2\epsilon})$, $O(h^{2a})$, $O(h^{2\epsilon})$, and the declared exit term.  Hence the failure probability is $O(h^{2a}+h^{2\epsilon})$.  The Sobolev reader gives
\begin{equation}
 T_{h,j}\le C h^{-a}K_h^{2-q},\qquad j=2,k+2.
 \label{eq:app-polynomial-tail}
\end{equation}
Consequently
\[
 \bar\varepsilon_{c,h}
 \le C\bigl(h^b+hK_h^3+h^{-a}K_h^{4-q}\bigr),
 \qquad \bar L_{h,k}\le C(1+|\log h|),
\]
and the terms in $\bar F_{h,k}$ have orders
\[
 h^bK_h^{k+1},\qquad hK_h^{k+4},\qquad
 h^{-a}K_h^{k+5-q},\qquad h^{-a}K_h^{k+4-q},
\]
up to the fixed logarithmic factor.  The fourth exponent is larger than the third by $\beta$; adding $\alpha_h=h^c$ gives precisely \eqref{eq:selected-polynomial-rate}.  Positive powers of a dyadic spacing times a fixed logarithmic power are summable.

\section{Gauge-covariant realization of the weak-sector reader}
\label{app:gauge-reader}

The Fourier reader used in the quantitative theorem is a coordinate object, but its control need not originate from ordinary derivatives of an arbitrary gauge frame.  We record a finite gauge-covariant entrance on a nonvanishing Higgs chart.  This appendix concerns the weak $SU(2)$ sector only; it is not a substitute for independent color, determinant-line or gravitational reader estimates.

On a periodic grid let $U_i(x)\in SU(2)$ and define covariant shifts and differences
\begin{equation}
 T_i^UF(x)=\rho(U_i(x))F(x+e_i),\qquad
 D_i^UF=h^{-1}(T_i^UF-F),\qquad D_I^U=D_{i_1}^U\cdots D_{i_\ell}^U.
 \label{eq:gauge-covdiff}
\end{equation}
For a fixed integer $L$ put
\begin{equation}
 \mathcal J_{h,L}(F;U)^2
 =\sum_{\ell=0}^L\ \sum_{|I|=\ell}\norm{D_I^UF}_{2,h}^2.
 \label{eq:gauge-jet-energy}
\end{equation}
Expanding the ordered product $h^{-\ell}(T_{i_1}^U-I)\cdots(T_{i_\ell}^U-I)$ gives, for each word $I$,
\begin{equation}
 D_I^UF(x)=h^{-\ell}\sum_{S\subset\{1,\ldots,\ell\}}(-1)^{\ell-|S|}F_S(x),
 \label{eq:gauge-path-expansion}
\end{equation}
where all transported endpoints $F_S(x)$ transform by the same unitary at $x$.  Hence $|D_I^UF(x)|^2$ is a fixed contraction of the positive mixed path Gram $G_I(x)_{ST}=\langle F_S(x),F_T(x)\rangle$.  The energy \eqref{eq:gauge-jet-energy} is therefore gauge invariant and can be audited from same-record mixed path data.

The discrete Kato inequality
\begin{equation}
 |\delta_i|F|(x)|\le |D_i^UF(x)|
 \label{eq:gauge-Kato}
\end{equation}
and ordinary Sobolev inequalities on the fixed grid imply, for $d\le4$,
\begin{equation}
 \norm F_{\infty,h}\le C_d\mathcal J_{h,3}(F;U),\qquad
 \max_{|I|\le L-3}\norm{D_I^UF}_{\infty,h}
 \le C_{d,L}\mathcal J_{h,L}(F;U),
 \label{eq:gauge-cov-embed}
\end{equation}
with constants independent of the links.

Assume now $|H_x|\ge\rho_*>0$.  Let $Q_x\in SU(2)$ be the algebraic frame sending $e_1$ to $H_x/|H_x|$ and put
\begin{equation}
 W_i(x)=Q_x^*U_i(x)Q_{x+e_i},\qquad
 a_i=h^{-1}(W_i-I),\qquad J_I=Q_x^*D_I^UH.
 \label{eq:gauge-normalized}
\end{equation}
The exact identities
\begin{equation}
 \delta_iJ_I=J_{iI}-a_iT_iJ_I,
 \qquad |H_{x+e_i}|=\bigl||H_x|e_1+hJ_i(x)\bigr|
 \label{eq:gauge-closed}
\end{equation}
close the ordinary difference hierarchy in terms of covariant jets.  On a fixed lower-radius chart, repeated finite-difference chain rules therefore give
\begin{equation}
 \mathcal J_{h,s+4}(H;U)+\sum_b\mathcal J_{h,s+3}(F_b;U)\le R
 \Longrightarrow
 \norm{|H|}_{C_h^{s+1}}
 +\sum_i\norm{h^{-1}\log W_i}_{C_h^s}
 +\sum_b\norm{Q^*F_b}_{C_h^s}
 \le C_{s,\rho_*,R}.
 \label{eq:gauge-normalized-reader}
\end{equation}
No derivative of the original $Q_x$ and no smallness of the original links is assumed.

The algebraic Higgs frame need not be temporal gauge.  Starting with $g_{0,x}=I$ and defining causally
\begin{equation}
 g_{j+1,x}=g_{j,x}W_0(j,x)
 \label{eq:gauge-temporal}
\end{equation}
sets the transformed temporal links to the identity.  Iterated product rules and discrete Gronwall show that one additional covariant derivative controls the transformed spatial logarithmic links through order $s$ on a fixed physical-time interval.  All transformed matter fields retain the same order of control.  Parseval and the discrete-difference symbol estimate then yield an all-mode $H^q$ bound and, by \eqref{eq:app-sharp-tail}, the tail estimate \eqref{eq:gauge-robust-tail}.

Finally, exact finite gauge invariance prevents this normalization from hiding the action row.  If $E_h$ is the Riesz gradient of the \emph{full} physical action in the mass norm, then for every fixed site gauge transformation $g$,
\begin{equation}
 E_h(g\cdot z)=T_gE_h(z),\qquad \norm{E_h(g\cdot z)}_{0,h}=\norm{E_h(z)}_{0,h}.
 \label{eq:gauge-full-covector}
\end{equation}
When $g=g(z)$ is the algebraic Higgs normalization, differentiating $g(z)\cdot z$ adds a vertical gauge tangent; the exact finite Ward identity annihilates that tangent, so the complete first variation is unchanged.  On the normalized logarithm chart $W=e^{h\mathsf A}$ with $h\norm{\mathsf A}\le c$, the physical left link variation and $\dot{\mathsf A}$ are related by
\begin{equation}
 J_h(\mathsf A)\dot{\mathsf A}
 =\int_0^1\operatorname{Ad}_{e^{th\mathsf A}}\dot{\mathsf A}\,dt,
 \qquad \norm{J_h-I}\le Cc,
 \label{eq:gauge-log-source}
\end{equation}
so the logarithmic and physical link-source norms are uniformly equivalent for fixed sufficiently small $c$.  A derivative restricted to the radial Higgs slice alone is not used as a substitute for this full covector.

\section{Generated finite dynamics from constrained data}
\label{app:generated-dynamics}

This appendix gives an independent nonemptiness construction for the controlled finite class.  It is local in time and deliberately high-regularity; sharp thresholds are not claimed.  Its role is not to replace the native source theorem of \Cref{thm:native-selected-closure}, but to show that the quantitative hypotheses admit a fully finite realization generated from constrained finite initial data.  The construction uses the positive symmetric actual-jet system of \Cref{app:coupled-stability} and the standard conformal strategy for the Einstein constraints \cite{ChoquetBruhatIsenbergPollack2007,IsenbergMaxwell2021}.

\begin{theorem}[Fully finite approximation branch]
\label{thm:generated-dynamics}
Fix integers $s\ge5$, $p\ge1$, and put $q=s+p+8$.  For the constraint-compatible initial class constructed below, each spatial cutoff $N$ has finite initial data $U_{0,N}$ satisfying
\begin{equation}
 \sup_N\norm{U_{0,N}}_{H^q}\le M_q,
 \qquad
 \norm{U_{0,N}-U_0}_{H^{s+4}}\le C N^{-p}.
 \label{eq:generated-initial-rate}
\end{equation}
Let
\begin{equation}
 G_N(U)=P_N\{F(U)-A^i(U)\partial_iU\},
 \qquad
 U^{j+1}=U^j+\tau G_N\!\left(\frac{U^{j+1}+U^j}{2}\right),
 \quad U^0=U_{0,N}.
 \label{eq:generated-midpoint}
\end{equation}
If
\begin{equation}
 \tau(N+1)\le c_*,\qquad \tau\le\tau_*,
 \label{eq:generated-CFL}
\end{equation}
then the fully finite recursion exists on a common interval $[0,T]$ and
\begin{equation}
 \max_j\norm{U^j}_{H^q}\le R_q.
 \label{eq:generated-uniform}
\end{equation}
Its cubic-Hermite physical reconstruction $z_{N,\tau}$ satisfies
\begin{align}
 \norm{z_{N,\tau}-z}_{C_tH^{s+2}}
 +\norm{\partial_tz_{N,\tau}-\partial_tz}_{C_tH^{s+1}}
 &\le C(N^{-p}+\tau^2),\label{eq:generated-first}\\
 \norm{\Riem(g_{N,\tau})-\Riem(g)}_{L^\infty_tH^s}
 +\norm{\mathcal R_{\rm B}(z_{N,\tau})}_{L^\infty_tH^s}
 &\le C(N^{-p}+\tau),\label{eq:generated-bosonic}\\
 \norm{\mathcal R_{\rm D}(z_{N,\tau})}_{C_tH^{s+1}}
 &\le C(N^{-p}+\tau^2).\label{eq:generated-Dirac}
\end{align}
Moreover the finite comparison action on the same generated record obeys
\begin{equation}
 |D S^{\rm cmp}_{N,\tau}[V]|
 \le C(N^{-p}+\tau)\norm V_{0,\tau}.
 \label{eq:generated-stationarity}
\end{equation}
The continuum solution $z$ is used only for the a-posteriori estimate; the finite initial solver and recursion do not sample its future values.
\end{theorem}

\subsection{A nonempty constrained initial class and its finite approximation}

Work on $\mathbb T^3$ near a smooth expanding constant background.  Fix a Higgs position $H_*$, a nonzero radial Higgs normal momentum $p$, and a smooth gauge background $A^*$ with bounded magnetic curvature.  Put
\begin{equation}
 A_*^0=\kappa p^2,
 \qquad G_*^0=\kappa|D_{A^*}H_*|^2,
 \qquad B_*^0=\kappa|B_{A^*}|^2,
 \label{eq:generated-background-coeffs}
\end{equation}
and choose the constant mean curvature so that the unperturbed Hamiltonian constraint holds.  In the conformal scalar equation the derivative at the unperturbed solution $w=1$ is
\begin{equation}
 L_*=-8\Delta+\mu_*,
 \qquad
 \mu_*=12A_*^0+4G_*^0+8B_*^0>0.
 \label{eq:generated-Lstar}
\end{equation}
Thus $L_*:H^{r+2}\to H^r$ is an isomorphism including the zero Fourier mode.

The Gauss and vector-momentum equations are solved first by their fixed elliptic inverses.  For a small transverse--traceless tensor seed $\sigma$ and a small standing spinor seed $\chi=\varepsilon f\chi_*$ with real $f$, the canonical spatial spinor momentum vanishes and the resulting source coefficients in the scalar conformal equation obey
\begin{equation}
 \norm{\delta\mathsf A}_{H^r}
 +\norm{\delta\mathsf B}_{H^r}
 +\norm{\mathsf Y}_{H^r}
 \le C\bigl(\norm\sigma_{H^{r+1}}^2+\norm\chi_{H^{r+1}}^2\bigr).
 \label{eq:generated-source-small}
\end{equation}
The Hamiltonian equation can be written
\begin{equation}
 -8\Delta w-\mathsf A w^{-7}-G_*^0w
 -\mathsf B w^{-3}-\mathsf Y w^{-1}+C_*^0w^5=0,
 \label{eq:generated-Hamiltonian}
\end{equation}
where $C_*^0$ is fixed by the background constraint.

\begin{proposition}[Constructive constrained initial data]
\label{prop:generated-initial}
For every integer $r\ge2$, sufficiently small $\sigma,\chi$ as above determine a unique solution $w$ of \eqref{eq:generated-Hamiltonian} in a fixed $H^{r+2}$ neighborhood of one, with
\begin{equation}
 \norm{w-1}_{H^{r+2}}
 \le C\bigl(\norm\sigma_{H^{r+1}}^2+\norm\chi_{H^{r+1}}^2\bigr),
 \qquad \tfrac12<w<\tfrac32.
 \label{eq:generated-w-bound}
\end{equation}
Together with the Gauss and vector solutions this gives exact Einstein, gauge and matter initial constraints.  The class contains nonzero-spinor data and, for nonconstant $f$, genuinely inhomogeneous metric response.

If the seeds have $p$ additional derivatives, Fourier projection of the seeds followed by the same finite elliptic construction gives finite initial tuples $U_{0,N}$ satisfying \eqref{eq:generated-initial-rate}; the Gauss and vector constraints remain exact and the scalar Hamiltonian residual is $O(N^{-p})$ in the indicated lower Sobolev norm.  No spacetime solution is used in this finite initial solver.
\end{proposition}
\begin{proof}
Write $u=w-1$ and subtract the linear term $L_*u$ from \eqref{eq:generated-Hamiltonian}.  The remaining pure-background nonlinearities are quadratic in $u$, while the perturbed source coefficients satisfy \eqref{eq:generated-source-small}.  Hence
\[
 u=L_*^{-1}\bigl[\delta\mathsf A(1+u)^{-7}
 +\delta\mathsf B(1+u)^{-3}+\mathsf Y(1+u)^{-1}-R_0(u)\bigr].
\]
On a small $H^{r+2}$ ball, Sobolev multiplication and the convergent power series for the negative integer powers give
\[
 \norm{\mathcal T(0)}_{H^{r+2}}\le C\eta,
 \qquad
 \norm{\mathcal T(u)-\mathcal T(v)}_{H^{r+2}}
 \le C(\eta+R)\norm{u-v}_{H^{r+2}},
\]
where $\eta=\norm\sigma^2+\norm\chi^2$.  Banach contraction proves \eqref{eq:generated-w-bound}.  The Gauss and vector equations use bounded elliptic inverses of their compatible sources, so all initial constraints hold after substitution into the conformal identities.  For the finite construction, project the finite seeds, solve those elliptic equations exactly, and apply the projected contraction to the scalar equation.  The contraction constant remains uniform because $P_N$ is a contraction in Sobolev space and commutes with $L_*$.  Subtracting the exact and projected fixed-point equations and using the ordinary Fourier tail estimate gives $O(N^{-p})$.  The same argument applies to the derived harmonic initial completion and actual first jets.
\end{proof}

\begin{lemma}[Physical identification of the symmetric extension]
\label{lem:generated-physical-identification}
Let $U_0$ belong to the constrained Sobolev class of \Cref{prop:generated-initial}, with sufficiently many finite derivatives for the estimates below.  The independent symmetric system
\begin{equation}
 U_t+A^i(U)\partial_iU=F(U)
 \label{eq:generated-extension}
\end{equation}
has a common local Sobolev solution, and on a possibly smaller interval its head fields satisfy the complete torsion-free Einstein--Standard-Model equations and the harmonic/temporal-gauge constraints.  Thus the physical solution $z$ used in \Cref{thm:generated-dynamics} is identified from the constrained initial data; it is not an external trajectory supplied to the finite recursion.
\end{lemma}
\begin{proof}
Approximate the finite-Sobolev seeds by Fourier-polynomial seeds and apply the exact, unprojected contraction of \Cref{prop:generated-initial} to each approximation.  This gives a sequence of analytic initial tuples satisfying all physical constraints exactly and converging to $U_0$ in the Sobolev topology.  For each analytic tuple, the local Einstein--Standard-Model Cauchy problem in harmonic coordinates and temporal internal gauge has a physical analytic germ.  Its actual first jets satisfy \eqref{eq:generated-extension} by \Cref{prop:actual-jet-writer}.  Uniqueness for the symmetric system therefore identifies the independent analytic solution with those actual jets on their common local interval.

The symmetric energy estimate for \eqref{eq:generated-extension} depends only on a bounded Sobolev ball and the fixed nondegenerate chart, not on the analytic radius.  Hence the analytic solutions exist on a common smaller interval and converge in one lower Sobolev order to the solution with initial datum $U_0$.  The defining-jet identities, Einstein and matter residuals, Gauss constraint and harmonic constraint vanish on every analytic approximant.  The convergence of fields, first derivatives and the equation then passes these identities distributionally to the Sobolev limit; the algebraic reconstruction identifies the actual connection, curvature and complete stress.  This proves physical identification.  It is an analytic-core argument for the constrained class, not a claim that arbitrary auxiliary initial jets are physical.
\end{proof}

\subsection{Finite symmetric evolution and midpoint construction}

Let the coupled actual-jet system of \Cref{prop:actual-jet-writer} be written on its fixed chart as
\begin{equation}
 U_t+A^i(U)\partial_iU=F(U),
 \qquad G(U)=F(U)-A^i(U)\partial_iU.
 \label{eq:generated-symmetric-system}
\end{equation}
For the finite initial data above define the semidiscrete system
\begin{equation}
 \dot U_N=P_NG(U_N),\qquad U_N(0)=U_{0,N}.
 \label{eq:generated-Galerkin}
\end{equation}
The symmetry of the principal matrices implies that the $H^q$ energy estimate for \eqref{eq:generated-Galerkin} has constants independent of $N$; the orthogonal projection disappears from the energy pairing.  A first-exit argument therefore gives a common existence interval and a uniform $H^q$ bound.  On that interval,
\begin{equation}
 \max_{0\le a\le2}
 \norm{\partial_t^a(U_N-U)}_{C_tH^{s+4-a}}
 \le C N^{-p},
 \label{eq:generated-Galerkin-rate}
\end{equation}
where $U$ is the physical Sobolev solution with initial datum $U_0$.  This comparison identifies the limit only after the finite ODE has been constructed.

For the fully finite evolution use \eqref{eq:generated-midpoint}.  On $\operatorname{Ran}P_N$ the inverse inequality gives
\[
 \norm{G_N(V)-G_N(W)}_{H^q}
 \le C_R(N+1)\norm{V-W}_{H^q}.
\]
Thus \eqref{eq:generated-CFL} makes the midpoint fixed-point map a contraction.  The exact discrete symmetric energy identity then removes the apparent $N$-dependence from the propagated bound and discrete Gronwall gives \eqref{eq:generated-uniform}.  If two midpoint trajectories have a normalized forcing difference $r_j$, the same energy calculation at any $3\le m\le q-1$ gives
\begin{equation}
 \max_{j\tau\le T}\norm{U^j-V^j}_{H^m}
 \le C_T\left(\norm{U^0-V^0}_{H^m}
 +\tau\sum_j\norm{r_j}_{H^m}\right),
 \label{eq:generated-causal}
\end{equation}
with $C_T$ independent of $N,\tau$.  Applying this estimate to the local truncation residual of the exact semidiscrete solution yields
\begin{equation}
 \max_j\norm{U^j-U_N(t_j)}_{H^{s+4}}\le C\tau^2,
 \qquad
 \max_j\norm{U^j-U(t_j)}_{H^{s+4}}
 \le C(\tau^2+N^{-p}).
 \label{eq:generated-time-rate}
\end{equation}

\subsection{Physical Hermite readout and equation residuals}

On each interval $[t_j,t_{j+1}]$, reconstruct the physical head by the cubic Hermite polynomial having endpoint values equal to the generated physical heads and endpoint slopes equal to the corresponding components of $G_N(U^j)$ and $G_N(U^{j+1})$.  The resulting $z_{N,\tau}$ is $C^1$ in time.  The Hermite basis, \eqref{eq:generated-time-rate}, and the uniform higher time derivatives of the semidiscrete system give
\begin{align}
 \norm{z_{N,\tau}-z}_{C_tH^{s+2}}
 +\norm{\partial_tz_{N,\tau}-\partial_tz}_{C_tH^{s+1}}
 &\le C(N^{-p}+\tau^2),\notag\\
 \norm{\partial_t^2z_{N,\tau}-\partial_t^2z}_{L^\infty_tH^s}
 &\le C(N^{-p}+\tau).
 \label{eq:generated-Hermite}
\end{align}
The Riemann tensor is linear in second metric derivatives plus smooth quadratic first-derivative terms, giving \eqref{eq:generated-bosonic}.  The wave-type Yang--Mills and Higgs rows have the same derivative count.  The Dirac row is first order and therefore inherits the sharper estimate \eqref{eq:generated-Dirac}.  The harmonic, Gauss and normal Einstein constraints use only the controlled first jets and spatial derivatives and are $O(N^{-p}+\tau^2)$ in the corresponding lower norms.

Finally let $S^{(1)}$ be the complete first-derivative torsion-free Einstein--Standard-Model action representative with fixed-collar boundary completion, and let $\mathcal R$ be the Hermite reconstruction from independent nodal values and slopes.  Define
\begin{equation}
 S^{\rm cmp}_{N,\tau}=S^{(1)}\circ\mathcal R,
 \qquad
 \norm V_{0,\tau}^2=\tau\sum_j
   \left(\norm{v_j}_2^2+\tau^2\norm{\dot v_j}_2^2\right).
 \label{eq:generated-comparison-action}
\end{equation}
At the generated record, integration by parts expresses $D S^{\rm cmp}_{N,\tau}$ through the physical rows just estimated.  Since $\norm{D\mathcal R[V]}_{L^2}\le C\norm V_{0,\tau}$ and the temporal gauge and all metric entries remain independent variations before evaluation, this gives \eqref{eq:generated-stationarity}.  Keeping the first-derivative variational representative and using \eqref{eq:generated-first} gives the sharper $O(N^{-p}+\tau^2)$ continuum first-variation comparison.  In particular the complete symmetric spinor stress is included and converges strongly.


\section{The explicit local finite action and its quantitative estimates}
\label{app:native-action}

This appendix specifies the finite action used in \Cref{sec:native-branch}.  It is the local common-action representative built from the finite gravitational and Standard-Model link data, not a least-squares functional or a discretization chosen after the continuum equations are known.  The consistency and leakage estimates below analyze this fixed action without adding counterterms or filtering its physical variables.  All contractions are based at one lattice vertex.  Let $e^a_\mu$ be an oriented invertible coframe and $q^a_{\lambda\nu}$ an independent first-jet coordinate.  Define
\begin{align}
 G_{\lambda,\mu\nu}(e,q)
 &=\eta_{ab}(q^a_{\lambda\mu}e^b_\nu+e^a_\mu q^b_{\lambda\nu}),\notag\\
 \Gamma^\rho_{\mu\nu}(e,q)
 &=\tfrac12g^{\rho\sigma}(G_{\mu,\nu\sigma}+G_{\nu,\mu\sigma}-G_{\sigma,\mu\nu}),\notag\\
 \Omega_\mu(e,q)^a{}_b
 &=e^a_\rho\Gamma^\rho_{\mu\nu}(e,q)e_b^\nu-q^a_{\mu\nu}e_b^\nu.
 \label{eq:native-LC-reader}
\end{align}
At $q=\partial e$ this is the Levi--Civita connection.  Put
\[
 \omega_{\mu,h}(x)=\Omega_\mu(e(x),\delta_h^+e(x)),
 \qquad L_\mu(x)=\exp(h\omega_{\mu,h}(x)),
 \qquad U_\mu(x)=\exp(hA_\mu(x)).
\]
For $\mu<\nu$ define the logarithmic plaquettes
\begin{align}
 F^h_{\mu\nu}(x)&=h^{-2}\log\!\bigl(U_\mu(x)U_\nu(x+h e_\mu)
 U_\mu(x+h e_\nu)^{-1}U_\nu(x)^{-1}\bigr),\notag\\
 R^h_{\mu\nu}(x)&=h^{-2}\log\!\bigl(L_\mu(x)L_\nu(x+h e_\mu)
 L_\mu(x+h e_\nu)^{-1}L_\nu(x)^{-1}\bigr).
 \label{eq:native-plaquettes}
\end{align}
The logarithm is the analytic branch near the identity.  Let
\begin{equation}
 K^h_\mu=\frac{\rho_H(U_\mu)H(x+h e_\mu)-H(x)}h,
 \qquad
 V_\mu=\exp(h\sigma(\omega_{\mu,h}))\rho_S(U_\mu),
 \label{eq:native-matter-links}
\end{equation}
and
\begin{equation}
 \nabla^h_\mu\Psi=\frac{V_\mu\Psi(x+h e_\mu)-\Psi(x)}h,
 \qquad
 \nabla^{h,\vee}_\mu\bPsi
 =\frac{\bPsi(x+h e_\mu)V_\mu^{-1}-\bPsi(x)}h.
 \label{eq:native-spin-differences}
\end{equation}
With $v(e)=\sqrt{-\det g}$ and the same coefficient normalization as \eqref{eq:total-action}, set
\begin{align}
 \mathcal L_{{\rm g},h}
 &=\frac{v(e)}{2\kappa}e_a^\mu e^{b\nu}(R^h_{\mu\nu})^a{}_b
    -\frac{\Lambda}{\kappa}v(e),\notag\\
 \mathcal L_{{\rm YM},h}
 &=-\tfrac14v(e)g^{\mu\rho}g^{\nu\sigma}
      \ip{F^h_{\mu\nu}}{F^h_{\rho\sigma}}_{\boldsymbol g},\notag\\
 \mathcal L_{{\rm H},h}
 &=-v(e)g^{\mu\nu}\Rea\ip{K^h_\mu}{K^h_\nu}-v(e)V(H),\notag\\
 \mathcal L_{{\rm D},h}
 &=v(e)\Rea\left\{\tfrac{\ii}{2}
 [\bPsi\gamma^\mu(e)\nabla^h_\mu\Psi
 -(\nabla^{h,\vee}_\mu\bPsi)\gamma^\mu(e)\Psi]
 -\bPsi\mathscr M_{\boldsymbol Y}(H)\Psi\right\}.
 \label{eq:native-densities}
\end{align}
The finite local common action is
\begin{equation}
 S_h^{\rm loc}=h^4\sum_{x\in\mathscr Q_h}
 (\mathcal L_{{\rm g},h}+\mathcal L_{{\rm YM},h}
 +\mathcal L_{{\rm H},h}+\mathcal L_{{\rm D},h}).
 \label{eq:native-local-action}
\end{equation}
It has all coframe components and temporal gauge links as variation variables; temporal gauge is imposed only on the evaluated records.  The Cartan and spin links share the same coframe-derived connection.

\begin{lemma}[Regularity of the plaquette map]
\label{lem:native-plaquette}
For bounded matrices $X,Y,a,b$ the map
\[
 \Phi(h;X,Y,a,b)=h^{-2}\log(e^{hX}e^{h(Y+ha)}e^{-h(X+hb)}e^{-hY})
\]
extends analytically through $h=0$ with $\Phi(0)=a-b+[X,Y]$.  All fixed-order derivatives are uniformly bounded on compact argument sets for sufficiently small $h$.
\end{lemma}
\begin{proof}
The product inside the logarithm is $I+h^2(a-b+[X,Y])+O(h^3)$.  Since the logarithm is analytic near the identity and the linear term vanishes, the removable quotient may be written as an integral of the second $h$ derivative.  This also proves uniform bounds for parameter derivatives.
\end{proof}

\subsection{Amplitude-preserving scaling}
\label{app:native-scaling}

\begin{lemma}[Differential-order scaling]
\label{lem:native-scaling}
Under \eqref{eq:growing-derivatives}, put $\rho=hK$, $\nu=K^{-1}$ and $\widetilde y(\xi)=y(\xi/K)$, without rescaling the field amplitudes.  The exact density \eqref{eq:native-densities} has the form
\begin{equation}
 \mathcal L_h(y)=K^2\widetilde{\mathcal L}^{\rm B}_{\rho,\nu}(\widetilde y)
 +K\widetilde{\mathcal L}^{\rm D}_{\rho,\nu}(\widetilde y),
 \label{eq:native-density-scaling}
\end{equation}
where both normalized stencil densities and their required finite jet derivatives are uniformly regular for small $\rho$ and $0\le\nu\le1$.  The normalized Euler rows carry the same factors $K^2$ and $K$.
\end{lemma}
\begin{proof}
The algebraic Cartan reader is linear in a first difference, so $\omega_h=K\widetilde\omega_\rho$ and $e^{h\omega_h}=e^{\rho\widetilde\omega_\rho}$.  The Cartan plaquette therefore has factor $K^2$.  For the gauge plaquette use \Cref{lem:native-plaquette} with arguments $\nu A$; the removable quotient shows $F^h=K\widetilde F_{\rho,\nu}$, so the Yang--Mills density has factor $K^2$.  The Higgs link is
\[
 K^h_\mu=K\frac{e^{\rho\nu\rho_H(A_\mu)}H(\xi+\rho e_\mu)-H(\xi)}\rho,
\]
which is a regular first-jet function times $K$.  The spin link is
$e^{\rho(\sigma(\widetilde\omega_\rho)+\nu\rho_S(A))}$, so each Dirac difference has factor $K$ and the bilinear density has factor $K$.  Potential, mass and cosmological coefficients appear with nonnegative powers of $\nu$ inside the normalized densities.  Since the varied head amplitudes are not rescaled, differentiation of the finite sum leaves exactly the factors in \eqref{eq:native-density-scaling}.
\end{proof}

\begin{proof}[Proof of \Cref{prop:native-consistency}]
At normalized mesh $\rho$, Taylor expansion of the regular shifted-jet densities and their first variations gives an $O(\rho)$ Euler error with constants uniform in $\nu$.  Restoring the factors of \Cref{lem:native-scaling} gives
$C(K^2+K)\rho\le ChK^3$.  The same normalized plaquette expansion gives $O(\rho)$ for Cartan curvature, which scales by $K^2$, proving \eqref{eq:native-Cartan}.  The metric row differentiates the Cartan reader and the spin link throughout; no connection is frozen in this argument.
\end{proof}

\subsection{Measured spectral leakage}
\label{app:native-tails}

For the proof only, write $u=P_{\le K}u+r$ and $z_h^{\rm lo}=\mathcal I_h^{\rm trig}P_{\le K}u_h$.  The complete record remains $u_h$.

\begin{lemma}[Finite-order tail bounds]
\label{lem:native-tail-transfer}
For $0\le j\le m=k+2$,
\begin{equation}
 \max_{|\alpha|\le j}K^{-|\alpha|}\norm{\partial^\alpha(z_h-z_h^{\rm lo})}_\infty
 \le C_j\tau_{h,j}(K).
 \label{eq:native-tail-derivatives}
\end{equation}
For sufficiently small $hK$ and $\tau_{h,m}$ the full and low-frequency fields stay in the same nondegenerate chart.  Moreover
\begin{align}
 \norm{E_h^{\rm raw}(u_h)-E_h^{\rm raw}(\mathsf S_hz_h^{\rm lo})}_{0,h;Q'}
 &\le CK^2\tau_{h,2},\label{eq:native-Euler-tail}\\
 \norm{\mathcal R_{\rm B}(z_h)-\mathcal R_{\rm B}(z_h^{\rm lo})}_{L^2_tH_x^k}
 +\norm{\mathcal R_{\rm D}(z_h)-\mathcal R_{\rm D}(z_h^{\rm lo})}_{L^2_tH_x^{k+1}}
 &\le CK^{k+2}\tau_{h,m}.\label{eq:native-source-tail}
\end{align}
The corresponding zero-order physical source difference is bounded by $CK^2\tau_{h,2}$.
\end{lemma}
\begin{proof}
Differentiate the finite Fourier tail and use
$|\ell^\alpha|\le K^{|\alpha|}(1+|\ell|_1/K)^j$.  Smooth algebraic dependence of the metric, inverse metric, frame and finite representation coefficients on the field heads gives \eqref{eq:native-tail-derivatives}.  The finite Euler row is, after the exact shifted-first-jet reduction in \Cref{lem:native-scaling}, a smooth function of finitely many shifted jets through second order with overall bosonic factor at most $K^2$; the mean-value theorem gives \eqref{eq:native-Euler-tail}.  The continuum bosonic Euler rows have differential degree at most two and the Dirac rows degree one.  Differentiate their difference $k$ and $k+1$ spatial times, respectively; every product contains at least one tail factor and at most $k+2$ total field derivatives.  Equation \eqref{eq:native-tail-derivatives} gives \eqref{eq:native-source-tail}.
\end{proof}

\begin{lemma}[Logarithmic derivative upgrade for a band-limited source]
\label{lem:log-source-upgrade}
Let $f(t,x)$ extend holomorphically in the three spatial variables to width $a/K$, be bounded there by $A$, and satisfy $\norm f_{L^2(Q)}\le\epsilon\le1$.  Then for every fixed integer $j\ge1$,
\begin{equation}
 \norm f_{L^2_tH_x^j(Q)}
 \le C_j\epsilon K^j
 \left[1+\log\left(2+\frac{C_jAK^{j+3}}\epsilon\right)\right]^j.
 \label{eq:log-source-upgrade}
\end{equation}
When $\epsilon=0$ the right side is interpreted as zero.
\end{lemma}
\begin{proof}
Spatial contour shifting gives $|\widehat f(t,\ell)|\le Ae^{-a|\ell|_1/K}$.  For a frequency threshold $B$, Plancherel bounds the low frequencies by $C_j(1+B)^j\epsilon$, while the exponentially small tail is bounded by $C_jAK^{j+3}e^{-aB/(2K)}$.  Choosing $B$ proportional to $K[1+\log(2+C_jAK^{j+3}/\epsilon)]$ proves the claim.
\end{proof}

\begin{proof}[Proof of \Cref{thm:native-source}]
The full field satisfies the growing $C^5$ reserve by \Cref{lem:native-tail-transfer}; hence \Cref{prop:native-consistency} and \eqref{eq:native-sigma} give \eqref{eq:native-zero-source} directly for the unfiltered field.  For the differentiated source, \eqref{eq:native-Euler-tail} first transfers the small finite row to the low-frequency field at cost $K^2\tau_{h,2}$.  The low-frequency field has a spatial analytic strip of width $a/K$, and its bosonic and Dirac source amplitudes are bounded by $CK^2$ and $CK$.  Apply \Cref{lem:log-source-upgrade} with $j=k$ and $j=k+1$.  Finally restore the measured full tail by \eqref{eq:native-source-tail}.  This is exactly the budget \eqref{eq:native-forcing-budget}.  The Cartan assertion is the second estimate of \Cref{prop:native-consistency} applied to the full field.
\end{proof}

\begin{proposition}[A sufficient nodal-accuracy condition]
\label{prop:native-rounding}
Let $u_h^b$ have Fourier support in $|\ell|_\infty\le K$ and let $u_h=u_h^b+e_h$ with $\max_x|e_h(x)|\le e$, preserving the affine field-chart constraints.  Then
\begin{equation}
 \tau_{h,j}(K)\le C_j e\,h^{-2}(hK)^{-j}.
 \label{eq:native-rounding-tail}
\end{equation}
In particular, with $m=k+2$, the sufficient tolerance
\begin{equation}
 e\le c\,h^{m+3}K^{m+2}
 \label{eq:native-rounding-tolerance}
\end{equation}
implies $\tau_{h,m}=O(hK^2)$ and, on the rate families of \Cref{cor:native-rate}, $\tau_{h,2}=O(hK)$.
\end{proposition}
\begin{proof}
Parseval for the normalized four-dimensional discrete Fourier transform gives $\sum_\ell|\widehat e_h(\ell)|^2\le e^2$.  There are $O(h^{-4})$ Fourier modes, so Cauchy--Schwarz bounds the absolute coefficient sum by $Ch^{-2}e$.  On the grid spectrum $(1+|\ell|_1/K)^j\le C_j(hK)^{-j}$, which proves \eqref{eq:native-rounding-tail}.  Substitute \eqref{eq:native-rounding-tolerance} for the final assertion.
\end{proof}


\section{Finite-prefix stationarity from the same local action}
\label{app:prefix-stationarity}

This appendix proves \Cref{thm:prefix-stationarity,cor:prefix-table}.  The argument does not define a new field evolution and does not replace the indefinite Lorentzian action by a least-squares cost.  It studies any accepted finite prefix using the same local action \eqref{eq:native-local-action} and a positive operational potential supplied by the source geometry.

\begin{lemma}[Shifted-first-jet normal form of the unchanged local action]
\label{lem:native-firstjet-normal-form}
On a bounded nondegenerate coframe/logarithm chart, the action \eqref{eq:native-local-action} has the exact interior representation
\begin{equation}
 S_h^{\rm loc}[y]
 =h^4\sum_{x\in Q'_h}F_h^{(1)}\!\left(\Xi_h^{(1)}y(x)\right)
 +\mathscr B_h[y],
 \label{eq:native-firstjet-normal-form}
\end{equation}
where $\Xi_h^{(1)}y$ consists of a fixed finite set of shifted field values and normalized first differences, $F_h^{(1)}$ is smooth with fixed-order derivatives bounded uniformly for small $h$, and $\mathscr B_h$ depends only on the fixed collars for the interior variation problem.  Therefore the interior Euler covector and all of its mixed physical derivatives are exactly those of the shifted-first-jet bulk representative; no continuum approximation is used in this identity.
\end{lemma}
\begin{proof}
Only the gravitational plaquette needs a reduction.  Write its curvature logarithm as the exact discrete curl of the coframe connection plus a remainder that is a regular function of shifted connections.  If $C^{\mu\nu}(e)$ is the Palatini coframe coefficient, the discrete product identity
\begin{equation}
 \ip{C}{\delta_\mu^+W}
 =-\ip{\delta_\mu^-C}{W}
   +\delta_\mu^-\ip{C}{T_\mu W}
 \label{eq:native-SBP}
\end{equation}
converts both curvature-curl terms into first-jet bulk expressions plus an exact lattice divergence.  Moreover, with $e_-=T_{-\mu}e$ and $p_-=T_{-\mu}\delta_\mu^+e$,
\begin{equation}
 \delta_\mu^-C(e)
 =\int_0^1 DC(e_-+thp_-)[p_-]\,\dd t,
 \label{eq:native-C-firstjet}
\end{equation}
which is a smooth shifted-first-jet function with no inverse power of $h$.  Summing the divergence leaves only the fixed collar functional $\mathscr B_h$.

The Yang--Mills plaquette is already a uniformly regular function of shifted gauge values and first differences by \Cref{lem:native-plaquette}; the Higgs and Dirac link terms in \eqref{eq:native-matter-links}--\eqref{eq:native-spin-differences} are shifted first-jet functions directly.  Thus all sectors share the representation \eqref{eq:native-firstjet-normal-form}.  Since the transformation is an exact discrete summation-by-parts identity, interior first and higher physical variations are unchanged.
\end{proof}

\begin{lemma}[Uniform upper Hessian scale of the local action]
\label{lem:native-action-Hessian}
On a compact convex subchart of the finite field coordinates, with all shifted first-jet arguments of \eqref{eq:native-densities} remaining in a common compact coefficient set, there are constants $C_1,K_0,D_0$ independent of $h$ and of the number of nodes such that
\begin{align}
 |D^2\mathcal A_h(y)[u,v]|&\le C_1\norm u_{1,h}\norm v_{1,h},\label{eq:action-Hessian-bilinear}\\
 \norm{\nabla_h^2\mathcal A_h(y)}_{0,h\to0,h}&\le K_0h^{-2},\label{eq:action-Hessian-scale}\\
 \operatorname{osc}_{\Theta_h}\mathcal A_h&\le D_0,\label{eq:action-span}
\end{align}
where
\[
 \norm u_{1,h}^2=\norm u_{0,h}^2
 +\sum_{\mu=0}^3\norm{\delta_\mu^+u}_{0,h}^2.
\]
These are upper bounds only; no positivity or inverse for the Lorentzian action Hessian is asserted.
\end{lemma}
\begin{proof}
By \Cref{lem:native-firstjet-normal-form}, the interior local density is exactly a smooth function of finitely many shifted first jets.  Differentiating the finite sum twice gives $h^4$ times a sum of uniformly bounded coefficient bilinear forms applied to shifted values and first differences of $u,v$.  Each shift occurs only a fixed number of times.  Cauchy--Schwarz gives \eqref{eq:action-Hessian-bilinear}.  The elementary inverse estimate $\norm{\delta_\mu^+u}_{0,h}\le2h^{-1}\norm u_{0,h}$ gives \eqref{eq:action-Hessian-scale}.  Finally the density is uniformly bounded on the compact chart and $h^4$ times the number of lattice cells in the fixed cylinder is uniformly bounded, giving \eqref{eq:action-span}.  The coframe/spin-connection chain is contained in the smooth density and is differentiated throughout.
\end{proof}

\begin{proof}[Proof of \Cref{thm:prefix-stationarity}]
Let $x=\theta_j$, $y=\theta_{j+1}$, $d=y-x$, and write $p_h=\nabla_h\Psi_h$.  By \eqref{eq:action-Hessian-scale}, Taylor's integral formula gives
\begin{equation}
 \mathcal A_h(x)-\mathcal A_h(y)
 \ge -(E_h(y),d)_{0,h}-\tfrac12K_0h^{-2}\norm d_{0,h}^2.
 \label{eq:prefix-descent1}
\end{equation}
The lower bound in \eqref{eq:prefix-Fisher} yields
\[
 (p_h(y)-p_h(x),d)_{0,h}\ge g_*\norm d_{0,h}^2.
\]
Substitute \eqref{eq:prefix-balance} and put
\[
 a_h=g_*\upsilon_h^{-1}-\tfrac12K_0h^{-2}=a_0h^{-2}>0,
 \qquad
 \ell_h=g^*\upsilon_h^{-1}+K_0h^{-2}=\ell_0h^{-2}.
\]
Young's inequality gives
\begin{equation}
 \mathcal A_h(x)-\mathcal A_h(y)
 \ge \tfrac12a_h\norm d_{0,h}^2
 -\frac1{2a_h}\norm{b_j}_{0,h}^2.
 \label{eq:prefix-descent2}
\end{equation}
Moreover
\[
 E_h(x)=b_j-\upsilon_h^{-1}(p_h(y)-p_h(x))
                 +(E_h(x)-E_h(y)),
\]
so
\[
 \norm{E_h(x)}_{0,h}\le\norm{b_j}_{0,h}+\ell_h\norm d_{0,h}.
\]
Combining with \eqref{eq:prefix-descent2}, squaring, and summing over $j$ gives
\[
 \sum_{j<J}\norm{E_h(\theta_j)}_{0,h}^2
 \le \frac{4\ell_h^2}{a_h}
       \{\mathcal A_h(\theta_0)-\mathcal A_h(\theta_J)\}
 +\left(2+\frac{2\ell_h^2}{a_h^2}\right)
       \sum_{j<J}\norm{b_j}_{0,h}^2.
\]
Since $4\ell_h^2/a_h=C_0h^{-2}$, the second coefficient is a cutoff-independent constant, and the telescoped action decrement is at most the uniform oscillation $D_0$, this is \eqref{eq:prefix-average}.  Taking the minimum of the finite list gives \eqref{eq:prefix-selected}.
\end{proof}

\begin{proof}[Proof of \Cref{cor:prefix-table}]
Let $G$ be the event that the accepted prefix stays inside the declared safe set.  On $G$, \eqref{eq:prefix-average} holds pathwise.  Multiplying by $\mathbf1_G$, taking expectations, and enlarging the nonnegative balance sum to all safe occupied edges gives
\[
 \mathbb E\!\left[\mathbf1_G\frac1J\sum_{j<J}
     \norm{E_h(\theta_j)}_{0,h}^2\right]
 \le C_0h^{-2}\frac{D_0}{J}+D_1\mathcal Q_{h,J}^{\rm bal}.
\]
If the best safe slope exceeds $h^\alpha$, the mean square exceeds $h^{2\alpha}$.  Markov's inequality therefore gives the last two terms of \eqref{eq:prefix-table-probability}; adding $\mathbb P(G^c)=p_h^{\rm exit}$ proves the bound.  Summability along dyadic cutoffs and the first Borel--Cantelli lemma yield eventual certification without independence.
\end{proof}


\end{document}
